\documentclass[11pt,a4paper]{article}
\usepackage{fontspec}
\IfFontExistsTF{Malgun Gothic}{\setmainfont{Malgun Gothic}}{\setmainfont{Noto Sans CJK KR}}
\XeTeXlinebreaklocale "ko"
\XeTeXlinebreakskip = 0pt plus 1pt
\usepackage{amsmath,amssymb,bm}
\setlength{\textwidth}{162mm}
\setlength{\textheight}{242mm}
\setlength{\oddsidemargin}{-1.4mm}
\setlength{\evensidemargin}{-1.4mm}
\setlength{\topmargin}{-12mm}
\renewcommand{\contentsname}{목차}
\emergencystretch=3em
\setlength{\parindent}{0pt}
\setlength{\parskip}{0.55em}
\setcounter{tocdepth}{2}
\title{주식, 채권, 파생금융상품 1: 이론\\\large 슬라이드별 학습 가이드}
\author{}
\date{}
\begin{document}
\maketitle
\tableofcontents
\clearpage

\section{1부 · 투자와 금융자산}

\subsection{슬라이드 1: 금융투자의 지도를 먼저 그리기}

금융시장은 서로 다른 현금흐름을 시간과 상태에 따라 교환하는 곳이다. 주식은 기업의 잔여 성과를, 채권은 약속된 지급을, 파생상품은 기초자산이나 사건에 연동된 지급을 제공한다. 이 셋의 가격을 이해하려면 먼저 \textbf{언제 얼마를 받는가}, \textbf{어떤 상태에서 못 받을 수 있는가}, \textbf{그 위험을 누가 떠안는가}를 구분해야 한다.
\par

이 가이드의 순서는 투자 대상의 분류에서 시작해 수익률의 확률분포를 정의하고, 마지막으로 서로 다른 자산을 결합할 때 위험이 어떻게 바뀌는지를 계산하는 것이다. 뒤의 포트폴리오 수식은 앞의 현금흐름·기대값·공분산을 재조합한 결과다.
\par

\paragraph{상태별 현금흐름으로 자산 읽기}

같은 평균 지급액이라도 \textbf{언제, 어떤 상태에서} 받는지가 다르면 투자자에게 주는 가치가 다르다. 호황과 불황의 확률이 각각 절반일 때 A가 두 상태에서 모두 100을 지급하고 B가 호황에 160, 불황에 40을 지급한다면 두 자산의 기대 지급액은 100으로 같다. 그러나 불황에 등록금이나 부채를 갚아야 하는 사람에게 B는 필요한 때 현금을 적게 준다. 가격을 비교하기 전에 상태별 지급액을 써 보아야 하는 이유다.
\par

상품을 비교할 때 지급액과 투자 수익률을 구분한다. \textbf{현금흐름}은 실제로 받거나 내는 금액이고, \textbf{수익률}은 그 금액을 최초 투자액으로 나눈 비율이다. 같은 10만원 지급도 매입가격이 100만원이면 10\%, 200만원이면 5\%에 해당한다.
\par

\paragraph{채권은 지급액이 약속돼 있는데 왜 상태별 현금흐름을 보나?}

약속과 실현은 다르기 때문이다. 액면 100인 회사채도 기업이 정상 영업하면 100을 지급하지만, 부도 상태에서는 회수액이 60일 수 있다. 주식은 그보다 뒤에 남는 금액을 받는다. 호황·불황의 두 칸에 실제 지급액을 적으면 '채권은 고정, 주식은 변동'이라는 이름만으로 놓치는 신용위험과 지급 우선순위가 드러난다.
\par

\subsection{슬라이드 2: 투자이론이 다루는 질문}

투자이론은 현재 소비를 미래 소비로 옮기는 선택, 불확실한 현금흐름의 가격, 투자자 사이의 위험 배분을 다룬다. 미시경제학의 기대효용은 위험을 감수할 의사를, 거시경제학의 시간 간 선택은 현재와 미래 소비의 교환비율을 이해하는 출발점이다.
\par

가령 같은 100만원을 오늘 소비할 수도, 1년 만기 예금에 둘 수도, 주식에 넣을 수도 있다. 세 선택을 비교하려면 명목 수익률뿐 아니라 물가, 세금, 유동성, 손실 가능성, 투자 기간을 같은 기준으로 맞춰야 한다.
\par

\paragraph{현재 소비와 미래 소비의 교환}

예금의 명목 연수익률이 8\%이고 1년 물가상승률이 5\%라면 구매력 기준 수익률은 단순히 3\%라고만 계산하면 근사치다. 정확히는 $1+r_{\mathrm{real}}=(1.08)/(1.05)$이므로 $r_{\mathrm{real}}\simeq2.86\%$다. 오늘 소비재 1단위를 포기해 내년에 몇 단위를 살 수 있는지를 계산한 것이다. 기간과 통화가 달라지면 같은 비교가 성립하지 않는다.
\par

위험한 자산을 평가할 때는 평균 지급액 외에 지급이 줄어드는 상태도 본다. 미시경제학의 한계효용은 소비가 적을 때 일반적으로 더 크다. 따라서 불황에 잘 지급하는 자산은 평균 수익이 낮아도 가치가 있을 수 있다. 이 관점에서 공분산은 \textbf{소비가 부족한 상태에 자산의 지급액이 어떻게 움직이는지}를 읽는 도구가 된다.
\par

\paragraph{기대수익률이 더 높은 투자라면 언제나 선택해야 하나?}

1년 뒤 등록금 1,000만원이 필요한 학생에게는 평균적으로 15\%를 버는 주식보다 만기에 필요한 금액을 확정하는 자산이 더 적합할 수 있다. 주식이 불황에 -30\%를 기록하면 평균 수익률이 높아도 그날의 지출을 감당하지 못한다. \textbf{필요한 날짜의 소비와 손실 상태}를 지정한 뒤 수익률을 비교해야 한다.
\par

\subsection{슬라이드 3: 자산가격과 자산배분}

\textbf{자산가격}은 미래 현금흐름의 현재 교환비율을 묻는다. \textbf{자산배분}은 주어진 가격에서 투자자의 자금을 자산별로 얼마나 나눌지를 묻는다. 개인은 은퇴 시점과 손실 감내 능력이 중요하고, 보험사는 약속된 보험금의 지급 시점과 자산의 만기가 맞는지도 중요하다.
\par

주식에 초점을 맞추더라도 채권과 파생상품을 함께 배우는 이유는 비교 기준이 필요하기 때문이다. 국채의 수익률은 주식의 초과수익을 계산하는 기준이 되고, 선물·옵션은 보유한 주식의 위험을 옮기거나 특정 상태의 지급을 분리한다.
\par

\paragraph{가격 질문과 선택 질문의 순서}

가격과 선택은 연결되지만 질문이 다르다. 채권의 미래 지급액 105가 확정이고 비교 가능한 1년 금리가 5\%라면 현재가격은 $105/1.05=100$이다. 이 가격이 주어진 뒤 투자자는 채권을 얼마나 살지 결정한다. 예상 현금지출이 곧 닥치는 사람은 가격이 공정해도 주식보다 현금성 자산을 더 보유할 수 있다.
\par

주식과 국채를 비교할 때 국채 수익률을 '모든 위험자산의 적정 수익률'로 쓰면 안 된다. 위험자산은 불황에 손실을 낼 수 있으므로 같은 기대 현금흐름이라도 더 낮은 가격에 거래될 수 있다. 반대로 위험을 줄이는 자산은 낮은 기대수익을 감수할 이유가 있다. \textbf{시장가격의 설명}과 \textbf{개인의 적합성 판단}을 구분하자.
\par

\paragraph{시장가격이 적정하면 모든 투자자가 같은 비중을 사야 하나?}

동일한 가격 100의 장기채라도 10년 뒤 보험금 지급 의무가 있는 보험사에는 부채와 만기를 맞추는 자산이고, 다음 달 생활비가 필요한 가계에는 중간 매도 위험이 있는 자산이다. 가격은 시장에서 하나로 관찰되지만 투자자의 부채, 기존 자산, 투자기간은 다르므로 최적 보유량도 달라진다.
\par

\subsection{슬라이드 4: 금융 분야의 연결 고리}

자산가격론은 위험이 가격에 어떻게 반영되는지, 기간구조론은 만기가 다른 채권의 수익률이 어떻게 연결되는지, 위험관리는 큰 손실의 가능성을 어떻게 측정하고 제한하는지를 연구한다. 매매전략은 이 이론을 거래비용·제약·시장 유동성이 있는 환경에서 실행하는 문제다.
\par

같은 경제 충격은 여러 시장에 전달된다. 금리 변화는 채권의 할인율과 주식의 미래 이익 전망에 함께 영향을 줄 수 있다. 각 자산의 가격 변화를 설명할 때는 \textbf{예상 현금흐름의 변화}와 \textbf{할인율의 변화}를 나누어 본다.
\par

\paragraph{한 금리 변화가 여러 시장에 전파되는 경로}

고정금리 채권의 가격은 약속된 미래 지급액을 현재로 할인한 합이다. 금리가 오르면 각 할인분모가 커져 기존 채권 가격이 떨어진다. 주식은 채권과 달리 미래 현금흐름 자체도 경기와 함께 변한다. 금리 상승기에 예상 기업이익이 동시에 크게 오르면 할인율 상승의 가격 하락 효과를 일부 상쇄할 수 있다. 따라서 '금리 상승이면 모든 주식이 하락'처럼 하나의 방향을 기계적으로 외우면 안 된다.
\par

위험관리에서는 같은 사건이 여러 자산에 동시에 닿는지도 본다. 금리 위험을 채권만의 문제로 보면, 대출금리와 부채의 현재가치가 같이 움직이는 은행이나 보험사의 \textbf{순자산} 변화를 놓친다. 개별 가격보다 자산과 부채를 합친 순현금흐름이 의사결정의 대상이다.
\par

\paragraph{금리 상승이 모든 금융기관에 똑같이 손실인가?}

금리가 1\%포인트 오를 때 장기채 자산 가치가 8\% 떨어지는 보험사를 생각하자. 같은 금리에 민감한 미래 보험금 부채의 현재가치도 9\% 줄어들면 순자산은 오히려 개선될 수 있다. 자산가격 하나의 방향으로 결론 내리지 말고 \textbf{자산 변화액에서 부채 변화액을 뺀 값}을 계산해야 한다. 은행처럼 단기 예금과 장기 대출을 조합한 기관은 또 다른 민감도를 갖는다.
\par

\subsection{슬라이드 5: 필요한 사전 개념}

필수 언어는 현재가치, 복리, 확률변수, 기대값, 미분이다. 현재가치 $PV=CF/(1+r)$는 한 기간 뒤 확정된 현금흐름 $CF$를 오늘의 금액으로 바꾼다. 불확실한 현금흐름에는 지급액의 평균과 함께 \textbf{어떤 경제 상태에서 지급되는지}가 가격에 반영된다. 소비가 부족한 불황에 지급액도 줄어드는 자산이라면 투자자는 더 큰 위험 보상을 요구한다.
\par

앞으로 수익률은 원칙적으로 \textbf{동일한 보유기간의 단순 수익률}로 비교한다. 일별, 월별, 연별 숫자를 섞으면 분산과 샤프비율의 단위가 달라져 잘못된 결론에 도달한다.
\par

\paragraph{현재가치 공식의 적용 범위}

두 기간 뒤 확정 현금흐름은 $PV=CF_2/(1+r)^2$로 할인한다. 기간마다 금리가 다르면 각 만기의 할인계수를 따로 쓴다. 예를 들어 1년 뒤 100과 2년 뒤 100은 같은 금액이어도 오늘 가격과 금리 민감도가 다르다.
\par

확률적인 지급액에는 먼저 상태별 현금흐름을 적는다. $E[CF]$를 무위험 금리로 할인하면 위험 보상이 사라지므로 일반적으로 올바른 가격이 아니다. 학부 수준에서는 '기대현금흐름을 위험에 맞게 할인한다'는 직관으로 시작하고, 구체적 모형에서는 어떤 위험이 가격에 반영되는지 따진다. \textbf{확률의 평균을 구하는 작업과 가격을 매기는 작업은 별개}다.
\par

\paragraph{미래 지급액의 평균을 예금금리로 할인해도 되나?}

1년 뒤 호황에 150, 불황에 50을 같은 확률로 받는 청구권의 평균 지급액은 100이다. 예금금리가 5\%라고 $100/1.05$를 가격으로 정하면 두 상태의 지급 시점은 맞추지만 위험에 대한 보상은 반영하지 못한다. 특히 불황에 50만 받는 손실이 투자자에게 크게 부담되면 같은 평균을 가진 확정 지급 100보다 낮은 가격을 요구할 수 있다.
\par

\subsection{슬라이드 6: 성과 측정의 구조}

투자 성과는 투자 결과만으로 판단할 수 없다. 예컨대 주식시장이 20\% 오른 해에 10\%를 벌었다면 절대수익은 양수지만 시장 대비 성과는 음수다. 평가에는 \textbf{기준지수(benchmark)}, \textbf{위험}, \textbf{평가기간}, \textbf{비용}을 함께 지정해야 한다.
\par

기준지수 대비 수익률은 $R_p-R_b$로 계산한다. 여기에 수익률의 변동성과 비용을 함께 놓으면 높은 수익이 더 큰 위험의 결과인지 살필 수 있다. 예컨대 같은 초과수익 2\%포인트라도 운용 결과의 흔들림이 큰 전략은 투자 목표를 달성할 확률이 낮을 수 있다.
\par

\paragraph{성과를 비교 가능한 숫자로 만들기}

운용 성과를 평가할 때 먼저 기준지수를 정한다. 포트폴리오 수익률이 12\%, 같은 기간 기준지수 수익률이 10\%라면 active return(기준지수 대비 수익률)은 2\%포인트다. 그러나 기준지수보다 훨씬 큰 위험을 감수했거나 비용을 더 냈다면 이 숫자만으로 운용 능력을 판단할 수 없다. 평가기간이 다른 두 성과를 비교할 때는 동일 기간으로 환산하고 복리 방식도 맞춘다.
\par

수수료가 자산의 1\%이고 거래비용이 0.5\%라면 투자자에게 남는 수익률은 총수익률과 다르다. \textbf{사전 기대 성과}, \textbf{사후 실현 성과}, \textbf{비용 차감 후 성과}를 따로 기록해야 평가의 질문이 분명해진다.
\par

\paragraph{같은 초과수익이라도 active risk가 다르면?}

두 운용자가 같은 기준지수와 기간을 썼다고 하자. A의 연평균 active return은 2\%포인트, 기준지수와의 수익률 차이의 표준편차인 tracking error는 4\%포인트다. B는 각각 1\%포인트와 1\%포인트다. active return을 tracking error로 나눈 information ratio는 A가 $2/4=.5$, B가 $1/1=1$이다. A의 초과수익이 더 크지만 \textbf{기준에서 벗어난 위험 한 단위당 초과수익}은 B가 높다. 비용 차감 전후 수익률과 표본기간도 맞춰 비교해야 한다.
\par

\subsection{슬라이드 7: 공부의 기준점}

투자론 교재를 읽을 때는 공식의 암기보다 각 공식이 쓰이는 \textbf{의사결정 문제와 가정}을 먼저 확인한다. 평균-분산 모형은 투자자가 수익률의 평균과 분산으로 선택을 요약할 수 있다는 가정 위에서 힘을 얻는다. 실제 수익률의 비대칭성, 극단적 손실, 거래비용은 이 가정 밖의 문제다.
\par

한 개념을 익힐 때는 현금흐름 예시 → 확률변수 → 기대값·분산 → 경제학적 해석 → 제약을 적용한 현실 사례 순으로 따라가면 수식이 고립되지 않는다.
\par

\paragraph{공식을 읽는 다섯 단계}

새 공식을 만나면 먼저 (1) 무엇을 선택하는 식인지, (2) 단위가 무엇인지, (3) 고정한 변수가 무엇인지, (4) 극단값에서 말이 되는지, (5) 제약이 바뀌면 어디가 달라지는지 확인한다. 예컨대 $\sigma_p^2=w^2\sigma_1^2+(1-w)^2\sigma_2^2+2w(1-w)c$에서 $w=1$을 넣으면 반드시 $\sigma_1^2$, $w=0$을 넣으면 $\sigma_2^2$가 나와야 한다. 이 경계 검산만으로 비중을 거꾸로 쓴 오류를 잡을 수 있다.
\par

수식의 결과가 음수가 될 수 없는지 같은 정의도 검산에 쓴다. 분산은 평균에서의 편차를 제곱한 기대값이므로 음수일 수 없다. 계산에서 음의 분산이 나오면 공분산의 단위, 상관행렬의 타당성, 대입한 숫자를 다시 확인해야 한다.
\par

\paragraph{공식이 맞는지 숫자 없이 확인하는 방법은?}

두 자산 비중을 $w$와 $1-w$로 둔 식이라면 $w=1$을 넣었을 때 자산 1 하나의 결과, $w=0$을 넣었을 때 자산 2 하나의 결과가 나와야 한다. 예를 들어 분산식에서 $w=1$인데 $\sigma_2^2$가 남으면 비중을 거꾸로 쓴 것이다. 또 수익률의 분산에는 수익률의 제곱 단위가, 표준편차에는 수익률 단위가 있어야 한다.
\par

\subsection{슬라이드 8: 금융자산의 큰 분류}

주식, 채권, 파생상품, 혼합증권, 대체자산은 법적 권리와 현금흐름의 형태가 다르다. 보통주는 기업의 잔여청구권, 국채는 정부의 지급약속, 선물은 미래 거래가격에 관한 계약, 옵션은 특정 가격에 거래할 \textbf{권리}다. 전환사채는 채권에 주식 전환권이 결합된 예다.
\par

자산군 이름만으로 위험을 단정할 수는 없다. 단기 국채는 명목 원금 변동이 작아도 물가 위험을 지고, 우량주라도 가격이 크게 흔들릴 수 있다. ABS·MBS는 기초 대출의 상환과 연체에 좌우되고, 대체자산은 평가 빈도가 낮아 관측 변동성이 작아 보일 수 있다.
\par

\paragraph{청구권의 우선순위와 혼합증권}

기업 자산을 처분해 80을 회수했는데 채권자에게 70을 지급해야 한다면, 단순한 청산 예에서 먼저 채권자에게 70을 지급하고 보통주 주주에게 10이 남는다. 회수액이 60이라면 보통주에는 0이 남고 채권자도 손실을 본다. 이 \textbf{지급 순위} 때문에 보통주를 잔여청구권이라 부른다. 실제 우선순위에는 담보권, 법적 비용, 여러 종류의 채무가 더해진다.
\par

전환사채는 고정 지급 약속과 주식으로 바꿀 권리를 함께 가진다. 따라서 '채권'이라는 이름만으로 금리 위험만 판단할 수 없다. 기업 주가가 오르면 전환권의 가치가 커지고, 신용이 나빠지면 채권 부분의 가치가 떨어질 수 있다. 상품을 분류한 다음 \textbf{지급액을 구성 요소로 분해}하면 위험이 보인다.
\par

\paragraph{전환권은 언제 행사할 가치가 있나?}

액면 100만원의 전환사채를 주식 2주로 바꿀 수 있다고 하자. 주가가 60만원이면 전환해 받는 주식의 시가는 120만원이다. 주가가 40만원이면 전환가치는 80만원이므로, 만기에 원금 100만원을 온전히 받을 수 있다면 전환하지 않는 편이 유리하다. 실제 선택에는 남은 이자, 만기까지의 시간, 발행사 신용위험이 더해진다. 전환권의 가치가 주가와 함께 커지는 이유를 \textbf{두 지급액의 비교}로 확인할 수 있다.
\par

\subsection{슬라이드 9: 주식 운용: 수동과 능동}

패시브(passive) 운용은 특정 시장지수를 낮은 비용으로 추종한다. 액티브(active) 운용은 가격이 본질가치에서 벗어났거나 시장 국면을 예측할 수 있다고 보고 포트폴리오를 바꾼다. 종목 선택(stock selection)은 기업 간 상대가격, 시장 타이밍(market timing)은 전체 주식 비중에 관한 판단이다.
\par

평균-분산 분석은 기대수익과 공분산을 이용해 위험 대비 조합을 찾는다. CAPM은 모든 투자자가 공통의 효율적 시장포트폴리오를 가진다는 강한 조건에서 개별 자산의 \textbf{시장과 함께 움직이는 위험}에 보상이 붙는다고 설명한다. 모멘텀과 역추세 전략은 과거 가격 패턴을 이용하지만 거래비용과 표본 밖 성과를 반드시 검토해야 한다.
\par

\paragraph{패시브 운용의 기준과 액티브 운용의 부담}

지수 추종 상품도 무엇을 추종하는지가 중요하다. 시가총액 가중지수는 큰 기업의 비중이 높고, 동일가중지수는 기업마다 같은 비중을 목표로 해 정기 재조정이 필요하다. 두 지수는 같은 국가의 주식을 담아도 산업 집중도와 거래비용이 달라진다. 패시브 운용의 핵심 결정은 \textbf{추종할 지수의 선택과 tracking error 관리}다. Tracking error는 여러 기간에 걸친 펀드와 기준지수의 수익률 차이(active return)의 표준편차다.
\par

능동 전략이 유리한지를 판단할 때는 예측이 맞은 횟수만 보지 않는다. 틀렸을 때의 손실 크기, 거래 횟수, 수수료, 세금, 시장에 거래 내용을 노출했을 때의 가격 영향까지 차감해야 한다. 모멘텀의 역사적 패턴을 발견했다면 자료를 고른 뒤 우연히 맞은 결과인지 다른 기간에서 검증하는 절차가 필요하다.
\par

\paragraph{지수형 ETF를 사면 지수 수익률과 정확히 같나?}

지수를 추종하는 펀드에도 운용보수와 매매비용이 들고, 구성종목을 바꾸는 날짜나 현금 보유 비중 때문에 지수와 수익률 차이가 생길 수 있다. 예를 들어 지수가 10\% 올랐는데 비용 등을 반영한 tracking difference가 -0.4\%포인트라면 펀드 수익률은 9.6\%다. 패시브 운용은 예측을 덜 요구하지만 \textbf{기준지수의 선택과 실제 추적 정도}를 확인해야 한다. 한 기간의 수익률 차이인 tracking difference와 여러 기간의 차이가 흔들리는 정도인 tracking error는 구분한다.
\par

\subsection{슬라이드 10: 지수 비교를 읽는 방법}

다우, 닛케이, 코스피 같은 지수는 국가별 주식시장의 일부를 요약한다. 지수 수준을 그대로 비교하면 기준일과 산출 방식이 달라 의미가 없다. 같은 시작점을 100으로 정규화하고 같은 통화·배당 재투자 기준을 적용해야 기간 수익률을 비교할 수 있다.
\par

1990년 말의 가격지수를 모두 100으로 놓은 비교에서는 2025년 말 무렵 S\&P 500이 약 2,073, 코스피가 약 605, FTSE 100이 약 463, 닛케이 225가 약 211로 표시된다. 이는 시작점에 대한 \textbf{각 시장의 자국통화 가격지수 배율}이다. 배당과 환율 효과가 빠져 있으므로 같은 기간 한국 투자자가 실제로 각 시장에 투자해 얻었을 원화 총수익률 순위라고 해석하면 안 된다. 1990년을 시작점으로 잡은 선택 자체도 일본 시장의 장기 비교에 큰 영향을 준다.
\par

국가별 수익률 차이는 환율, 산업 구성, 배당, 지수 편입 규칙과 경기의 차이를 반영한다. 예를 들어 원화 투자자의 미국 주식 수익률은 달러 주가 수익률에 원/달러 환율 변화가 결합된다. 단순수익률로 쓰면 $1+R_{\mathrm{KRW}}=(1+R_{\mathrm{USD}})(1+R_{\mathrm{FX}})$이다.
\par

\paragraph{지수와 투자자의 실제 수익률 사이}

가격지수가 100에서 120으로 올랐으면 가격 상승률은 20\%다. 같은 기간 현금배당 3을 받아 재투자하지 않았다면 단순 총수익률은 $(120+3-100)/100=23\%$다. 배당을 중간에 재투자했다면 재투자 시점의 가격이 필요하므로 최종 가격과 배당 합계만으로 정확한 누적 수익률을 계산할 수 없다. \textbf{가격지수와 총수익지수}를 비교할 때 이 차이가 핵심이다.
\par

해외지수를 비교하면 통화가 하나 더 들어간다. 달러 지수가 20\% 오르고 달러가 원화 대비 10\% 하락하면 원화 기준 가격수익률은 $1.20\times.90-1=8\%$다. 달러 수익 20\%를 곧바로 한국 투자자의 20\%라고 읽을 수 없다. 과거의 지수 끝값은 배당·환율·세금·투자 가능한 상품을 확인해야 비로소 투자 성과로 바뀐다.
\par

\paragraph{가격가중지수와 시가총액가중지수는 왜 다르게 움직이나?}

가격가중지수에서는 주당 가격이 높은 종목의 변화가 크게 반영된다. 시가총액가중지수에서는 주가에 발행주식 수를 곱한 기업 규모가 비중을 정한다. 시가총액 1조원인 기업의 주가가 10\% 오르고 9조원인 기업이 그대로라면 두 종목의 시가총액가중지수는 약 1\% 오른다. 두 주식의 주당 가격이 같았다면 가격가중지수는 약 5\% 오른다. 시작값을 모두 100으로 맞춰도 \textbf{가중 규칙}이 다르면 지수 수익률이 달라진다.
\par

\subsection{슬라이드 11: 채권 운용의 지도}

채권은 \textbf{현금흐름의 지급 시점}이 핵심이다. 무이표채 가격 $P=F/(1+y)^T$에서 수익률 $y$가 오르면 가격은 내려간다. 듀레이션은 이 관계의 일차 기울기, 컨벡시티는 곡률을 근사한다. 현금흐름 매칭(cash-flow matching)은 미래 지급 의무와 채권 지급액을 맞추는 전략이다.
\par

국채와 회사채를 비교할 때는 같은 만기의 기준 금리에 부도확률, 부도 시 회수율, 유동성을 더해 살핀다. 구조형 신용모형은 기업 자산가치와 부채의 관계로, 강도모형은 시간에 따른 부도 도착률로 신용위험을 표현한다.
\par

\paragraph{채권의 금리 민감도를 직접 미분하기}

만기 $T$의 무이표채가 만기에 $F$를 지급하고 연복리 수익률이 $y$라면 $P(y)=F(1+y)^{-T}$다. 미분하면 $dP/dy=-TF(1+y)^{-T-1}=-TP/(1+y)$이므로, 가격에 대한 상대 민감도는 $-(1/P)dP/dy=T/(1+y)$다. 금리가 1\%포인트 오를 때 가격의 대략적인 하락률은 이 값에 $.01$을 곱한 것이다. 만기가 길수록 같은 금리 변화에 더 민감한 이유가 수식에 드러난다.
\par

쿠폰채는 여러 날짜의 지급액을 더한 것이므로 각 현금흐름에 같은 논리를 적용해 가중평균한다. 다만 듀레이션은 작은 금리 변화의 \textbf{일차 근사}다. 변화 폭이 커지면 가격 곡선의 굽음, 곧 컨벡시티를 고려해야 한다. 채권을 만기까지 보유할 계획이라도 중간에 팔거나 담보로 쓸 가능성이 있다면 시가 변동은 중요하다.
\par

\paragraph{쿠폰율 5\%인 채권의 투자 수익률도 5\%인가?}

액면 100, 1년 뒤 쿠폰 5와 원금 100을 주는 채권을 95에 샀다면 만기까지 보유할 때 약정 지급액 105를 받는다. 단순 만기수익률은 $105/95-1\simeq10.53\%$다. 쿠폰율 5\%는 \textbf{액면 대비 쿠폰}이고, 투자 수익률은 \textbf{실제 매입가격 대비 총지급액}을 사용한다. 부도나 매매비용이 없다는 조건도 함께 적어야 한다.
\par

\subsection{슬라이드 12: 금리 그래프의 독해}

국가별 10년물 금리의 과거 추이는 각국의 통화정책, 물가 기대, 성장 전망, 안전자산 수요가 서로 다르게 움직였음을 보여준다. 장기금리는 단기 정책금리와 같지 않다. 장기간의 예상 단기금리, 기간 프리미엄 등이 합쳐져 형성된다.
\par

2000\textasciitilde{}2009년의 비교에서 한국 10년물 금리는 대체로 미국보다 높고 일본은 훨씬 낮은 구간에 있다. 2008년 전후에는 미국 금리가 급격히 내려가는 모습도 보인다. 한 시점의 국가 간 금리 격차를 곧바로 '추가 이익'이라고 볼 수 없는 이유는 투자 통화와 기대 환율, 물가, 위험 프리미엄이 다르기 때문이다. 원화 투자자가 달러채를 사면 달러 표시 금리와 원/달러 변동을 함께 받는다.
\par

\paragraph{장기금리를 한 숫자로 읽지 않기}

10년물 수익률은 앞으로 10년 동안 매년 지급될 돈을 현재가치로 만드는 가격의 요약치다. 오늘의 정책금리 변화가 장기금리에 미치는 영향은 앞으로의 단기금리 예상과 기간 프리미엄에 달려 있다. 장기금리에는 미래 금리 경로에 대한 예상뿐 아니라 긴 기간 동안 자금을 묶어 두는 위험과 수급의 영향도 담길 수 있다.
\par

원화 국채와 달러 국채의 표시 수익률을 비교하는 예를 생각해 보자. 원화 투자자가 달러 국채를 사려면 현재 원화를 달러로 바꾸고, 만기에는 달러를 원화로 바꾼다. 미래 환율이 불확실하면 달러 채권의 원화 수익도 불확실하다. 환헤지를 한다면 선도환율과 헤지 비용을 포함해야 하므로 \textbf{표시 금리 차이만으로 차익}을 결론 내릴 수 없다.
\par

\paragraph{만기까지 들고 있으면 금리 상승은 무관한가?}

만기에 받을 명목 원금이 확정되고 실제로 만기까지 보유한다면 중간의 시가 하락을 매각 손실로 실현하지 않을 수 있다. 그러나 갑자기 현금이 필요해 매도하거나 채권을 담보로 맡겼다면 가격 변화가 바로 중요하다. 쿠폰을 재투자하는 금리도 달라진다. 따라서 \textbf{중간 시가 위험과 만기 현금흐름 위험}은 구분하되 둘 다 투자 계획에 포함한다.
\par

\subsection{슬라이드 13: 파생상품의 핵심 역할}

선도·선물은 미래 가격을 오늘 정하고, 옵션은 정해진 가격에 사고팔 수 있는 권리를 제공한다. 콜옵션의 만기 지급액은 $\max(S_T-K,0)$, 풋옵션은 $\max(K-S_T,0)$다. 가격은 미래 주가의 단순한 평균 예측만으로 정해지지 않고, 복제 포트폴리오와 무차익 조건에 의해 제약된다.
\par

국내의 코스피200 선물은 지수 노출을 조절하는 구체적 예다. 주식 포트폴리오를 팔지 않고 선물의 매도 포지션으로 시장 하락에 대한 민감도를 줄일 수 있다. 다만 헤지 대상과 선물 기초지수의 차이, 증거금, 만기 교체 비용이 남는다.
\par

\paragraph{지급액과 순손익을 구별하기}

행사가격 $K=100$인 콜옵션을 프리미엄 8에 샀다면 만기 주가 $S_T=120$에서 지급액은 20이지만 프리미엄을 뺀 만기 단순 순손익은 12다. $S_T=80$이면 행사하지 않아 지급액은 0이고 프리미엄 8을 잃는다. \textbf{지급액} $\max(S_T-K,0)$과 \textbf{초기 비용을 반영한 손익}을 구분해야 손익분기점을 정확히 읽을 수 있다.
\par

선물 매도 헤지는 보유 주식의 시장 하락 손실을 일부 상쇄할 수 있지만, 주식 포트폴리오와 선물의 기초지수가 정확히 같이 움직인다는 보장은 없다. 헤지를 설계할 때는 계약 승수, 포트폴리오의 지수 민감도, 만기, 증거금 유동성을 함께 정한다.
\par

\paragraph{선물 계약에 큰 초기 매입대금이 없으면 위험도 작은가?}

코스피200 선물은 지수 변화에 따른 손익을 주는 계약이다. 처음에 기초 주식 전체의 매입대금을 내지 않아도 지수가 반대 방향으로 움직이면 계약 규모에 비례한 손실이 난다. 증거금은 거래 이행을 위한 담보이며, 손실 규모는 기초지수의 변동과 계약 규모에 달려 있다. 주식 1억원의 시장 노출을 줄이려고 선물을 매도했다면 계약 수량과 기존 포트폴리오의 지수 민감도를 함께 계산해야 한다.
\par

\subsection{슬라이드 14: 소비와 저축에서 투자로}

투자 결정은 오늘의 소비를 포기하고 불확실한 미래 소비를 얻는 선택이다. 내일 110을 확실히 받기 위해 오늘 100을 내면 명목 수익률은 10\%다. 그러나 미래 물가가 10\% 오르면 실질 소비량은 거의 늘지 않는다. 불확실한 110이라면 손실 상태에서의 소비 감소까지 고려해야 한다.
\par

투자의 기준은 필요한 시점에 필요한 소비를 마련할 가능성이다. 기존 소득과 부채, 투자기간을 함께 놓으면 같은 자산도 투자자마다 다른 역할을 한다.
\par

\paragraph{투자기간이 위험의 뜻을 바꾼다}

1년 뒤 확정적으로 100만원을 내야 하는 사람에게 중요한 위험은 그날 100만원을 마련하지 못하는 것이다. 장기 평균수익률이 높은 자산이라도 지급 직전 가격이 하락할 수 있으면 목적에 맞지 않는다. 반면 30년 뒤의 지출을 위해 저축하는 사람은 단기 가격 변동을 감수할 여지가 있지만, 장기간의 물가 상승과 지속적인 저축 부족도 위험이다.
\par

명목 원금의 안전과 실질 구매력의 안전도 다르다. 1년 뒤 105만원을 확정적으로 받더라도 물가가 10\% 오르면 오늘의 100만원으로 살 수 있던 물건보다 적게 살 수 있다. 따라서 \textbf{목표 지출의 날짜와 단위}를 먼저 정한 뒤 수익률을 평가한다.
\par

\paragraph{물가 위험을 줄이는 채권도 있나?}

미국의 물가연동국채 TIPS는 소비자물가 지표에 따라 원금을 조정하는 구조의 예다. 명목 지급액이 고정된 채권과 달리 물가가 오르면 조정 원금이 커질 수 있다. 그렇다고 원화 투자자에게 완전한 실질 무위험자산은 아니다. 달러·원 환율, 매입가격, 투자기간 중 금리 변화가 남기 때문이다. \textbf{어떤 물가와 어떤 통화의 구매력}을 보호하는지 확인해야 한다.
\par

\subsection{슬라이드 15: 다섯 가지 기본 질문}

투자를 해석할 때 다음 순서를 쓴다. \textbf{무엇을 사는가 → 현금흐름은 어디서 나오는가 → 언제 지급되는가 → 어떤 위험을 지는가 → 다른 자산과 함께 보유하면 어떻게 달라지는가.} 같은 자산도 단독 보유와 포트폴리오 안에서의 위험이 다르다.
\par

\paragraph{다섯 질문을 한 상품에 적용하기}

예를 들어 3년 만기 회사채를 산다면 무엇을 사는가는 발행 기업의 채무 청구권, 현금흐름은 쿠폰과 만기 원금, 지급 시점은 약정된 이자일과 만기다. 위험에는 기업의 부도, 금리 상승에 따른 중간 매도가격 하락, 거래가 잘 안 되는 유동성 위험이 있다. 마지막 질문은 이 채권이 기존 주식이나 부채와 함께 보유될 때 전체 위험을 늘리는지 줄이는지다.
\par

이 순서는 높은 표시 수익률에만 눈이 가는 실수를 막는다. 두 회사채의 만기수익률이 같아도 한 채권이 특정 산업과 기존 노동소득에 함께 노출된다면 투자자 개인에게는 더 위험할 수 있다. \textbf{상품 자체의 수익률과 전체 가계의 위험}을 구분해야 한다.
\par

\paragraph{같은 회사채가 사람마다 다른 위험인 이유는?}

자동차 회사에 재직하는 사람은 급여와 고용이 자동차 경기의 영향을 받는다. 그 회사의 주식이나 채권을 많이 보유하면 경기 악화 시 노동소득과 금융자산이 함께 나빠질 수 있다. 다른 산업의 투자자는 같은 채권을 자신의 기존 소득과 덜 연결된 자산으로 볼 수 있다. 표시 수익률은 같아도 \textbf{전체 가계 현금흐름의 공변동}이 다르다.
\par

\subsection{슬라이드 16: 투자와 청구권}

실물자산은 토지·건물·기계·지식처럼 재화와 서비스를 생산한다. 금융자산은 그 생산에서 생기는 현금흐름에 대한 \textbf{청구권}이다. 기업의 주식을 산다고 공장을 하나 더 만드는 것은 아니지만, 발행시장에서 기업이 자금을 조달해 공장을 만들 수 있고 유통시장은 투자자의 자금 회수를 쉽게 해 발행을 돕는다.
\par

경제 전체로 보면 금융청구권의 한쪽 자산은 다른 쪽 부채 또는 지분이다. 그래서 금융자산의 가치가 모든 주체의 순부를 더해 무한히 증가시키는 것은 아니다. 최종적으로 소비 가능한 재화·서비스를 만들어 내는 실물자산과 생산성이 부의 기반이다.
\par

\paragraph{유통시장이 발행시장에 주는 간접 효과}

기업이 신주 100만 주를 주당 1만원에 발행하면 발행시장에서 조달하는 총액은 비용을 무시할 때 100억원이다. 그다음 기존 주주가 다른 투자자에게 같은 주식을 판 대금은 기업이 다시 받는 돈이 아니다. 그러나 유통시장에서 쉽게 매도할 수 있다면 새 주식을 인수하는 투자자가 요구하는 유동성 보상이 작아질 수 있어 미래 발행 조건에 영향을 준다.
\par

주가는 '기업이 보유한 현금'과도 다르다. 발행 뒤 주가가 오르면 보유 주주의 지분가치가 바뀌지만 기업 금고에 같은 액수의 현금이 자동으로 들어오지 않는다. 추가 발행, 자사주 매각, 담보 조달처럼 별도의 거래가 있어야 현금흐름이 바뀐다.
\par

\paragraph{자사주 매각은 누구에게 현금이 들어오나?}

기업이 보유한 자사주 10만 주를 주당 1만2천원에 팔면 비용을 무시한 매각대금 12억원이 기업으로 들어온다. 같은 가격에 기존 주주가 10만 주를 다른 투자자에게 팔면 12억원은 매도 주주가 받는다. 화면에 보이는 거래가격과 주식 수가 같아도 \textbf{누가 매도인인지}에 따라 기업의 현금흐름이 달라진다.
\par

\paragraph{발행시장과 유통시장은 어떻게 다른가?}

발행시장에서는 기업이 새 주식이나 채권을 팔아 자금을 받는다. 유통시장에서는 이미 발행된 증권을 투자자끼리 사고팔기 때문에 보통 그 거래대금이 기업에 직접 들어가지는 않는다. 그렇지만 유통시장이 활발하면 투자자가 나중에 증권을 처분하기 쉬워지고, 신규 발행을 받아들일 때 요구하는 유동성 프리미엄이 줄어 기업의 자금 조달에도 영향을 준다.
\par

기업이 주식을 새로 발행하면 기존 주주의 지분율이 희석될 수 있지만, 조달자금으로 생산성이 높은 사업을 진행하면 기업 전체 가치가 늘 수도 있다. 따라서 '신주 발행 = 기존 주주에게 무조건 손해'라는 결론도 성립하지 않는다.
\par

\subsection{슬라이드 17: 재무상태표로 이해하는 실물·금융자산}

가계가 기업 주식 100을 보유하면 가계에는 금융자산 100, 기업에는 지분 청구권 100이 대응한다. 기업이 그 자금으로 설비를 사면 기업에는 실물자산이 남는다. 가계와 기업을 합친 통합 재무상태표에서는 내부 금융청구권이 상쇄되고 설비가 경제의 순실물자산으로 남는다.
\par

통합 재무상태표는 청구권을 중복 계산하지 않도록 돕는다. 주가 변화는 미래 이익 기대나 할인율의 변화를 반영하며 투자·담보·소비 결정에도 영향을 준다. 따라서 \textbf{청구권의 회계적 상쇄}와 \textbf{생산능력의 시장가치 변화}를 함께 살펴야 한다.
\par

\paragraph{통합 재무상태표의 계산}

가계가 현금 100을 새 회사채에 투자하고 기업이 받은 100으로 기계를 구입했다고 하자. 가계에는 회사채 자산 100, 기업에는 기계 자산 100과 회사채 부채 100이 있다. 두 주체를 합치면 회사채 자산과 부채가 서로 지워지고 기계가 남는다. \textbf{금융자산 100과 기계 100을 모두 경제 전체의 순자산으로 더하면 100을 중복 계산}한다.
\par

기업의 생산성이 올라 기계가 창출할 미래 현금흐름의 가치가 150으로 재평가될 수는 있다. 경제적 부의 변화는 예상 생산력이나 할인율의 변화에서 나온다. 회계상의 상쇄와 시장가치 변화의 경제적 효과를 함께 이해해야 한다.
\par

\paragraph{할인율만 바뀌어도 지분가치가 달라지나?}

매년 5억원의 현금흐름이 영구히 같은 수준으로 이어진다고 단순화하면, 할인율 5\%에서 현재가치는 $5/.05=100$억원이다. 미래 현금흐름 전망을 그대로 두고 할인율만 4\%로 낮추면 가치는 $5/.04=125$억원이 된다. 생산설비의 물리적 수량이 바뀌지 않아도 \textbf{미래 현금흐름을 평가하는 가격}은 달라질 수 있다. 이 청구권의 125억원과 기초 사업의 125억원을 별개의 부로 더하지는 않는다.
\par

\paragraph{금융자산은 왜 경제 전체의 순자산이 아닌가?}

가계가 기업 채권 100을 보유한다면 가계의 자산 100은 기업의 부채 100과 짝을 이룬다. 두 재무상태표를 합하면 내부 채권·채무가 사라진다. 주식도 가계의 지분자산과 기업의 자기자본 청구권이 맞물린다. 그래서 모든 경제주체의 자산만 더하고 부채를 빼지 않으면 부를 이중 계산한다.
\par

국제 거래가 있으면 외국에 대한 순채권은 국내 경제의 순금융자산이 될 수 있다. 또한 실물자산의 시장가치와 장부가치는 다를 수 있다. 통합 회계 항등식은 \textbf{누구의 청구권이 누구의 의무인지}를 확인하는 도구다. 가격 변화의 경제적 효과는 별도로 평가한다.
\par

\subsection{슬라이드 18: 투자관리의 다섯 단계}

먼저 미래 지급 목표와 투자 기간을 정한다. 다음으로 허용 손실, 유동성, 세금, 규제, 투자 가능 자산을 정책으로 만든다. 그 정책 안에서 수동·능동·구조화 전략을 선택하고, 기대수익과 공분산을 이용해 실제 비중과 종목을 정한다. 마지막으로 기준지수와 위험을 감안한 성과를 평가한다.
\par

예를 들어 1년 뒤 등록금이 필요한 자금은 30년 뒤 은퇴자금과 같은 포트폴리오가 적절하지 않다. 기대수익이 높아도 등록금 지급 직전의 손실을 감당할 수 없다면 목표에 맞지 않는다. \textbf{목표 → 제약 → 자산배분}의 순서를 거꾸로 하지 않는 것이 중요하다.
\par

\paragraph{목표에서 평가까지 한 번에 연결하기}

투자정책은 추상적인 '안전하게 운용'보다 날짜와 금액으로 적을 때 실행 가능하다. 예를 들어 2년 뒤 등록금 2,000만원을 마련해야 하고 이미 1,800만원이 있다면, 첫 질문은 목표를 달성할 확률과 부족액의 허용 범위다. 그다음 투자기간, 유동성 필요, 허용 가능한 손실을 적고 자산군 비중을 정한다. 종목 선택은 이 정책 안에서 이루어진다.
\par

평가도 목표와 같은 단위로 해야 한다. 1년 수익률이 기준지수를 이겼더라도 2년 뒤 필요한 현금을 확보하지 못했다면 목표 달성 관점에서는 성공이라 보기 어렵다. 반대로 시장이 크게 오를 때 덜 올랐더라도 지급 의무를 안정적으로 충족했다면 정책에는 충실할 수 있다.
\par

\paragraph{확정 지출을 위한 부족액은 얼마인가?}

1년 뒤 등록금 2,000만원이 필요하고 현재 1,900만원이 있다면 필요한 확정 수익률은 $2{,}000/1{,}900-1\simeq5.26\%$다. 같은 기간의 확정 수익률이 4\%라면, 위험한 자산에서 차이를 메우려 하기 전에 오늘 추가로 저축할 금액을 계산할 수 있다. $2{,}000/1.04-1{,}900\simeq23.08$만원을 더 넣으면 4\% 수익으로 목표 금액에 도달한다. 목표 시점이 정해진 지출에서는 \textbf{투자와 저축을 함께 설계}한다.
\par

\subsection{슬라이드 19: 주식·채권·파생상품의 현금흐름}

주식은 기업의 잔여이익과 청산가치에 대한 지분이고 만기가 정해져 있지 않다. 채권은 원금과 이자 등의 지급이 계약으로 약속되지만 발행자의 부도 위험은 남는다. 파생상품은 주가·금리·환율·지수 등 기초변수에 따라 지급액이 달라진다.
\par

만기 1년 이하의 단기채는 통상 단기 금융시장, 장기채와 주식은 자본시장에서 논의된다. 그러나 '고정수익'이란 가격이 고정됐다는 뜻이 아니다. 고정 쿠폰 채권도 시장금리가 바뀌면 현재 가격이 바뀐다.
\par

\paragraph{같은 이름의 자산도 다른 손익을 만든다}

액면 100을 만기에 지급하는 국채와 회사채는 약속한 숫자가 같아도 실제 지급 가능성이 다르다. 주식은 채권자에게 지급한 뒤 남는 가치에 대한 청구권이므로 상방이 열려 있지만 큰 손실도 먼저 흡수한다. 파생상품의 경우 미래 지급액이 다른 가격이나 사건의 함수로 정해져, 계약 규모에 비해 초기 지출이 작을 수 있다.
\par

현금흐름을 볼 때는 \textbf{계약상 약속}, \textbf{실제로 실현될 수 있는 금액}, \textbf{오늘 지급하는 가격}을 따로 쓴다. 고정 쿠폰 5를 지급하는 채권을 90에 샀다면 쿠폰율은 액면 대비 5\%일 수 있어도 투자자가 얻는 만기수익률은 5\%와 다르다. 이 구별은 다음 장의 수익률 계산으로 이어진다.
\par

\paragraph{'고정수익' 채권의 가격이 왜 움직이나?}

1년 뒤 105를 주는 채권은 시장에서 요구하는 수익률이 5\%면 가격이 100이다. 같은 지급액에 요구수익률이 10\%가 되면 가격은 $105/1.10\simeq95.45$로 떨어진다. 약속된 현금흐름은 그대로지만 \textbf{오늘 그 현금흐름과 교환하는 가격}이 달라진 것이다. 채권의 고정 지급 약속과 시장가격의 고정은 같은 말이 아니다.
\par

\paragraph{파생상품이 모두 투기 상품인가?}

농산물 생산자가 가격 하락을 걱정해 선물을 매도하면 기존 생산물의 가격 위험을 줄이는 헤지다. 반대로 기초자산 없이 같은 선물을 매도하면 가격 하락에 새로 노출되는 포지션이다. 계약은 같아도 투자자의 기존 자산·부채에 따라 전체 위험에 미치는 효과가 다르다.
\par

옵션 매수자는 프리미엄을 지불하고 비대칭 지급을 얻는다. 옵션 매도자는 프리미엄을 받는 대신 특정 상태의 큰 지급 의무를 질 수 있다. 따라서 상품 이름보다 만기 지급액과 기존 포트폴리오를 함께 봐야 한다.
\par

\subsection{슬라이드 20: 기관투자자의 서로 다른 제약}

은행·보험사·연기금·기금·투자회사·헤지펀드·사모펀드는 모두 큰 자금을 운용하지만 목표가 다르다. 보험사는 보험금 지급과 만기·금리 민감도를 맞추고, 연기금은 장기 지급 의무를 고려한다. 은행은 예금 인출과 규제 자본을, 헤지펀드는 계약상 레버리지와 환매 조건을 고려한다.
\par

같은 수익률 예측을 가지고도 제약이 다르면 최적 비중이 달라진다. 다음 장의 평균-분산 해는 이런 현실 제약을 제거한 \textbf{기준 모형}이며, 실제 운용에서는 유동성·공매도·보유한도·거래비용을 추가한다.
\par

\paragraph{부채가 투자 전략을 바꾸는 예}

보험사가 10년 뒤 확정적인 보험금 100을 지급해야 한다고 하자. 자산만 보아 장기채의 가격 변동이 커 보일 수 있지만, 금리가 하락할 때 보험금 부채의 현재가치도 커진다. 자산과 부채의 금리 민감도를 비슷하게 맞추면 순자산의 변동은 오히려 줄어들 수 있다. 주식 비중이 같은 두 기관이라도 지급 약속의 시점이 다르면 적절한 채권 만기가 달라진다.
\par

은행은 예금자의 단기 인출 요구도 고려해야 하므로 기대수익이 높은 장기 비유동 자산을 무제한 보유할 수 없다. 평균-분산 그래프의 '효율적'이라는 표현은 먼저 정한 투자 가능 집합 안에서만 성립한다. 규제와 유동성 제약을 바꾸면 선택 가능한 점 자체가 달라진다.
\par

\paragraph{자산과 부채의 금리 민감도가 다르면?}

연기금의 자산 현재가치가 100억원, 부채 현재가치가 90억원이라고 하자. 자산의 수정듀레이션이 5년, 부채가 10년이면 금리가 1\%포인트 하락할 때 가치변화의 1차 근사는 자산 $+5$억원, 부채 $+9$억원이다. 순자산은 10억원에서 약 6억원으로 줄어든다. 자산 가격이 올랐다는 사실만으로 재무 상태가 개선된 것은 아니다. \textbf{부채와 함께 측정한 금리 노출}이 중요한 이유다.
\par

\subsection{슬라이드 21: 기관 규모 표를 읽을 때}

운용자산 규모(AUM)는 수익률이나 경제적 순부와 같지 않다. 한 고객의 돈이 펀드와 재간접펀드에 중복 집계될 수 있고, 시장가격 상승만으로도 AUM이 증가한다. 규모가 큰 운용사는 거래비용을 낮추기도 하지만 특정 시장에서는 주문 자체가 가격을 움직일 수 있다.
\par

기관 순위를 볼 때는 기준 연도, 환율, 운용과 수탁의 구분, 중복 집계, 법인 단위와 그룹 단위를 확인한다. 순위보다 중요한 질문은 기관이 어떤 의무를 지고 어떤 위험을 흡수하는가이다.
\par

\paragraph{AUM 변화의 세 부분}

운용자산이 연초 100억원, 한 해 투자수익률이 10\%이고 신규 고객 순유입이 연말에 20억원이라면 연말 AUM은 단순화하면 130억원이다. 30억원 증가 중 기존 자산의 운용 성과는 10억원, 고객 자금 유입은 20억원이다. 유입 시점이 연중이라면 새 돈에도 수익이 붙으므로 현금흐름 날짜가 필요하다. AUM 증가율 30\%를 운용수익률 30\%라고 읽을 수 없는 이유다.
\par

규모가 크면 주문을 여러 번 나누거나 거래량이 풍부한 자산을 선택해야 할 수 있다. 동일한 전략도 소규모에서는 가능하지만 대규모에서는 가격을 움직여 수익이 줄 수 있다. 규모 순위와 운용 능력은 별개의 지표다.
\par

\paragraph{큰 펀드가 작은 시장에서 매매하면 어떤 비용이 생기나?}

매일 5억원어치 거래되는 종목을 대형 펀드가 하루에 20억원 사려고 하면 매수 주문이 호가를 밀어 올릴 수 있다. 예상한 가격보다 비싸게 체결된 차이가 \textbf{시장충격(market impact)}이다. 나중에 같은 물량을 팔 때도 가격 하락을 유발할 수 있다. AUM이 커지면 고정비를 낮출 여지가 있지만 거래 규모와 유동성의 관계가 전략의 실현 수익률을 좌우한다.
\par

\subsection{슬라이드 22: 다음 단계: 수익률을 확률변수로 보기}

투자 대상의 차이를 알았으니 이제 수익률을 동일한 보유기간의 확률변수로 정의한다. 기대수익은 가능한 결과의 평균, 분산은 평균 주위의 흩어짐, 공분산은 두 자산이 함께 움직이는 정도다. 이 세 도구를 익히면 분산투자 효과와 효율적 포트폴리오가 계산으로 연결된다.
\par

핵심 연결식은 $R_p=\sum_i w_iR_i$이다. 기대수익은 비중의 가중합으로, 위험은 각 자산의 분산과 자산쌍의 \textbf{공분산}으로 계산된다. 다음 장에서 이 입력값을 측정한 뒤 최적 비중을 구한다.
\par

\paragraph{공분산이 필요한 이유를 미리 보기}

자산 A와 B가 호황·불황에서 각각 $(+10\%,-10\%)$와 $(-10\%,+10\%)$의 수익률을 내고 두 상태의 확률이 같다고 하자. 각각의 평균은 0이고 변동은 크지만 절반씩 보유한 포트폴리오는 두 상태에서 모두 0\%다. 반대 방향의 움직임이 서로 상쇄되기 때문이다. 이 예의 공분산은 음수이며 두 자산의 개별 표준편차만 더해서는 전체 위험 0을 설명할 수 없다.
\par

반대로 두 자산이 매번 같은 방향으로 움직이면 절반씩 보유해도 같은 공통 위험이 남는다. 다음 장에서는 이 '함께 움직임'을 $\operatorname{Cov}(R_A,R_B)$로 측정하고, 세 번째 장에서는 비중을 바꿨을 때 위험이 어떻게 변하는지 계산한다.
\par

\paragraph{자산을 절반씩 나누면 위험도 절반이 되나?}

두 자산의 표준편차가 모두 10\%이고 상관이 1이면 절반씩 나누어도 포트폴리오 표준편차는 10\%다. 같은 충격을 함께 받기 때문이다. 상관이 0이라면 표준편차는 $10\%/\sqrt2\simeq7.07\%$, 상관이 -1이면 같은 비중에서 0이 될 수 있다. \textbf{함께 움직이는 정도}가 분산투자 효과를 정한다.
\par

\clearpage

\section{2부 · 수익률과 위험}

\subsection{슬라이드 1: 위험과 수익률의 공통 언어}

서로 가격도 지급 방식도 다른 자산을 비교하려면 투자한 돈 한 단위가 기간 말에 얼마나 되었는지 측정해야 한다. 그 비율이 \textbf{보유기간 수익률}이다. 미래 수익률은 아직 확정되지 않았으므로 확률변수 $R$로 나타내고, 기대값과 분산으로 분포의 일부 특징을 요약한다.
\par

\textbf{변동성}은 예상에서 벗어난 정도를 상승·하락 양쪽에서 측정한다. \textbf{손실 위험}은 결과의 왼쪽 꼬리에 초점을 맞춘다. 두 기준을 구분해 두면 뒤의 표준편차와 VaR가 답하는 질문이 분명해진다.
\par

\paragraph{수익률이 서로 다른 가격을 비교하게 하는 이유}

1만원짜리 주식이 1,100원 오르고 10만원짜리 채권이 5,000원 올랐다면 원화 이익만으로 투자 효율을 비교할 수 없다. 최초 투입금으로 나누면 각각 11\%와 5\%다. 다만 주식의 배당과 채권의 쿠폰을 빼면 총수익률이 아니며, 두 투자기간이 다르면 기간도 맞춰야 한다. 수익률은 \textbf{금액을 비율로 표준화}하지만 위험과 현금흐름의 시점까지 자동으로 맞추지는 않는다.
\par

미래 수익률을 확률변수로 본다는 말은 결과가 아직 여러 값 중 하나로 결정될 수 있다는 뜻이다. 가능한 수익률의 집합과 확률을 정하면 평균과 분산을 계산할 수 있다. 손실 가능성이 같은 두 투자라도 한쪽은 작은 손실이 자주, 다른 쪽은 큰 손실이 드물게 나타날 수 있으므로 평균·표준편차와 손실 분위수를 함께 읽는다.
\par

\paragraph{10\% 손실 뒤 10\% 수익이면 원금을 되찾나?}

100만원이 10\% 하락하면 90만원이고, 여기서 10\% 상승하면 99만원이다. 두 수익률은 각각 \textbf{그 시점의 투자금}을 분모로 삼는다. 90만원에서 100만원으로 돌아가려면 $10/90=11.11\%$의 수익이 필요하다. 기간별 수익률을 더하는 계산과 실제 재산이 늘어나는 복리 계산을 구분해야 한다.
\par

\subsection{슬라이드 2: 역사적 자료와 확률모형}

과거 수익률 표본은 미래를 위한 정보이지만 미래의 분포 그 자체는 아니다. 경제체제, 통화정책, 산업 구성, 거래비용이 바뀔 수 있기 때문이다. 과거 표본에서 계산한 평균과 공분산을 미래의 모수 추정치로 쓸 때에는 표본 기간, 생존편향, 인플레이션과 환율을 점검해야 한다.
\par

확률모형에서는 가능한 상태 $s$의 확률 $h(s)$와 상태별 수익률 $R(s)$를 먼저 정한다. 역사적 추정에서는 관측 시점 $t=1,\ldots,T$를 표본으로 취급한다. 상태확률과 표본 빈도를 같은 기호처럼 사용하면 모집단의 모수와 추정량을 혼동하기 쉽다.
\par

\paragraph{관측값과 가능한 상태를 구별하기}

가능한 상태가 호황·불황이고 각각의 발생확률이 주어졌다면 확률가중 평균은 모형 안에서 정의한 \textbf{모집단 평균}이다. 반면 과거 10년의 연간 수익률 열 개는 10개의 관측값이며, 그 산술평균은 알려지지 않은 미래 평균의 추정치다. 과거에 호황이 7번 있었다고 해서 미래 호황확률이 정확히 70\%라는 결론은 나오지 않는다.
\par

같은 분포에서 독립적으로 뽑혔다는 가정은 계산을 간단하게 하지만 금융자료에서는 점검 대상이다. 기업의 사업 구조, 거래제도, 물가 체계가 바뀌면 오래된 관측치가 지금의 분포를 잘 대표하지 않을 수 있다. 특히 위기 이전 자료만으로 상관을 추정하면 위기 때 공동 하락하는 정도를 과소평가할 수 있다.
\par

\paragraph{과거 10년 중 7년 상승했다면 내년 상승 확률은 70\%인가?}

관측된 상승 빈도 $7/10$은 그 10년의 \textbf{표본 비율}이다. 내년의 확률이 70\%라고 말하려면, 자료가 같은 환경에서 반복 추출되었고 미래도 같은 확률분포를 따른다는 가정이 더 필요하다. 정책, 산업 구성, 금리 체계가 달라지면 그 가정이 깨질 수 있다. 과거 빈도는 확률모형을 추정할 자료이지 미래에 대한 보증이 아니다.
\par

\subsection{슬라이드 3: 측정 순서}

첫째, 배당과 이자를 포함해 기간 수익률을 계산한다. 둘째, 단위가 같은 수익률들의 평균과 표준편차를 구한다. 셋째, 두 자산을 함께 보유할 계획이면 공분산과 상관계수를 본다. 개별 자산의 표준편차만으로는 결합된 포트폴리오의 위험을 알 수 없다.
\par

마지막으로 분포의 비대칭과 두꺼운 꼬리를 확인한다. 평균·분산은 정규분포에서 분포 전체를 결정하지만, 실제 손실분포가 비정규적이면 같은 평균과 표준편차를 가진 두 투자가 전혀 다른 극단손실을 가질 수 있다.
\par

\paragraph{단위부터 맞춘 다음 통계량을 구하기}

주식 A의 월수익률과 채권 B의 연수익률을 한 행에 두고 공분산을 계산하면, 같은 시점에 무엇이 함께 움직였는지 정의되지 않는다. 먼저 두 자산의 수익률을 같은 시작일과 종료일, 같은 통화와 배당 처리 기준으로 맞춘다. 그다음 각 열의 평균·표준편차와 두 열의 공분산을 구한다. \textbf{정렬된 두 수익률 쌍}이 있어야 공분산이 의미를 가진다.
\par

수익률이 10\%, 0\%, -10\%인 세 기간의 평균은 0\%지만 순서와 복리 성장은 별도의 정보다. 처음 100이 110, 110, 99로 변하므로 최종 부는 시작보다 1\% 적다. 통계량의 계산 순서는 분석 질문—한 기간의 불확실성을 묻는지, 여러 기간의 부를 묻는지—에 맞춰야 한다.
\par

\paragraph{해외자산 수익률을 원화로 보면?}

같은 한 달 동안 미국 ETF의 달러 가격이 5\% 오르고 달러/원 환율이 4\% 하락했다고 하자. 환헤지를 하지 않은 원화 투자자의 수익률은 $1.05\times.96-1=.008$, 곧 0.8\%다. 달러 수익률 5\%를 국내 주식의 원화 수익률과 바로 나란히 놓으면 통화 노출이 빠진다. \textbf{동일한 시작일·종료일과 같은 통화}로 각 자산의 수익률 열을 만든 뒤 공분산을 계산한다.
\par

\subsection{슬라이드 4: 보유기간 수익률을 현금흐름에서 유도}

기초 가격이 $P_{t-1}$, 기말 가격이 $P_t$, 기간 중 받은 현금배당이 $D_t$라면 기말 총가치는 $P_t+D_t$다. 순이익은 $P_t+D_t-P_{t-1}$이므로 최초 투입금으로 나누면 다음과 같다.
\par

\[
R_t=\frac{P_t-P_{t-1}+D_t}{P_{t-1}}=\frac{P_t+D_t}{P_{t-1}}-1.
\]

주당 500달러에 100주를 사서 525달러에 팔고 배당이 없으면 순이익은 $100(525-500)=2{,}500$달러, 최초 투자액은 $50{,}000$달러다. 따라서 수익률은 $2{,}500/50{,}000=5\%$다. 주당 이익은 \textbf{25달러}, 최초 투자액 대비 수익률은 \textbf{5\%}다. 매매수수료와 세금을 넣으면 실현 수익률은 낮아진다.
\par

\paragraph{가격 변화와 현금 지급을 한 분자로 묶기}

처음 100에 산 주식이 끝에 106이고 중간 배당 4를 받았다면 가격수익률은 6\%, 배당수익률은 4\%, 단순 보유기간 총수익률은 10\%다. 식으로는 $R=(106-100)/100+4/100=.10$이다. 배당을 지급받은 날 다시 투자했다면 재투자 뒤 주식 수가 달라지므로 여러 기간 총수익은 지급 날짜와 재투자가격까지 알아야 정확히 계산된다.
\par

한 기간에 여러 현금 유입·유출이 생기면 단순한 기말가격 공식만으로 투자자의 실제 성과를 설명하기 어렵다. 추가 납입은 투자수익이 아니고 인출은 손실이 아니다. 이때는 현금흐름 시점을 반영하는 내부수익률과, 운용 자체의 성과를 보기 위해 외부 현금흐름을 제거하는 시간가중수익률을 목적에 따라 구분한다.
\par

\paragraph{액면분할 뒤 주가가 반토막 나면 50\% 손실인가?}

한 주가 100달러였던 주식을 2대 1로 분할하면 기존 한 주는 두 주가 된다. 분할 뒤 주가가 51달러라면 두 주의 가치는 102달러이므로, 배당이 없을 때 보유기간 수익률은 $(102-100)/100=2\%$다. 화면에 보이는 주당 가격 100→51만 비교하면 허위 손실을 만든다. 배당, 분할, 유상증자 등 \textbf{현금흐름과 보유 주식 수}를 함께 조정해야 한다.
\par

\paragraph{주식분할·배당·환율은 수익률에 어떻게 넣나?}

주식이 1주에서 2주로 분할되면 주당 가격이 절반이 되어도 투자자의 보유 가치가 절반이 된 것은 아니다. 분할 전후의 주식 수와 배당을 조정한 \textbf{총수익 지수}를 쓰면 경제적 수익률을 더 정확히 계산할 수 있다. 현금배당은 수익률 분자에 포함하며, 배당 재투자 여부를 명시해야 한다.
\par

해외자산은 현지통화 자산수익률과 환율수익률을 곱해 연결한다. 예를 들어 달러 자산이 10\% 오르고 달러가 원화 대비 5\% 오르면 원화 수익률은 $1.10\times1.05-1=15.5\%$다. 단순히 10\%+5\%=15\%로 더하면 두 효과의 교차항 0.5\%포인트를 빠뜨린다.
\par

\subsection{슬라이드 5: 분포의 네 가지 모멘트}

평균은 중심, 분산은 평균 주위의 퍼짐, 왜도는 비대칭, 첨도는 평균에서 멀리 떨어진 결과의 빈도를 설명한다. 기대값이 $\mu$, 표준편차가 $\sigma$이면 표준화된 왜도는 $E[(R-\mu)^3]/\sigma^3$, 초과첨도는 $E[(R-\mu)^4]/\sigma^4-3$이다.
\par

양의 왜도는 큰 양의 결과가 상대적으로 긴 꼬리를 만들고, 음의 왜도는 큰 손실 쪽 꼬리가 길다는 뜻이다. 초과첨도가 높으면 정규분포보다 극단값이 자주 나타날 수 있다. 그러나 첨도 한 숫자만으로 어느 쪽 꼬리가 위험한지 알 수 없어 왜도와 손실 분위수를 함께 봐야 한다.
\par

\paragraph{네 모멘트가 답하는 서로 다른 질문}

평균은 '중심이 어디인가', 분산은 '평균에서 얼마나 퍼졌는가', 왜도는 '어느 쪽 꼬리가 긴가', 첨도는 '극단값에 얼마나 민감한가'를 묻는다. 제곱을 쓰는 분산은 양쪽 편차를 모두 양수로 만들고, 세제곱을 쓰는 왜도는 편차의 부호를 남기며, 네제곱을 쓰는 첨도는 큰 편차에 매우 높은 가중치를 준다. 이 \textbf{거듭제곱의 차이}가 통계량의 역할을 만든다.
\par

표본 왜도나 첨도는 작은 자료에서 특히 불안정하다. 예를 들어 큰 손실 한 번을 표본에 추가하면 네제곱 항이 급증할 수 있다. 첨도가 높다고 모든 손실 분위수가 정확히 얼마인지 결정되는 것은 아니며, 한 숫자를 해석할 때는 분포 그림과 관측 기간을 같이 살핀다.
\par

\paragraph{표준편차가 같으면 위험도 같나?}

한 상품은 매달 ±5\% 부근에서 움직이고, 다른 상품은 대부분 1\%를 벌지만 드물게 40\%를 잃을 수 있다. 둘의 표준편차가 비슷해도 손실의 방향과 꼬리 모양은 다르다. 표준편차는 평균에서 벗어나는 정도를 \textbf{양쪽 방향 모두} 세므로, 큰 손실에 특히 민감한 투자자는 왜도와 꼬리 위험도 확인해야 한다.
\par

\subsection{슬라이드 6: 기대수익, 변동성, 손실 측도}

기대수익 $\mu=E[R]$은 가능한 수익률의 확률가중 평균이다. 분산은 $\sigma^2=E[(R-\mu)^2]$, 표준편차는 $\sigma=\sqrt{\sigma^2}$이다. 제곱 후 평균을 내므로 예상보다 높은 수익도 변동성에 기여한다. 같은 평균이어도 $\sigma$가 큰 투자는 결과가 더 넓게 퍼진다.
\par

손실만 알고 싶다면 손실변수 $L=-R$의 분위수로 VaR를 정의할 수 있다. 다만 수익률의 왼쪽 분위수 $q_\alpha(R)$를 음수로 표시하는 방식과 양의 손실액 $\mathrm{VaR}_\alpha=-q_\alpha(R)$로 표시하는 방식이 공존한다. 부호와 신뢰수준을 먼저 밝히지 않은 VaR 숫자는 해석할 수 없다.
\par

\paragraph{손실의 단위와 신뢰수준}

원금 1억원에 대한 수익률 $R=-.08$은 손실 800만원을 뜻한다. 손실비율을 $L=-R$로 정의하면 $L=.08$이고 손실액은 $1억원\times L$이다. 수익률로 구한 VaR를 원화로 바꾸려면 포지션 규모를 곱해야 한다. 여러 포지션의 VaR를 각각 구해 단순히 더하면 서로 상쇄되는 움직임을 무시할 수 있다.
\par

97.5\% VaR는 손실분포의 97.5\% 분위수이며 2.5\% 꼬리가 그보다 더 나쁠 수 있다. \textbf{확률의 방향}을 먼저 적어야 한다. 수익률 분포에서는 왼쪽 2.5\% 분위수를 찾고, 양의 손실 분포에서는 오른쪽 97.5\% 분위수를 찾는 두 표현이 같은 사건을 가리킨다.
\par

\paragraph{VaR의 퍼센트를 실제 손실 금액으로 어떻게 옮기나?}

1일 95\% VaR가 투자금의 4\%이고 포지션이 1억원이라면 손실 기준금액은 400만원이다. 같은 수익률 분포를 가정하면서 포지션을 2억원으로 늘리면 금액 VaR도 약 800만원으로 늘어난다. 단, 거래비용과 시장충격 때문에 큰 주문의 손실이 항상 비례하는 것은 아니다. \textbf{신뢰수준·기간·투자금액}을 함께 적어야 숫자가 의미를 갖는다.
\par

\subsection{슬라이드 7: 정규분포 VaR와 숫자 점검}

$R\sim N(0.10,0.15^2)$이면 2.5\% 왼쪽 분위수는 표준정규분포의 $z_{0.025}\simeq-1.96$을 이용해 계산한다.
\par

\[
q_{0.025}(R)=\mu+z_{0.025}\sigma=0.10-1.96(0.15)=-0.194.
\]

따라서 \textbf{왼쪽 2.5\% 수익률 분위수는 -19.4\%}, 같은 사건을 손실로 표시한 \textbf{97.5\% VaR는 19.4\%}다. 임계값 바깥에 남은 2.5\%의 손실 규모는 이 분위수만으로 알 수 없다.
\par

\paragraph{변동성 추정이 VaR를 바꾸는 크기}

왼쪽 꼬리확률을 $\alpha$라 하면 정규모형의 수익률 분위수는 $q_\alpha=\mu+z_\alpha\sigma$이고, 손실비율 VaR는 $-q_\alpha$다. 평균을 고정하면 $\partial\operatorname{VaR}/\partial\sigma=-z_\alpha>0$이다. 예를 들어 97.5\% 신뢰수준에서 변동성 전망이 15\%에서 20\%로 높아지면 VaR는 $1.96(.20)-.10=29.2\%$가 된다. \textbf{위험 추정치의 변화}가 보고하는 손실 경계에 직접 전달된다.
\par

이 관계는 지정한 한 보유기간의 정규분포 가정 아래 성립한다. 월별 입력으로 계산한 값을 연간 손실 한도로 바꾸려면 기간별 수익률의 결합과 상관을 다시 모델링해야 한다.
\par

\paragraph{신뢰수준을 바꾸면 VaR가 얼마나 달라지나?}

같은 정규분포 가정 $\mu=10\%,\sigma=15\%$에서 95\% VaR는 왼쪽 5\% 분위수 $z_{.05}\simeq-1.645$를 쓴다. 수익률 임계값은 $.10-1.645(.15)=-.14675$이므로 손실비율 VaR는 약 14.7\%다. 왼쪽 2.5\%를 쓰는 97.5\% VaR 19.4\%보다 작다. \textbf{더 드문 손실까지 포착하려는 높은 신뢰수준}일수록 경계가 더 멀리 이동한다.
\par

\paragraph{VaR의 부호와 확률을 다시 확인하기}

수익률 왼쪽 2.5\% 분위수가 -19.4\%라면 $P(R\le-.194)=.025$라는 뜻이다. 손실변수 $L=-R$를 쓰면 $P(L\ge.194)=.025$가 되고, 이때 양의 손실 VaR는 19.4\%다. 수익률 임계값과 손실액을 같은 이름으로 부르기도 하므로 부호가 혼란스럽다.
\par

또한 97.5\% 신뢰수준의 하루 VaR 19.4\%는 매일 최대 손실이 19.4\%라는 뜻이 아니다. 모형이 맞고 관측일이 충분하다면 약 40일에 하루꼴로 이 기준을 넘는 손실이 나타날 수 있다는 확률적 설명이다. 실제 금융위기에는 정규분포보다 더 자주 발생할 수 있다.
\par

\subsection{슬라이드 8: 정규분포 가정의 힘과 한계}

정규분포에서는 평균과 분산만 알면 분포가 결정된다. $\mu\pm\sigma$ 사이에 약 68.3\%, $\mu\pm1.645\sigma$ 사이에 약 90\%, $\mu\pm1.96\sigma$ 사이에 약 95\%의 확률이 있다. 이는 한 번의 미래 관측치가 해당 구간에 들어갈 확률이지 \textbf{추정한 평균의 신뢰구간}이 아니다.
\par

정규분포 밖에서는 평균과 표준편차가 같아도 손실 분위수가 달라질 수 있다. \textbf{Expected Shortfall (ES)}은 VaR 경계를 넘은 손실들의 조건부 평균으로 꼬리의 크기를 요약한다. VaR가 손실의 경계라면 ES는 그 경계 바깥에서 예상되는 손실 수준이다.
\par

\paragraph{VaR 바깥의 손실을 보는 방법}

손실 $L=-R$이고 $R\sim N(\mu,\sigma^2)$라면 오른쪽 꼬리확률 $\alpha$에 대한 정규모형 ES는 $-\mu+\sigma\phi(z_{1-\alpha})/\alpha$다. 여기서 $\phi$는 표준정규 밀도다. $\mu=.10,\sigma=.15,\alpha=.025$이면 $\phi(1.96)\simeq.0584$이므로 ES는 약 $-.10+.15(.0584/.025)\simeq25.0\%$다.
\par

이 숫자는 \textbf{정규모형이 맞는 경우} 꼬리 2.5\%에 속한 손실의 평균이다. 정규분포보다 손실 꼬리가 두껍다면 실제 ES가 더 클 수 있다. 예측구간, 평균의 신뢰구간, 손실분포의 분위수는 모두 다른 확률적 대상을 설명한다.
\par

\paragraph{VaR가 같은 두 전략은 비슷하게 안전한가?}

두 전략의 95\% VaR가 모두 10\%여도 첫 전략의 기준 초과 손실은 대체로 11\%이고, 둘째는 드물게 50\%일 수 있다. VaR는 경계 \textbf{바깥 손실의 크기}를 보여주지 않는다. 그 구간의 평균 손실인 Expected Shortfall (ES)을 함께 보면 차이가 드러난다. 5\%의 나쁜 경우에 얼마를 잃는지도 투자자의 자금 조달 능력에 중요하다.
\par

\paragraph{예측구간과 평균의 신뢰구간}

$\mu\pm1.96\sigma$는 분포 모수를 알고 정규성을 가정할 때 \textbf{미래 수익률 한 번}의 가운데 약 95\%를 담는 구간이다. 반면 독립적인 관측치 $T$개로 평균을 추정할 때 $\bar R\pm1.96s/\sqrt T$는 큰 표본에서 \textbf{모평균}을 추정하는 구간이다. 뒤 구간은 표본이 커질수록 좁아지지만, 미래 한 번의 투자 결과 자체가 거의 확정되는 것은 아니다.
\par

예를 들어 연간 수익률 표준편차가 20\%이고 100년의 독립 자료가 있으면 평균 추정의 표준오차는 약 2\%지만, 1년 후 수익률의 표준편차는 여전히 20\% 수준일 수 있다.
\par

\subsection{슬라이드 9: 오래된 장기 수익률 표본 읽기}

1926년 말 1의 부에서 출발한 누적부 그래프는 1995년 무렵 미국 소형주 약 3,425, S\&P 500 약 974, 장기 회사채 약 44, 장기 국채 약 37, 단기 국채 약 13, 소비자물가 약 9의 지수 수준을 보여준다. 세로축은 \textbf{로그 눈금}이므로 1에서 10으로 올라가는 거리와 10에서 100으로 올라가는 거리가 같다. 주식의 높은 끝값만 보면 1930년대의 큰 손실과 오랜 회복 기간을 놓친다.
\par

이 그림은 매 기간의 총수익을 재투자한 \textbf{사후 복리 결과}다. 끝값에는 표본의 시작·끝 시점, 물가, 거래비용, 생존한 시장만 남는 편향이 함께 작용한다. 높은 주식 수익률은 경기 악화기에 현금흐름도 나빠지는 위험과 연결해 해석할 수 있다.
\par

\paragraph{로그축에서 기울기가 의미하는 것}

부가 매년 같은 성장률 $g$로 증가하면 $W_T=W_0(1+g)^T$이고 로그를 취하면 $\log W_T=\log W_0+T\log(1+g)$다. 따라서 로그 세로축에서 일정한 복리 성장률은 대체로 직선의 기울기로 나타난다. 지수 수준 10에서 20으로의 변화와 100에서 200으로의 변화는 모두 100\% 증가이므로 로그축에서는 같은 세로 거리다.
\par

누적부 곡선의 끝점만 비교할 때는 투자 기간 중 겪은 손실도 확인한다. 원금이 절반으로 줄면 남은 금액에서 100\% 상승해야 원래 수준으로 돌아간다. 사후 최고 수익 자산이 그 경로를 견딜 수 있었는지와 미래에도 같은 성장률이 가능한지는 별도 문제다.
\par

\paragraph{가격이 원래 수준으로 돌아오면 장기 수익률도 0인가?}

가격이 100→50→100이면 전체 기간의 가격 수익률은 0\%다. 하지만 중간 수익률은 -50\%와 +100\%로 크게 달랐고, 50에서 매도해야 했다면 회복의 혜택을 얻지 못한다. 같은 시작점과 끝점이라도 경로는 필요한 현금, 담보, 심리적 부담을 바꾼다. 장기 성과표와 별도로 \textbf{maximum drawdown(MDD, 최고점 대비 최대 하락률)}을 보는 이유다.
\par

\subsection{슬라이드 10: 수익률 분포의 꼬리}

히스토그램은 평균 주변 관측치가 얼마나 모이는지와 극단적인 상승·하락이 얼마나 자주 생기는지 보여준다. 정규 곡선과 겹쳐 볼 때 꼬리의 차이를 확인할 수 있지만, 구간 폭과 관측기간에 따라 그림이 달라진다.
\par

1926\textasciitilde{}1997년의 \textbf{월별} 수익률 표에서는 S\&P 지수의 월평균 약 1.03\%, 월표준편차 약 5.66\%, 소형주의 월평균 약 1.36\%, 월표준편차 약 8.66\%, 장기 국채의 월평균 약 0.45\%, 월표준편차 약 2.23\%가 함께 나온다. 소형주가 더 높은 평균을 보인 동시에 분포도 훨씬 넓었다는 뜻이다. 표의 소형주 최소·최대 월수익률은 약 -36.74\%와 +73.46\%로, 평균 하나가 결과의 범위를 가려 버릴 수 있음을 보여준다.
\par

왜도와 초과첨도는 다음처럼 표준화한다. $\gamma_1=E[(R-\mu)^3]/\sigma^3$, $\gamma_2=E[(R-\mu)^4]/\sigma^4-3$. 정규분포는 둘 다 0이다. 표본에서는 극단값 몇 개가 고차 모멘트를 크게 바꿀 수 있어 수치의 불확실성도 함께 판단해야 한다.
\par

\paragraph{히스토그램의 모양과 표본의 한계}

월수익률 표본 120개에서 관측되지 않은 손실도 미래에는 발생할 수 있다. 히스토그램의 양 끝은 관측된 범위만 보여 준다. 구간을 넓게 잡으면 두꺼운 꼬리가 가려지고, 너무 좁게 잡으면 우연한 표본 잡음이 패턴처럼 보일 수 있다.
\par

극단 손실의 확률을 추정할 때는 표본 빈도, 분포 가정, 경제적 시나리오를 서로 보완적으로 쓴다. 표준편차가 같은 두 자산이라도 한쪽의 수익률이 완만한 반면 다른 쪽이 가끔 크게 하락하면 투자자의 지급 목표에 미치는 영향이 다르다. 평균 주변의 흔들림과 꼬리의 손실을 구분하는 이유다.
\par

\paragraph{히스토그램 막대가 달라도 자료의 꼬리는 같나?}

월수익률 120개 중 6개가 -10\% 아래라면 관측 빈도는 $6/120=5\%$다. 막대 폭을 2\%포인트에서 5\%포인트로 바꾸면 히스토그램 모양은 달라지지만, \textbf{-10\% 아래가 6개라는 사실}은 바뀌지 않는다. 꼬리 확률을 비교할 때는 막대의 높이만 보지 말고 같은 손실 임계값 아래의 관측 개수를 직접 세어야 한다. 이 5\%도 미래 확률의 불확실한 추정치다.
\par

\subsection{슬라이드 11: 위험-수익 산점도}

월별 자료를 자산군별 막대로 비교하면 단기 국채의 월평균과 표준편차가 모두 작고, S\&P 지수와 특히 소형주는 두 수치가 모두 큰 모습이 보인다. 위험 막대와 수익 막대를 같은 축에 놓아도 둘의 \textbf{단위는 모두 월간 백분율}이며 서로 다른 통계량이다. 월평균 1\%와 월표준편차 5\%를 연간 1\%와 5\%로 읽으면 안 된다.
\par

높은 평균과 높은 변동성이 함께 보인다고 개별 자산마다 '표준편차 1\%당 기대수익 얼마'라는 보편 법칙이 성립하는 것은 아니다. 가격결정에서 보상받는 위험은 분산투자로 제거할 수 없는 부분일 수 있다. 총변동성과 시장과의 공변동은 다르며, 막대의 역사적 평균은 사후 자료이지 사전 기대수익이 아니다.
\par

\paragraph{월간 숫자를 연간 숫자로 바꾸는 조건}

월별 단순수익률의 평균을 $\mu_m$, 표준편차를 $\sigma_m$이라 하면, 월별 수익률이 독립이고 분산이 안정적이라는 근사 아래 12개월 \textbf{수익률 합계의} 평균은 $12\mu_m$, 표준편차는 $\sqrt{12}\sigma_m$이다. 실제 연간 단순수익률은 $\prod_{m=1}^{12}(1+R_m)-1$이므로 합계와 정확히 같지 않다. 평균을 12배 하고 표준편차를 12배 하는 방식도 틀리다.
\par

예를 들어 월표준편차가 5\%라면 제곱근 규칙의 연간 근사는 $5\%\sqrt{12}\simeq17.3\%$다. 월수익률이 연속적으로 같은 방향으로 움직이거나 위기 때 변동성이 커지면 공분산 항이 생겨 이 근사와 실제 연간 표준편차가 달라진다.
\par

\paragraph{월별 수익률에 양의 자기상관이 있으면?}

두 달의 월표준편차가 각각 5\%이고 수익률 공분산이 $.001$이라면, 두 달 \textbf{수익률 합계의} 분산은 $.05^2+.05^2+2(.001)=.007$이고 표준편차는 약 8.37\%다. 독립이라면 $\sqrt{2}(.05)\simeq7.07\%$다. 함께 움직이는 달을 독립으로 처리하면 위험을 낮게 추정한다. 실제 두 달 누적 단순수익률은 곱으로 계산하므로 이 비교는 합계 근사에 관한 것이다.
\par

\subsection{슬라이드 12: 과거 관계를 미래 예측으로 바꿀 때}

가로축을 월표준편차, 세로축을 월산술평균으로 바꾼 산점도에서 단기 국채는 왼쪽 아래, 장기채는 중간, S\&P 지수와 소형주는 오른쪽 위에 놓인다. 이 점들은 1926\textasciitilde{}1995년의 서로 다른 \textbf{자산군별 표본 요약}이며, 한 투자자가 그 사이의 모든 점을 반드시 거래할 수 있었다는 뜻이 아니다. 직선을 그어도 그것이 최적 포트폴리오의 자본배분선이나 회귀선이 자동으로 되는 것은 아니다.
\par

기간을 바꾸면 점의 위치와 순서도 바뀔 수 있고, 채권의 장기 수익률은 금리 하락기와 상승기에 달라진다. 투자 결정에는 같은 기간·통화·배당 재투자·세전후 기준을 맞춘 뒤 관측 평균의 불확실성을 생각해야 한다. 예를 들어 독립적인 연간 수익률 $T$개라면 표본평균의 표준오차는 대략 $\sigma/\sqrt T$다. 월별 자료로 추정했다면 월별 $\sigma$와 월별 관측 개수 $T$를 써야 한다.
\par

\paragraph{산점도의 점과 최적선은 다르다}

자산군별 과거 평균과 표준편차를 한 점씩 그린 그림은 \textbf{관측된 자산군의 성과 요약}이다. 두 자산을 실제로 섞어 만들 수 있는 모든 포트폴리오 점은 공분산이 있어야 계산된다. 자산 A와 B의 점을 직선으로 이어 그린다고 가운데 점들이 그 직선에 놓인다는 보장은 없다. 평균은 비중에 따라 선형으로 움직여도 표준편차는 보통 비선형이기 때문이다.
\par

표본평균의 오차도 큰 편이다. 독립된 월 관측치 120개, 월표준편차 5\%라면 평균의 표준오차는 대략 $5\%/\sqrt{120}\simeq0.46$\%포인트다. 표본 월평균 1\%와 1.3\%의 차이를 발견해도 그 차이가 미래의 확정적 우위를 뜻한다고 볼 수 없다.
\par

\paragraph{산점도의 오른쪽 위 점은 무조건 더 좋은 투자처인가?}

어떤 펀드 A의 기대수익률·변동성이 8\%·12\%, B가 10\%·12\%라면 같은 위험에서 B가 우월하다. 반면 C가 11\%·20\%라면 B와 C의 선호는 위험 회피 정도에 달려 있다. 그래프의 모든 점이 \textbf{효율적 투자기회}는 아니다. 숫자가 과거 추정치라면 그 우월 관계 역시 추정 오차의 영향을 받는다.
\par

\subsection{슬라이드 13: 주식 변동성의 시간 변화}

왼쪽 그래프는 S\&P 지수의 \textbf{월별 수익률}을 한 달씩 보여주며 1930년대에 특히 큰 양·음의 움직임이 눈에 띈다. 오른쪽은 같은 자료의 \textbf{120개월 이동평균(moving average)}이다. 120개월은 약 10년이므로 한 달의 급등락이 평균 과정에서 완화된다. 오른쪽 세로축도 수익률이므로 이를 '10년 변동성'이나 10년 동안 손실이 사라졌다는 증거로 읽어서는 안 된다.
\par

별도로 변동성을 측정하려면 각 120개월 구간의 \textbf{표준편차}를 계산해야 한다. 주식의 변동성 자체도 고정 상수가 아니며 조용한 시기와 큰 움직임이 몰리는 시기가 있다. 일별 표준편차를 $\sqrt{252}$로 연율화하는 관행은 일별 수익률이 서로 상관되지 않고 분산이 안정적일 때 성립한다. 자기상관과 변동성 군집이 있으면 실제 연간 위험과 다를 수 있다.
\par

\paragraph{이동평균과 위험 축소를 구분하기}

월수익률을 120개월 이동평균(moving average)으로 바꾸면 인접한 두 이동평균은 119개월 자료를 공유한다. 그래서 그래프가 매끄럽게 보이는 것은 당연하다. 원자료의 큰 손실도 120개월 평균에 나누어 반영된다. 이동평균 계열의 표준편차를 원래 월수익률의 표준편차와 비교하면 서로 다른 변수를 비교하게 된다.
\par

변동성이 시간에 따라 달라지는지를 보려면 고정 길이의 창에서 월수익률의 표준편차를 반복 계산하거나 제곱수익률의 시간 패턴을 살핀다. 다만 과거 창에 근거한 추정치는 미래 변동성을 정확히 알려주지 않는다. 특히 창의 길이를 12개월로 할지 120개월로 할지에 따라 급격한 변화에 반응하는 속도가 달라진다.
\par

\paragraph{가격이 거의 움직이지 않는 자산은 정말 안전한가?}

상장주식은 매일 거래되지만, 비상장 부동산은 분기별 감정평가만 받을 수 있다. 평가가격이 늦게 반영되면 관측 수익률의 표준편차가 낮아 보인다. 실제 매도 시점에는 할인과 거래비용이 한꺼번에 나타날 수 있다. 따라서 자산 간 변동성을 비교할 때는 \textbf{가격을 관측하는 빈도와 거래 가능성}을 함께 확인해야 한다.
\par

\subsection{슬라이드 14: 모집단 평균·분산: 상태표를 직접 계산}

상태별 확률이 $(.1,.2,.4,.2,.1)$, 수익률이 $(-.05,.05,.15,.25,.35)$이면 확률의 합은 1이다. 기대수익은 각 수익률에 확률을 곱해 더한다.
\par

\[
\mu=\sum_{s=1}^{5}h(s)R(s)=(-.005)+(.010)+(.060)+(.050)+(.035)=.150.
\]

평균에서의 편차는 $(-.20,-.10,0,.10,.20)$이다. 따라서 $\sigma^2=.1(.20)^2+.2(.10)^2+.4(0)^2+.2(.10)^2+.1(.20)^2=.012$, $\sigma=\sqrt{.012}\simeq10.95\%$다. 이 계산은 실제 상태확률을 알고 있다는 \textbf{모집단} 설정이다.
\par

\paragraph{분산의 다른 계산법으로 검산하기}

분산을 편차제곱으로 계산하는 대신 항을 펼치면 $\operatorname{Var}(R)=E[(R-\mu)^2]=E[R^2]-2\mu E[R]+\mu^2=E[R^2]-\mu^2$다. 이 상태표에서는 $E[R^2]=.0345$이고 $\mu^2=.0225$이므로 분산은 $.012$다. 편차표로 얻은 값과 일치해야 한다. 이 등식은 \textbf{평균을 뺀 뒤 제곱}하는 계산의 교차항을 이해하는 데도 유용하다.
\par

원금 100만원을 투자했다면 기대 기말 자산은 배당 등 다른 현금흐름이 없을 때 115만원이다. 표준편차 10.95\%는 수익률 분포가 평균 주변에 흩어진 정도를 나타낸다. 상태별 수익률 중 실제 최소값은 -5\%이고, 이 예에서는 상태표를 알고 있으므로 손실 확률도 직접 더할 수 있다.
\par

\paragraph{나쁜 상태가 드물어도 기대수익률이 양수일 수 있나?}

90\% 확률로 10\%를 벌고 10\% 확률로 40\%를 잃으면 기대수익률은 $0.9(10\%)+0.1(-40\%)=5\%$다. 평균이 양수여도 열 번 중 한 번꼴 큰 손실이 있을 수 있다. 기대값만으로 투자 여부를 판단하지 않는 이유다. 특히 손실 때 빚을 갚아야 하는 사람에게는 같은 평균이라도 훨씬 부담스럽다.
\par

\paragraph{확률의 합과 분산 단위를 검사하는 법}

상태확률의 합 $\sum_sh(s)=1$인지 먼저 확인한다. 평균은 수익률과 같은 단위이고 분산은 수익률 단위의 제곱이다. 수익률을 15\%처럼 퍼센트 숫자로 입력하면 분산의 숫자는 소수 입력 때보다 $100^2$배 커진다. 표준편차로 제곱근을 취해야 다시 수익률과 같은 단위가 된다.
\par

분산을 $E[R^2]-E[R]^2$로도 계산할 수 있다. 이 예의 $E[R^2]=.0345$이고 $\mu^2=.0225$이므로 $.0345-.0225=.012$다. 편차제곱 방식과 같은 결과가 나와야 표의 숫자가 일관된다.
\par

\subsection{슬라이드 15: 표본평균과 표본분산은 왜 다른가}

같은 다섯 수익률을 각각 한 번씩 관측했다면 각 관측치에 $1/5$의 가중치를 주는 표본평균은 $\bar R=.15$다. 이는 앞 장의 상태확률이 다섯 값에 균등하지 않았다는 점과 다르다. 표본분산은 다음과 같다.
\par

\[
s^2=\frac{1}{T-1}\sum_{t=1}^{T}(R_t-\bar R)^2=\frac{.04+.01+0+.01+.04}{4}=.025.
\]

분모 $T-1$은 표본평균을 이미 같은 자료에서 추정해 자유도 하나를 사용했기 때문에, 독립·동일분포 표본의 모집단 분산에 대해 불편추정량을 만들기 위한 것이다. 여기서는 모집단의 확률가중 분산 $.012$와 표본분산 $.025$를 숫자만 비교하면 안 된다. \textbf{가중 체계와 질문이 다르다.}
\par

\paragraph{자유도 보정을 항등식으로 확인하기}

참 평균 $\mu$와 표본평균 $\bar R$에 대해 $\sum_t(R_t-\bar R)^2=\sum_t(R_t-\mu)^2-T(\bar R-\mu)^2$가 성립한다. 독립·동일분포 표본이라면 첫 항의 기대값은 $T\sigma^2$, 마지막 항의 기대값은 $T\operatorname{Var}(\bar R)=\sigma^2$다. 따라서 편차제곱합의 기대값은 $(T-1)\sigma^2$여서 $T-1$로 나누는 표본분산이 불편추정량이 된다.
\par

표본분산 $.025$는 상태확률을 알고 계산한 분산 $.012$와 다르지만, 앞의 다섯 관측에 상태확률 $.1,.2,.4,.2,.1$을 적용하지 않았으므로 모순이 아니다. \textbf{모형의 확률}과 \textbf{표본의 관측 빈도}를 어떤 질문에 쓰는지 먼저 정한다.
\par

\paragraph{같은 관측치도 분모에 따라 얼마나 달라지나?}

수익률 관측치가 2\%, 4\%, 6\%이면 표본평균은 4\%, 편차제곱합은 $(-.02)^2+0^2+(.02)^2=.0008$이다. 세 값이 \textbf{모집단 전체}라면 분산은 $.0008/3\simeq.000267$이다. 더 큰 모집단에서 뽑은 \textbf{표본}으로 모분산을 추정한다면 $.0008/(3-1)=.0004$를 쓴다. 숫자가 달라지는 이유는 계산 오류가 아니라 어떤 대상을 설명하거나 추정하는지의 차이다.
\par

\paragraph{왜 표본분산은 T-1로 나누나?}

표본평균을 계산하면 편차들의 합 $\sum_t(R_t-\bar R)=0$이 자동으로 성립한다. 따라서 앞의 $T-1$개 편차가 정해지면 마지막 편차는 자유롭게 움직일 수 없다. 모평균이 알려진 상황에서는 분모 $T$를 사용할 수 있지만, 같은 표본으로 모평균을 추정한 뒤 분산까지 추정하면 분모 $T$는 평균적으로 분산을 작게 만든다.
\par

서로 다른 관측 시점이 독립적이고 같은 분포에서 왔다는 가정도 중요하다. 금융수익률의 변동성이 시간에 따라 변하거나 관측치가 종속적이면 $T-1$만 사용한다고 미래 위험 추정 문제가 모두 해결되지는 않는다.
\par

\subsection{슬라이드 16: 산술평균과 기하평균}

여러 기간의 부는 수익률을 더하지 않고 곱한다. $W_T=W_0\prod_{t=1}^{T}(1+R_t)$이므로 매 기간 같은 수익률 $g$로 환산하면 $1+g=[\prod_t(1+R_t)]^{1/T}$다. $(10\%,25\%,-20\%,25\%)$의 산술평균은 10\%이고 기하평균은 $[1.10\times1.25\times.80\times1.25]^{1/4}-1\simeq8.29\%$다.
\par

산술평균은 \textbf{한 기간의 기대수익}을 추정할 때, 기하평균은 \textbf{실현된 장기 복리 성장률}을 요약할 때 적합하다. 산술-기하평균 부등식 때문에 같은 표본에서 산술평균은 기하평균 이상이다. 변동성이 커질수록 복리 누적은 더 크게 깎일 수 있다. 예를 들어 +50\% 뒤 -50\%면 산술평균은 0이지만 원금은 75\%가 된다.
\par

\paragraph{변동성이 복리 성장률을 낮추는 이유}

복리 성장의 관점에서는 $\log(1+g)=T^{-1}\sum_t\log(1+R_t)$다. 로그함수는 아래로 굽어 있으므로 Jensen 부등식에 의해 $E[\log(1+R)]\le\log(1+E[R])$다. 수익률의 흔들림이 장기 복리 성장률을 낮추는 수학적 이유다.
\par

수익률 변동이 작을 때 $\log(1+R)\simeq R-R^2/2$이므로 평균 로그성장률은 대략 $\mu-(\sigma^2+\mu^2)/2$다. 자주 쓰는 $g\simeq\mu-\sigma^2/2$는 추가 근사까지 포함한다. 큰 손실이나 레버리지가 있으면 원래의 수익률 곱셈식으로 계산한다.
\par

\paragraph{산술평균이 5\%면 매년 5\%씩 자란다는 뜻인가?}

첫해 +20\%, 다음 해 -10\%의 산술평균은 5\%다. 100만원의 최종금액은 $100\times1.2\times0.9=108$만원이고 연평균 복리성장률은 $\sqrt{1.08}-1\simeq3.92\%$다. 두 해의 결과를 한 해 수익률로 환산할 때는 실제 최종금액에 제곱근을 적용한다.
\par

\paragraph{변동성이 장기 성장률을 낮추는 이유}

로그를 취하면 기간당 실현 복리 성장률은 $\frac1T\sum_t\log(1+R_t)$이다. 로그함수는 오목하므로 젠슨 부등식에서 $E[\log(1+R)]\le\log(1+E[R])$이다. 평균 수익률을 고정해도 수익률의 퍼짐이 크면 평균 로그성장률이 낮아지는 이유다.
\par

작은 수익률에 대해 테일러 전개를 적용하면 $E[\log(1+R)]\approx E[R]-\frac12E[R^2]$이다. 자주 쓰이는 '기하평균 \ensuremath{\approx} 산술평균 - 분산의 절반'은 평균과 고차항에 관한 추가 근사를 포함하므로, 변동이 큰 레버리지 상품에는 정확한 항등식처럼 적용하지 않아야 한다.
\par

\subsection{슬라이드 17: 공분산: 함께 움직이는 정도}

두 자산의 공분산은 평균에서 벗어난 정도의 곱을 확률로 가중한다.
\par

\[
\operatorname{Cov}(R_i,R_j)=E[(R_i-\mu_i)(R_j-\mu_j)].
\]

둘 다 평균 위나 아래인 상태의 곱은 양수, 한쪽만 위인 상태는 음수다. 예시에서 두 평균이 각각 5\%, 표준편차가 각각 10\%이고 공분산이 $20$ 퍼센트포인트 제곱이라면 소수 단위에서는 $0.002$다. 퍼센트 숫자 20을 소수 공분산 20으로 쓰면 단위 오류가 생긴다. 공분산은 방향을 알려주지만 자산의 변동성 크기에 따라 수치가 달라져 직접 비교가 어렵다.
\par

\paragraph{공분산의 부호를 상태별로 해석하기}

경기 호황과 불황에서 자산 A가 평균보다 각각 $+10,-10$\%포인트, 자산 B가 각각 $+5,-5$\%포인트 움직이면 두 상태의 편차곱은 모두 양수다. 두 자산은 같은 방향으로 움직이므로 공분산이 양수다. 반대로 B의 편차가 $-5,+5$\%포인트라면 편차곱이 모두 음수이고 한 자산의 손실을 다른 자산이 일부 메운다.
\par

공분산은 수익률 단위의 \textbf{제곱}을 갖는다. 계산에 10과 5를 넣었다면 결과는 퍼센트포인트 제곱이고, $.10$과 $.05$를 넣었다면 소수 수익률의 제곱이다. 자산마다 변동성이 다른 경우 비교하기 어렵기 때문에 다음 슬라이드에서 표준편차로 나눈 상관계수를 쓴다.
\par

\paragraph{공분산이 양수면 두 자산을 함께 보유하면 안 되나?}

양의 공분산은 함께 오르내리는 경향을 뜻하지만, 두 자산의 상관이 1보다 작으면 한 자산만 보유할 때보다 분산이 줄 수 있다. 예컨대 표준편차가 각각 10\%이고 상관이 0.2인 두 자산을 절반씩 담으면 포트폴리오 표준편차는 $\sqrt{0.5^2(0.1)^2+0.5^2(0.1)^2+2(0.5)^2(0.2)(0.1)^2}\simeq7.75\%$다. 분산 효과에는 음의 상관이 필수는 아니다.
\par

\paragraph{삼성전자와 코스피의 공동 움직임}

삼성전자 주식과 코스피 지수가 같은 날 함께 오르거나 내리는 경향은 \textbf{시장과의 공변동}을 떠올리게 한다. 그러나 삼성전자는 코스피 구성종목이므로 두 수익률의 상관에는 지수 구성상 기계적인 부분도 있다. 코스피를 '시장 전체'의 완전한 대리변수로 보아서는 안 되고, 다른 국내·해외 자산 및 투자자의 기존 소득과의 관계도 중요하다.
\par

공동 움직임을 한 숫자로 요약하려면 $R_{i,t}=\alpha_i+\beta_iR_{M,t}+\varepsilon_{i,t}$라는 선형 투영을 생각할 수 있다. 평균제곱오차 $E[(R_i-\alpha-\beta R_M)^2]$를 $\alpha,\beta$에 대해 최소화하면 첫 조건에서 $\alpha=\mu_i-\beta\mu_M$가 된다. 이를 두 번째 조건에 넣으면 $\operatorname{Cov}(R_i,R_M)-\beta\operatorname{Var}(R_M)=0$이므로 다음을 얻는다.
\par

\[
\beta_i=\frac{\operatorname{Cov}(R_i,R_M)}{\operatorname{Var}(R_M)}.
\]

베타는 시장 수익률이 1\%포인트 움직일 때 개별주식의 \textbf{선형 예측 수익률}이 평균적으로 몇 \%포인트 움직이는지를 뜻한다. 이것은 회귀 기울기이지 시장 변화가 해당 기업의 수익률을 인과적으로 변화시킨다는 증명은 아니다. 회사 고유의 뉴스가 남는 잔차 $\varepsilon_i$와 시장에 함께 노출되는 부분을 구분하면 분산투자 논리와 연결된다.
\par

\subsection{슬라이드 18: 상관계수와 회귀의 R²}

상관계수는 공분산을 각 표준편차로 나눈 무차원 값이다. $\rho_{ij}=\operatorname{Cov}(R_i,R_j)/(\sigma_i\sigma_j)$이고 코시-슈바르츠 부등식에 의해 $-1\le\rho\le1$이다. 앞 예시에서 $20/(10\cdot10)=.20$이다. 단위가 퍼센트이든 소수이든 일관되게 계산하면 결과가 같다.
\par

상수항을 포함한 \textbf{단순 선형회귀}에서는 표본 $R^2$가 표본 상관계수의 제곱과 같다. $\rho=.2$이면 관측된 선형 관계가 종속변수 변동의 4\%를 설명한다. 인과관계는 이 통계량만으로 판단할 수 없다. 표본 공분산의 분모는 보통 $T-1$이고, 회귀 잔차분산에서는 절편과 기울기를 추정한 자유도를 반영해 $T-2$를 쓴다.
\par

\paragraph{회귀 설명력의 수학적 연결}

절편이 있는 단순회귀에서 기울기는 $b=\operatorname{Cov}(X,Y)/\operatorname{Var}(X)$다. 예측값의 분산은 $b^2\operatorname{Var}(X)=\operatorname{Cov}(X,Y)^2/\operatorname{Var}(X)$이므로 이를 $\operatorname{Var}(Y)$로 나누면 $R^2=\operatorname{Cov}(X,Y)^2/[\operatorname{Var}(X)\operatorname{Var}(Y)]=\rho^2$가 된다. $\rho=.2$일 때 4\%라는 숫자가 나오는 과정을 한 단계씩 확인할 수 있다.
\par

이 4\%는 \textbf{해당 표본의 선형 회귀 설명력}이다. $X$가 $Y$를 원인으로 움직이게 한다는 뜻도, 미래 예측 정확도가 반드시 4\%라는 뜻도 아니다. 기간을 바꾸거나 추가 변수를 넣으면 설명력이 달라질 수 있다.
\par

\paragraph{음의 상관도 회귀 설명력은 높을 수 있나?}

절편이 있는 단순회귀에서 표본 상관계수가 $\rho=-.6$이면 회귀선의 기울기는 음수지만 $R^2=(-.6)^2=.36$이다. 상관의 \textbf{부호는 관계의 방향}, $R^2$는 표본의 선형 변동 중 설명한 비율을 나타낸다. $\rho=.6$인 경우도 $R^2=.36$이다. 두 경우의 예측 방향은 정반대이므로 설명력 숫자만 읽으면 중요한 정보를 놓친다.
\par

\paragraph{상관계수 제곱의 해석이 제한되는 이유}

단순 최소제곱 회귀 $R_j=a+bR_i+\varepsilon$에 절편 $a$를 포함하고 같은 관측 표본을 사용하면 $R^2=\rho_{ij}^2$가 된다. 절편을 제거하거나 통계 모형을 바꾸면 일반적으로 이 등식이 성립하지 않는다. $R^2$는 그 표본에서 선형 예측이 설명한 변동 비율이며 인과관계를 말하지 않는다.
\par

상관계수가 0이어도 정보가 없는 것은 아니다. 예를 들어 불확실성이 클 때 두 자산이 극단적으로 함께 움직이는 비선형 의존이 있으면 평상시 선형 상관계수 하나는 손실 위험을 충분히 보여주지 못한다.
\par

\subsection{슬라이드 19: 상관관계를 그림으로 읽기}

산점도가 우상향이면 양의 상관, 우하향이면 음의 상관이다. $\rho=0$이면 선형 공변동이 없지만 곡선형 관계는 남을 수 있다. 예를 들어 평균이 0인 대칭 변수 $X$와 $X^2$는 일반적으로 의존하면서 공분산이 0일 수 있다.
\par

상관이 +1이면 한 자산의 변동이 다른 자산의 변동에 정확히 비례해, 단순한 결합으로 위험을 줄일 여지가 제한된다. -1이면 적절한 양의 비중만으로 변동을 완전히 상쇄할 수 있는 경우가 있다. 실제 시장의 상관은 시간에 따라 변하고 위기 때 바뀌기도 한다.
\par

\paragraph{상관 0의 구체적 반례}

$X$가 $-1,0,1$을 같은 확률로 갖고 $Y=X^2$라고 하자. $Y$는 $X$를 알면 완전히 결정되므로 둘은 독립이 아니다. 그러나 $E[X]=0$, $E[XY]=E[X^3]=0$이므로 $\operatorname{Cov}(X,Y)=0$이다. 상관 0은 \textbf{직선적인 공동 움직임이 없다}는 뜻이지 관계가 전혀 없다는 뜻이 아님을 보여준다.
\par

금융시장에서도 평상시 수익률의 상관은 낮지만 두 자산이 큰 손실 상태에서 함께 하락할 수 있다. 평균 주변의 상관계수와 꼬리에서의 공동 손실은 다른 질문이다. 분산투자 효과를 평가할 때 산점도의 극단 영역을 함께 보는 이유다.
\par

\paragraph{상관 0인 두 자산을 반씩 보유하면 위험은 얼마인가?}

두 자산의 표준편차가 각각 10\%와 20\%, 상관이 0이면 절반씩 보유한 분산은 $.5^2(.10)^2+.5^2(.20)^2=.0125$이고 표준편차는 약 11.18\%다. 교차 공분산 항이 0이어도 개별 분산 항은 남는다. 상관 0의 정확한 뜻과 자산을 결합했을 때의 위험을 구분할 수 있다.
\par

\subsection{슬라이드 20: 역사적 상관행렬의 활용}

평균·표준편차·상관행렬은 포트폴리오 위험 계산의 입력값이다. 대각선의 상관은 1이고 행렬은 대칭이다. 공분산은 $\sigma_{ij}=\rho_{ij}\sigma_i\sigma_j$로 복원할 수 있다. 상관이 낮은 자산을 결합하면 분산투자 이익을 기대할 수 있지만, 낮은 역사적 상관이 미래에도 유지된다는 보장은 없다.
\par

역사적 연간 수익률 표에서 S\&P 지수의 산술평균은 약 12.96\%, 기하평균은 약 11.00\%, 표준편차는 약 20.32\%다. 소형주는 각각 약 17.73\%, 12.71\%, 33.89\%로 평균이 더 높지만 변동도 훨씬 크다. 두 주식군의 상관은 약 0.81인 반면 S\&P와 장기 국채의 상관은 약 0.19, 장기 회사채와 장기 국채의 상관은 약 0.94로 제시된다. 따라서 '채권'이라는 이름만으로 분산효과가 같지 않으며, 같은 자산군 안에서도 공동 움직임이 다르다.
\par

상관행렬은 모든 포트폴리오의 분산을 0 이상으로 만드는 \textbf{양의 준정부호} 행렬이어야 한다. 추정치가 불안정하면 극단적인 최적 비중이 나올 수 있어 축소추정과 비중 제약을 적용하기도 한다.
\par

\paragraph{상관행렬이 만족해야 하는 조건}

상관행렬의 대각선은 각 자산과 자기 자신의 상관이므로 1이고, $\rho_{ij}=\rho_{ji}$로 대칭이다. 더 중요한 조건은 아무 비중 벡터 $\boldsymbol w$에 대해서도 $\boldsymbol w^\top\Sigma\boldsymbol w\ge0$이어야 한다는 것이다. 좌변이 바로 포트폴리오 분산이므로 음수가 나오면 그 숫자들의 조합은 유효한 공분산행렬이 아니다.
\par

두 자산에서는 $|\rho|\le1$이면 이 조건을 확인하기 쉽지만, 여러 자산의 상관을 서로 독립적으로 추측하면 각 쌍은 그 범위 안에 있어도 전체 행렬이 불가능할 수 있다. 과거 표본에서 추정한 행렬도 미래에는 변하므로 최적화 결과가 지나치게 극단적이면 추정 오차와 비중 제약을 검토한다.
\par

\paragraph{각 상관계수가 -1과 1 사이면 상관행렬은 언제나 유효한가?}

세 자산 A·B, A·C의 상관이 각각 0.9인데 B·C의 상관이 -0.9라는 표를 생각해 보자. A와 강하게 함께 움직이는 B와 C가 서로 강하게 반대로 움직인다는 조건은 동시에 성립할 수 없다. 유효한 상관행렬은 모든 비중 벡터에 대해 계산한 분산이 음수가 아니어야 한다. 추정한 행렬이 이 조건을 어기면 포트폴리오 계산 전에 수정해야 한다.
\par

\subsection{슬라이드 21: 포트폴리오 분산식의 유도}

상수 $a,b$에 대해 $E[aR_i+bR_j]=a\mu_i+b\mu_j$다. 평균에서의 편차를 $X=R_i-\mu_i$, $Y=R_j-\mu_j$라 두면 다음처럼 제곱을 전개한다.
\par

\[
\begin{aligned}
\operatorname{Var}(aR_i+bR_j)&=E[(aX+bY)^2]\\
&=a^2E[X^2]+b^2E[Y^2]+2abE[XY]\\
&=a^2\sigma_i^2+b^2\sigma_j^2+2ab\sigma_{ij}.
\end{aligned}
\]

마지막 교차항이 두 자산을 함께 보유할 때 생기는 효과다. 공분산이 작을수록 같은 개별 위험과 비중에서 전체 분산이 작다. $\operatorname{Cov}(aR_i+bR_j,R_k)=a\operatorname{Cov}(R_i,R_k)+b\operatorname{Cov}(R_j,R_k)$도 기대값의 선형성에서 곧바로 나온다.
\par

\paragraph{교차항이 두 번 나오는 이유}

두 자산의 평균 편차를 $X,Y$로 쓰면 $(aX+bY)^2=a^2X^2+abXY+abYX+b^2Y^2$다. 숫자의 곱에서는 $XY=YX$이므로 가운데 항이 $2abXY$가 된다. 기대값을 취하면 $E[XY]$가 공분산이고, 포트폴리오 분산식의 \textbf{2배 교차항}이 나온다. 이 계수를 빠뜨리면 분산투자 효과를 절반만 반영하게 된다.
\par

비중이 각각 절반이고 두 자산의 분산이 같은 $\sigma^2$라면 $\sigma_p^2=.5\sigma^2+.5\operatorname{Cov}(R_1,R_2)$다. 완전히 같은 방향이면 공분산이 $\sigma^2$여서 전체 분산도 $\sigma^2$다. 독립이면 공분산 0이어서 분산은 절반으로 줄지만 표준편차는 절반이 아닌 $1/\sqrt2$배다.
\par

\paragraph{두 자산을 반씩 담으면 분산은 어떻게 계산하나?}

각 자산의 표준편차가 10\%, 상관이 0.2이면 각 분산은 0.01, 공분산은 $0.2(0.1)(0.1)=0.002$다. 절반씩 담은 분산은 $0.25(0.01)+0.25(0.01)+2(0.5)(0.5)(0.002)=0.006$이고, 표준편차는 $\sqrt{0.006}\simeq7.75\%$다. 마지막 항을 빠뜨리면 두 자산이 함께 움직이는 효과를 무시하게 된다.
\par

\subsection{슬라이드 22: 두 자산 배분으로 넘어가기}

수익률 측정과 공분산을 익혔으면 선택 문제를 풀 수 있다. 두 자산 비중을 $w$와 $1-w$로 두고 기대수익과 분산을 $w$의 함수로 나타낸다. 그러면 가능한 모든 조합을 그리는 투자기회집합, 위험을 가장 작게 하는 최소분산 포트폴리오, 무위험자산이 있을 때 가장 가파른 자본배분선을 차례로 구할 수 있다.
\par

앞의 통계식은 \textbf{입력값을 어떻게 측정할지}, 다음의 최적화식은 \textbf{그 입력값으로 무엇을 선택할지} 답한다. 추정 오차가 선택에 영향을 준다는 점도 함께 기억한다.
\par

\paragraph{통계 입력에서 선택 문제로}

두 자산의 평균, 표준편차, 상관을 구했다면 비중 $w$에 대해 $\mu_p=w\mu_1+(1-w)\mu_2$와 $\sigma_p^2=w^2\sigma_1^2+(1-w)^2\sigma_2^2+2w(1-w)\rho\sigma_1\sigma_2$를 동시에 계산한다. \textbf{평균만 높은 자산을 전부 사는 선택}과 \textbf{위험만 가장 낮추는 선택}은 일반적으로 다르다. 목표 수익률이나 위험 허용 범위가 있어야 투자자의 문제를 완성할 수 있다.
\par

실제 분석의 평균과 공분산은 추정치다. 추정 평균의 작은 변화가 최적 비중을 크게 바꾸면 선택 문제가 입력값에 민감하다는 뜻이다. 다음 장에서는 먼저 입력값을 안다고 가정한 기준 해를 구하고, 그 해를 현실에 적용할 때 제약과 추정 오차가 어디에 들어가는지 따진다.
\par

\paragraph{최적 비중을 소수점까지 계산하면 그대로 실행해야 하나?}

두 자산의 추정 기대수익률이 10.0\%와 10.2\%라면 0.2\%포인트의 작은 차이에 최적화 결과가 민감하게 반응할 수 있다. 실제 기대수익률을 그보다 정확히 알기 어렵다면 계산된 47.36\% 같은 비중에 과도한 의미를 부여하면 안 된다. 먼저 가정이 조금 바뀔 때 비중과 위험이 얼마나 바뀌는지 \textbf{민감도}를 확인하고, 거래비용·세금·보유 제약도 반영한다.
\par

\clearpage

\section{3부 · 두 자산 포트폴리오}

\subsection{슬라이드 1: 두 자산으로 만드는 투자기회집합}

두 자산을 섞으면 평균은 비중의 가중합으로, 위험은 분산과 공분산의 결합으로 결정된다. 가능한 비중마다 $(\sigma_p,\mu_p)$를 계산해 그린 점들의 집합이 \textbf{투자기회집합}이다. 이 곡선에서 같은 위험에 더 높은 기대수익을 주는 점을 찾는다.
\par

먼저 무위험자산 하나와 위험자산 하나를 결합한다. 이후 두 위험자산의 곡선, 마지막으로 그 곡선과 무위험자산의 접선을 구한다. 세 단계는 각각 직선, 곡선, 최적 직선이라는 기하적 연결을 이룬다.
\par

\paragraph{투자기회집합의 각 점이 뜻하는 것}

평균-표준편차 평면의 각 점은 \textbf{정해진 비중의 포트폴리오}다. 두 자산의 비중을 바꿀 때 각각의 미래 수익률이 어떻게 합쳐지는지 계산해 점을 얻는다. 가로축이 오른쪽일수록 수익률의 흔들림이 크고, 세로축이 위쪽일수록 기대수익이 높다. 투자자는 보통 같은 위험이면 높은 평균을, 같은 평균이면 낮은 위험을 선호한다.
\par

이 그림만으로 특정 투자자의 최종 선택이 결정되지는 않는다. 같은 효율적 집합 위에서도 손실을 싫어하는 정도와 투자기간이 다르면 고르는 점이 다르다. 먼저 \textbf{가능한 조합의 집합}을 구하고, 그다음 선호와 제약을 이용해 하나를 고른다는 두 단계를 구분한다.
\par

\paragraph{효율적 포트폴리오는 하나뿐인가?}

효율적 프런티어에는 위험과 기대수익률이 서로 다른 여러 포트폴리오가 놓인다. 변동성 10\%에 기대수익률 8\%인 A와 변동성 15\%에 기대수익률 11\%인 B 중 어느 쪽이 더 나은지는 투자자가 추가 3\%포인트의 기대수익을 위해 추가 5\%포인트의 위험을 감수할지에 달려 있다. \textbf{효율적}은 같은 위험에서 더 높은 평균을 가진다는 비교 기준이지, 모두에게 하나의 정답을 뜻하지 않는다.
\par

\subsection{슬라이드 2: 자산배분 문제의 변수}

총 자산 중 위험자산 비중을 $w$, 무위험자산 비중을 $1-w$로 둔다. $0<w<1$은 분산 보유, $w>1$은 무위험 금리로 차입해 위험자산을 더 산 상태다. $w<0$은 위험자산을 공매도하고 무위험자산을 더 가진 상태로 해석할 수 있다.
\par

투자자가 선택할 수 있는 범위는 현실의 차입금리, 증거금, 공매도 가능 여부에 달려 있다. 이 장의 첫 식은 차입과 대출에 같은 무위험 수익률이 적용되고 세금·비용이 없다는 기준 모형이다.
\par

\paragraph{비중이 자본과 차입을 함께 나타내는 방식}

자기자본 100 중 위험자산에 120을 투자하면 $w=1.2$이고 무위험자산 비중은 $1-w=-.2$다. 음의 무위험 비중은 20을 빌려서 총 120을 샀다는 회계적 표현이다. 반대로 $w=-.2$라면 위험자산 20을 공매도하고 그 대금까지 합쳐 무위험자산 120을 가진 것으로 해석한다. 언제나 비중의 합은 1이지만 각 비중이 0과 1 사이일 필요는 없다.
\par

그러나 이 비중만으로 실제 거래가 가능한지는 알 수 없다. 공매도 담보와 대차료, 차입금리, 레버리지 한도에 따라 가능한 $w$의 범위가 달라진다. 수학적으로 계산한 직선 전체와 투자자가 실행할 수 있는 구간을 구별해야 한다.
\par

\paragraph{음의 비중은 손익에 어떻게 들어가나?}

자기자본 100만원으로 A를 130만원 매수하고 B를 30만원 공매도하면 비중은 $(1.3,-.3)$이다. 이후 A가 10\% 하락하고 B가 20\% 상승하면 계좌 수익률은 $1.3(-.10)-.3(.20)=-.19$, 곧 -19\%다. B 가격 상승은 공매도자에게 손실이므로 음의 비중과 양의 수익률을 곱한 항이 음수가 된다. 거래비용과 대차료를 빼기 전 계산이다.
\par

\subsection{슬라이드 3: 계산을 위한 핵심 질문}

무위험자산을 추가하면 어떤 위험-수익 조합이 가능한가? 두 위험자산의 상관이 낮아지면 분산이 얼마나 감소하는가? 최소분산 조합은 어떻게 찾는가? 무위험자산까지 허용하면 위험자산끼리 어떤 비중으로 묶는 것이 좋은가?
\par

이 질문들의 답은 서로 다른 최적화다. \textbf{최소분산}은 수익률을 보지 않고 위험만 최소화한다. \textbf{접점 포트폴리오}는 무위험 금리보다 높은 기대수익을 위험 한 단위당 가장 많이 제공하는 위험자산 조합이다.
\par

\paragraph{세 가지 최적화 문제를 분리하기}

최소분산 포트폴리오는 $\min_w\sigma_p^2(w)$를 풀며 평균을 목적함수에 넣지 않는다. 목표 평균 12\%를 정한 문제는 $\min_w\sigma_p^2(w)$에 $\mu_p(w)=.12$라는 제약을 추가한다. 무위험자산이 있을 때의 접점 문제는 위험자산 조합 $P$의 $(\mu_P-r_f)/\sigma_P$를 가장 크게 만든다. 비슷한 그래프를 사용하지만 \textbf{최적화 대상과 제약이 각각 다르다}.
\par

각 문제를 풀기 전에는 선택 변수가 무엇인지 적는다. 두 위험자산 사이의 상대 비중과, 만들어진 위험자산 묶음에 전체 재산을 얼마나 배분할지는 다른 변수다. 위험자산 내부 비중을 먼저 구하고 전체 위험자산 비중을 곱하면 실제 계좌 비중이 나온다.
\par

\paragraph{최소분산 포트폴리오가 기대수익도 최고인가?}

최소분산 포트폴리오는 가능한 조합 중 \textbf{분산이 가장 작은 점}이다. 무위험자산과 결합할 때 샤프비율이 가장 높은 접점 포트폴리오와 목적함수가 다르다. 예컨대 최소분산점의 기대수익률이 9\%이고 다른 조합이 12\%라면, 후자는 더 위험할 수 있어도 위험 한 단위당 초과수익이 더 클 수 있다. 최소분산과 최대 샤프비율을 같은 뜻으로 읽지 않는다.
\par

\subsection{슬라이드 4: 무위험자산의 의미}

\textbf{무위험자산}은 정해진 투자기간의 명목 수익률이 확정된 자산이다. 투자기간과 만기가 일치하는 단기 국채를 근사치로 쓸 수 있다. 만기 전 매도에는 금리 위험이, 외화 국채에는 환율 위험이 더해진다. 물가가 변하면 명목 지급액이 같아도 구매력은 달라진다.
\par

한편 위험자산은 개별 주식일 수도 여러 주식의 포트폴리오일 수도 있다. 여기서 $\sigma_f=0$은 \textbf{선택한 투자기간의 명목 수익률}에 관한 가정이다.
\par

\paragraph{무위험이라는 말의 조건}

1년 뒤 원화 100만원을 확정적으로 지급하는 자산은 1년 뒤 원화 지출을 준비하는 투자자에게 명목상 무위험에 가깝다. 그러나 같은 자산을 3개월 뒤 팔아야 하면 그때의 시장금리에 따라 매도가격이 달라진다. 1년 뒤에도 물가가 크게 올라 구매력이 변할 수 있다. 무위험의 기준은 \textbf{보유기간과 평가 통화, 명목·실질 목표}를 함께 지정해야 정해진다.
\par

수익률이 확정인 자산과 가격이 안정적으로 보이는 자산도 다르다. 만기일의 지급액이 정해진 무이표채는 중간 시가가 변하고, 가격이 자주 평가되지 않는 비유동 자산은 관측 변동이 작아 보여도 위험이 없는 것은 아니다.
\par

\paragraph{국채라면 언제나 무위험자산인가?}

만기가 1년인 국채를 3개월 뒤 매도할 투자자에게는 매도 당시 금리가 가격을 바꾸므로 3개월 수익률이 확정되지 않는다. 반면 3개월 만기 무이표국채를 만기까지 보유할 때는 명목 지급액이 확정된다. 그래도 물가가 오르면 \textbf{실질 구매력}은 변한다. 무위험은 통화, 보유기간, 명목·실질 기준을 정한 뒤에야 의미가 분명해진다.
\par

\subsection{슬라이드 5: 무위험·위험자산 혼합의 평균과 분산}

위험자산의 평균과 표준편차가 $\mu_P=15\%,\sigma_P=22\%$, 무위험 수익률이 $r_f=7\%$라고 하자. 포트폴리오 수익률은 $R_c=(1-w)r_f+wR_P$다. 기대값의 선형성과 상수의 분산이 0이라는 사실을 쓰면 다음과 같다.
\par

\[
\mu_c=(1-w)r_f+w\mu_P=.07+.08w,\qquad
\sigma_c^2=w^2\sigma_P^2,\quad \sigma_c=|w|\sigma_P.
\]

$w=.75$이면 평균 13\%, 표준편차 16.5\%다. $w$가 음수인 공매도 구간에서도 분산은 $w^2\sigma_P^2$, 표준편차는 $|w|\sigma_P$로 계산된다.
\par

\paragraph{평균과 분산을 따로 전개하기}

$R_c=(1-w)r_f+wR_P$에서 기대값을 취하면 $E[R_c]=(1-w)r_f+wE[R_P]$다. 평균으로부터의 편차를 계산하면 $R_c-E[R_c]=w(R_P-\mu_P)$이므로 제곱의 기대값은 $w^2\sigma_P^2$다. 무위험 수익률은 모든 상태에서 같아 편차가 0이라는 점이 핵심이다. 마지막으로 음수가 아닌 제곱근을 취해 $\sigma_c=|w|\sigma_P$를 얻는다.
\par

예시의 $w=.75$는 자본 100 중 75를 위험자산에, 25를 무위험자산에 둔다. 기대수익은 $25(.07)+75(.15)=13$의 자본 100당 기대이익으로 확인할 수 있다. 위험자산의 손익 변화는 전체 포트폴리오에 75\%만 전달되므로 표준편차는 $22\%\times.75=16.5\%$다.
\par

\paragraph{위험자산을 공매도하면 손익 방향은?}

무위험수익률이 4\%이고 위험자산 비중 $y=-.5$라면 무위험자산 비중은 1.5다. 위험자산의 실현수익률이 -20\%인 해에는 계좌 수익률이 $1.5(.04)-.5(-.20)=16\%$다. 반대로 위험자산이 20\% 오르면 $1.5(.04)-.5(.20)=-4\%$다. 공매도는 \textbf{위험자산 수익률에 대한 노출 방향}을 뒤집지만 표준편차는 여전히 양수다.
\par

\paragraph{음의 비중에서 왜 절댓값을 쓰나?}

$R_c=(1-w)r_f+wR_P$에서 상수 $r_f$는 분산에 기여하지 않으므로 $\operatorname{Var}(R_c)=w^2\sigma_P^2$다. 표준편차는 음수가 아닌 제곱근으로 정의되기 때문에 $\sqrt{w^2\sigma_P^2}=|w|\sigma_P$다. $w=-.5$는 위험자산을 50\% 공매도한 상태인데, 수익률의 방향은 뒤집혀도 변동성 크기는 $.5\sigma_P$다.
\par

위험자산의 기대 초과수익이 양수라면 $w<0$의 조합은 무위험자산과 위험자산의 양의 비중 조합에 의해 평균-분산 관점에서 지배될 수 있다. 다만 기존 부채나 다른 포지션을 헤지하는 목적이라면 공매도 비중의 경제적 의미가 달라진다.
\par

\subsection{슬라이드 6: 자본배분선의 기울기}

$w\ge0$에서 $w=\sigma_c/\sigma_P$이므로 앞 식에서 $w$를 제거하면 $\mu_c=r_f+[(\mu_P-r_f)/\sigma_P]\sigma_c$다. 가로축 $\sigma_c$, 세로축 $\mu_c$인 평면에서는 무위험점 $(0,r_f)$과 위험자산점 $(\sigma_P,\mu_P)$을 잇는 직선이다.
\par

기울기는 $(15-7)/22=8/22\simeq.3636$이다. 이는 위험 1 단위를 추가할 때 얻는 기대 초과수익의 비율이다. 직선 위의 $w=.5$는 평균 11\%, 표준편차 11\%이고, $w=0$은 무위험자산 100\%다.
\par

\paragraph{위험자산 공매도는 반대쪽 반직선을 만든다}

본문의 자본배분선은 위험자산 비중 $w\ge0$을 전제로 $\sigma_c=w\sigma_P$를 쓴다. $w<0$이면 표준편차는 $\sigma_c=-w\sigma_P$이고 평균은 $\mu_c=r_f-w_{\!\mathrm{abs}}(\mu_P-r_f)$이다. 여기서 $w_{\!\mathrm{abs}}=|w|$다. 따라서 무위험점의 왼쪽으로 가는 것이 아니라, \textbf{같은 위험축의 오른쪽 아래}로 향하는 별도 반직선이 생긴다.
\par

$\mu_P>r_f$라면 이 아래쪽 반직선은 같은 표준편차에서 무위험자산보다 평균이 낮아 평균·분산 선호 투자자에게 지배된다. 공매도 비중의 부호를 표준편차에도 그대로 붙이면 음의 위험이라는 잘못된 그래프가 나온다.
\par

\paragraph{같은 위험에서 어느 직선이 위에 있나?}

무위험수익률이 3\%일 때 위험포트폴리오 A의 평균·표준편차가 $(9\%,10\%)$, B가 $(12\%,20\%)$라면 기울기는 각각 $.6$과 $.45$다. 계좌의 목표 표준편차를 5\%로 맞추면 A 비중은 $.5$여서 평균 6\%, B 비중은 $.25$여서 평균 5.25\%다. \textbf{같은 위험}에서 A의 직선이 위에 있다. B의 단독 기대수익률이 더 높다는 사실만으로는 순위가 정해지지 않는다.
\par

\subsection{슬라이드 7: 차입을 통한 레버리지}

$w=1.5$이면 자기자본 100을 넣고 50을 빌려 위험자산 150을 보유한다. 모형상 평균은 $7+8(1.5)=19\%$, 표준편차는 $22(1.5)=33\%$다. 가격이 20\% 떨어지면 투자자 자기자본의 대략 30\%가 줄어들며 차입이자와 추가 증거금 부담은 별도로 남는다.
\par

현실에서 차입금리는 무위험 대출금리보다 높고, 담보가치가 하락하면 강제 청산될 수 있다. 따라서 위험자산점 위의 직선 연장은 이론적 기회집합이고 실제 거래 가능선은 더 낮거나 중간에 끊길 수 있다.
\par

\paragraph{차입금리가 높아질 때 자본배분선}

예금으로 운용할 때 7\%를 받고 차입할 때 10\%를 내는 경우를 보자. $w=1.5$, 위험자산 기대수익 15\%라면 계좌의 기대수익은 $1.5(.15)-.5(.10)=17.5\%$다. 차입도 7\%에 가능하다고 가정한 직선의 19\%보다 1.5\%포인트 낮다. \textbf{차입·대출 금리의 차이}가 자본배분선의 차입 구간을 꺾는다.
\par

레버리지의 실행 가능성에는 담보도 들어간다. 가격 하락으로 추가 증거금을 내야 하면 만기 손익이 회복되더라도 중간에 강제 매도될 수 있다. 투자자의 유동성 제약은 기대수익·표준편차 두 숫자만으로 표현되지 않는다.
\par

\paragraph{레버리지는 평균만 비례해서 높이나?}

자기자본 100만원에 50만원을 연 7\%로 빌려 주식 150만원을 샀다고 하자. 주식이 20\% 하락하면 주식 손실은 30만원, 차입 이자는 3.5만원이므로 자기자본 손익은 -33.5만원, 수익률은 -33.5\%다. 주식 하락률보다 손실률이 크다. 담보비율 조건이 있으면 회복을 기다리기 전에 \textbf{강제 매도}가 발생할 수도 있다.
\par

\subsection{슬라이드 8: 샤프비율은 무엇을 비교하나}

위험자산 $P$의 샤프비율은 $S_P=(\mu_P-r_f)/\sigma_P$다. 예시의 $S_P=(.15-.07)/.22\simeq.364$다. 분자는 무위험자산 대신 위험을 감수한 대가, 분모는 수익률 변동성이므로 \textbf{위험 한 단위당 기대 초과수익}으로 읽는다.
\par

같은 기간·통화의 수익률로 계산해야 하며, 분모에 총변동성을 쓰므로 분산이 충분한 시장포트폴리오가 아닌 자산에는 해석상 주의가 필요하다. 역사적 표본 샤프비율은 실제 미래 샤프비율의 추정치일 뿐이다. $\sigma=0$인 무위험자산 자체의 샤프비율은 분모가 0이라 정의되지 않는다.
\par

\paragraph{같은 직선 위의 샤프비율}

위험자산 $P$와 무위험자산을 섞어 $w>0$을 택하면 초과기대수익은 $w(\mu_P-r_f)$이고 표준편차는 $w\sigma_P$다. 따라서 조합의 샤프비율은 $w(\mu_P-r_f)/(w\sigma_P)=(\mu_P-r_f)/\sigma_P$로 동일하다. 무위험자산과의 혼합만으로 \textbf{위험 한 단위당 기대 보상}이 개선되지는 않고, 같은 효율의 위험 규모를 조절한다.
\par

두 위험자산을 새 비율로 묶으면 분자와 분모가 서로 다르게 변하므로 샤프비율 자체가 달라질 수 있다. 높은 수익률 자산에 돈을 더 넣는 것과 샤프비율이 높은 조합을 찾는 것은 같은 문제로 단순화되지 않는다. 수익률이 음의 초과수익이거나 공매도 비중을 허용하면 기하적 해석에도 조건이 붙는다.
\par

\paragraph{샤프비율이 높으면 반드시 돈을 버나?}

샤프비율은 기대 초과수익을 표준편차로 나눈 \textbf{사전 비교 지표}다. 무위험수익률 3\%, 기대수익률 9\%, 표준편차 12\%면 샤프비율은 0.5지만, 실제 한 해 수익률은 -10\%가 될 수도 있다. 지표가 높은 전략도 손실 가능성은 있으며, 차입금리·수수료·실현된 수익률이 가정과 다르면 실제 샤프비율도 달라진다.
\par

\subsection{슬라이드 9: 두 위험자산 수익률의 결합}

자산 1에 $w$, 자산 2에 $1-w$를 투자하면 $R_p=wR_1+(1-w)R_2$다. 평균은 $\mu_p=w\mu_1+(1-w)\mu_2$이고, 평균에서의 편차를 제곱해 기대값을 취하면 아래 식을 얻는다.
\par

\[
\sigma_p^2=w^2\sigma_1^2+(1-w)^2\sigma_2^2+2w(1-w)\rho_{12}\sigma_1\sigma_2.
\]

마지막 항이 분산효과를 결정한다. 평균은 $w$의 일차함수이지만 분산은 이차함수다. $w$가 0과 1 사이여도 포트폴리오 표준편차는 두 표준편차의 단순 가중평균과 같지 않다.
\par

\paragraph{공분산 항을 현금흐름에서 다시 유도하기}

평균에서 뺀 두 자산의 수익률을 $X=R_1-\mu_1$, $Y=R_2-\mu_2$라고 두자. 포트폴리오 편차는 $wX+(1-w)Y$다. 이를 제곱하면 $w^2X^2+(1-w)^2Y^2+2w(1-w)XY$이고 기대값을 취하면 첫 두 항은 개별 분산, 마지막 항은 공분산이다. 상관계수 정의 $c=\rho\sigma_1\sigma_2$를 넣으면 슬라이드의 식이 된다.
\par

분산효과는 단순히 자산을 두 개 샀기 때문에 생기는 것이 아니다. 두 자산이 매번 같은 방향으로 완전히 움직이면 공통 위험이 남는다. \textbf{서로 다른 상태에서 지급되는 현금흐름}을 조합할 수 있을 때 같은 평균 부근에서 전체 변동을 줄일 여지가 생긴다.
\par

\paragraph{상관이 음수면 포트폴리오 손실이 사라지나?}

상관계수 -0.5는 평균적으로 서로 반대로 움직이는 경향을 뜻한다. 그러나 두 자산이 특정 날 함께 손실을 낼 가능성까지 없애지는 않는다. 분산식으로 계산한 위험은 \textbf{확률분포의 요약값}이지, 매 거래일의 손익을 고정하는 계약이 아니다. 환율 헤지나 선물 헤지도 기초자산과 헤지수단의 가격이 정확히 반대로 움직이지 않으면 잔여 위험이 남는다.
\par

\subsection{슬라이드 10: 숫자로 보는 분산 효과}

자산 1의 $(\mu_1,\sigma_1)=(.10,.15)$, 자산 2의 $(\mu_2,\sigma_2)=(.14,.20)$, 상관계수는 $.20$이라고 하자. 공분산은 $.20(.15)(.20)=.006$이다. 따라서 $\mu_p=.14-.04w$이고 분산은 다음처럼 정리된다.
\par

\[
\sigma_p^2=.0225w^2+.04(1-w)^2+.012w(1-w)=.0505w^2-.068w+.04.
\]

$w=.6$이면 기대수익 $.116$, 분산 $.01738$, 표준편차 약 $.1318$이다. 두 자산 중 변동성이 낮은 자산 1의 $.15$보다도 작다. 두 자산이 완전히 함께 움직이지 않기 때문에 가능한 결과다.
\par

\paragraph{완전 동조와 비교한 분산투자 이익}

비중 $w=.6$에서 두 자산이 완전히 함께 움직인다면 표준편차는 $.6(.15)+.4(.20)=.17$이고 분산은 $.0289$다. 실제 예시의 상관 $.20$에서는 분산이 $.01738$이므로, 완전 동조와 비교한 \textbf{분산투자 이익}은 $.0289-.01738=.01152$다.
\par

이 차이는 두 자산의 개별 분산 항에서 생기지 않는다. 공분산 항이 완전 동조일 때 $.0144$, 상관 $.20$일 때 $.00288$로 달라진 결과다. 평균은 두 경우 모두 11.6\%다. 분산효과를 비교할 때는 비중과 개별 평균·표준편차를 고정하고 상관만 바꾼다.
\par

\paragraph{음의 상관은 같은 비중의 위험을 얼마나 낮추나?}

두 자산의 표준편차가 각각 20\%이고 절반씩 보유한다고 하자. 상관이 $\rho=-.5$이면 분산은 $.25(.04)+.25(.04)+.5(-.02)=.01$, 표준편차는 10\%다. 같은 비중에서 $\rho=1$이면 분산은 $.04$, 표준편차는 20\%다. 두 자산의 개별 변동성이 같아도 \textbf{공분산 항} 하나 때문에 전체 위험은 크게 달라진다.
\par

\paragraph{분산 계산을 한 줄씩 전개하기}

$w=.6$이라면 첫 자산 자체 위험의 기여는 $.6^2(.15)^2=.00810$, 둘째 자산은 $.4^2(.20)^2=.00640$이다. 교차항은 $2(.6)(.4)(.2)(.15)(.20)=.00288$이다. 합은 $.01738$이고 제곱근은 약 $.13183$이다.
\par

상관이 1이었다면 교차항은 $.0144$가 되어 분산은 $.0289$, 표준편차는 $.17$이다. 상관을 .2로 낮추면 같은 투자비중과 개별 표준편차에서도 전체 위험이 17\%에서 13.18\%로 낮아진다. 이것이 분산투자 효과를 숫자로 분리해 보는 방법이다.
\par

\subsection{슬라이드 11: 분산 축에 그린 기회집합}

$w=1,.8,.6,.4,.2,0$으로 움직이면 기대수익은 $.10$에서 $.14$로 선형 증가하지만 분산은 먼저 감소했다가 증가한다. 분산이 $.0225$에서 $.01738$까지 줄어드는 동안 평균은 $.10$에서 $.116$으로 오르므로 자산 1 단독 보유보다 명백히 유리한 조합이 나타난다.
\par

분산을 가로축으로 놓으면 비중에 따른 점들이 최소분산점을 기준으로 접히는 곡선을 이룬다. \textbf{같은 분산에 두 비중}이 나올 때는 기대수익이 높은 쪽이 효율적이다. 두 해는 슬라이드 13의 이차방정식으로 계산한다.
\par

\paragraph{그림의 가로축을 숫자로 확인하기}

이 예에서 $w=1$이면 $(\mu_p,\sigma_p^2)=(.10,.0225)$다. $w=.8$이면 평균은 $.108$, 분산은 $.01792$이고, $w=.6$이면 평균 $.116$, 분산 $.01738$이다. 자산 2 비중을 처음 늘리는 구간에서는 평균이 올라가면서 분산은 내려간다. 이런 구간의 시작점은 평균-분산 관점에서 다른 조합에 지배된다.
\par

분산을 가로축에 썼다면 $.01738$은 \textbf{수익률 제곱 단위의 분산}이다. 해당 표준편차는 $\sqrt{.01738}\simeq13.18\%$다. 그림마다 가로축 표시가 $\sigma^2$인지 $\sigma$인지 확인하지 않으면 점의 위치와 위험 해석이 달라진다.
\par

\paragraph{같은 분산에 서로 다른 비중이 나올 수 있나?}

분산을 비중 $w$의 함수로 그리면 최소분산점 $w_G$를 중심으로 한 포물선이다. $w_G\simeq0.673267$인 예에서 $w=0.8$과 $w\simeq0.546535$는 최소점으로부터 거리가 같아 분산이 같다. 두 자산의 기대수익률이 다르면 두 조합의 기대수익률은 달라진다. 같은 위험의 두 점 중 평균이 낮은 쪽이 \textbf{지배}된다.
\par

\subsection{슬라이드 12: 표준편차 축에 그린 기회집합}

표준편차는 분산의 제곱근이므로 가로축을 $\sigma_p$로 바꾸면 기회집합의 모양이 바뀐다. 평균과 표준편차의 평면에서는 일반적으로 굽은 가지 두 개가 나타난다. 높은 기대수익을 갖는 위쪽 가지가 평균-분산 관점의 \textbf{효율적 프런티어}다.
\par

예시에서 $w=.6$의 점은 $(\sigma_p,\mu_p)\simeq(13.18\%,11.6\%)$다. 표준편차 축에서는 분산 축보다 가로 좌표가 커 보이지만 경제적 선택은 같다. 같은 위험에서는 위쪽 점을 선택한다.
\par

\paragraph{곡선의 두 가지가 생기는 까닭}

분산이 비중 $w$의 볼록한 이차식이면 최솟값 양쪽에 같은 높이의 분산을 갖는 두 비중이 있을 수 있다. 표준편차는 분산에 제곱근을 취한 것이므로 \textbf{같은 표준편차에도 두 비중}이 대응한다. 두 비중의 기대수익이 다르면 더 높은 평균을 주는 점만 효율적 가지에 남는다.
\par

예시에서 자산 2의 기대수익이 더 높으므로 최소분산점에서 자산 1 비중을 줄이는 방향이 위쪽 가지다. 효율적인 가지는 실제로 허용된 비중과 기대수익의 순서에 따라 정해진다. 공매도 금지처럼 허용 비중을 제한하면 곡선 중 일부만 실제 기회집합이 된다.
\par

\paragraph{그림의 곡선 전체를 실제로 살 수 있나?}

두 자산의 비중을 $w$와 $1-w$로 놓으면 수학적으로는 $w<0$이나 $w>1$도 그릴 수 있다. 공매도와 차입을 금지하면 가능한 비중은 $0\leq w\leq1$뿐이며, 곡선에서도 두 자산을 잇는 구간만 실제 투자기회가 된다. 상품별 최소 거래단위, 세금, 거래비용이 있으면 그 구간의 점도 정확히 구현하기 어렵다.
\par

\subsection{슬라이드 13: 가중치·평균·분산의 수학적 연결}

$c=\operatorname{Cov}(R_1,R_2)$, $A=\sigma_1^2+\sigma_2^2-2c=\operatorname{Var}(R_1-R_2)$라 두면 분산은 $Aw^2+2(c-\sigma_2^2)w+\sigma_2^2$다. $A>0$에서 완전제곱을 하면
\par

\[
\sigma_p^2=A(w-w_G)^2+\sigma_G^2,\qquad
w_G=\frac{\sigma_2^2-c}{A},\qquad
\sigma_G^2=\frac{\sigma_1^2\sigma_2^2-c^2}{A}.
\]

주어진 분산 $v$에 대한 비중은 $w=w_G\pm\sqrt{(v-\sigma_G^2)/A}$다. 따라서 실제 비중 해가 존재하려면 $v\ge\sigma_G^2$여야 하며, 최소점보다 위험이 크면 일반적으로 두 비중이 나온다. 이것은 분산식을 $w$의 이차방정식으로 풀 때 판별식이 0 이상이어야 한다는 조건과 같다.
\par

$\mu_1\ne\mu_2$이면 $\mu_p=\mu_2+(\mu_1-\mu_2)w$이고 $\mu_G=\mu_2+(\mu_1-\mu_2)w_G$다. 두 식에서 $w-w_G=(\mu_p-\mu_G)/(\mu_1-\mu_2)$를 앞의 분산식에 넣으면
\par

\[
\sigma_p^2=\sigma_G^2+\frac{A}{(\mu_1-\mu_2)^2}(\mu_p-\mu_G)^2.
\]

분산을 가로축에 놓으면 옆으로 열린 포물선이고, 표준편차 $\sigma_p$를 가로축에 놓으면 쌍곡선의 두 가지다. 비중이라는 매개변수를 없애 얻은 같은 기회집합이다. $A=0$이나 $\mu_1=\mu_2$이면 나눈 값이 0이 되므로 각각 따로 분석한다.
\par

\paragraph{공매도 제약과 퇴화하는 경우}

본문의 두 해 $w_G\pm\sqrt{(v-\sigma_G^2)/A}$는 공매도를 허용할 때의 수학적 해다. 공매도를 금지하면 $0\le w\le1$에 들어오는 해만 실제 투자할 수 있다. 예컨대 목표 분산을 매우 높게 잡으면 한 해가 음수, 다른 해가 1보다 커져 두 해 모두 허용되지 않을 수 있다. 그림에서는 원래 쌍곡선 중 두 단일 자산을 잇는 부분만 남는다.
\par

$A=0$이면 $R_1-R_2$가 일정해 비중을 바꿔도 분산이 변하지 않는다. $\mu_1=\mu_2$이면 비중을 바꿔도 평균이 변하지 않는다. 이 두 경우에는 일반적인 포물선·쌍곡선 그림 대신 실제로 움직이는 좌표를 먼저 확인해야 한다.
\par

\paragraph{두 자산의 상관이 1이면 최소분산 비중을 유일하게 정할 수 있나?}

두 자산 수익률이 항상 $R_1=R_2+k$라면 둘의 표준편차가 같고 상관은 1이다. 어느 비중으로 섞어도 공통 변동 부분이 그대로 남아 분산이 같다. 따라서 분산을 최소화하는 비중은 \textbf{유일하지 않다}. 이 경우 일반 유도식의 분모 $\sigma_1^2+\sigma_2^2-2\sigma_{12}$가 0이므로 식에 수치를 억지로 대입해서는 안 된다.
\par

\subsection{슬라이드 14: Global minimum variance (GMV) 포트폴리오 유도}

분산식의 도함수는 $d\sigma_p^2/dw=2Aw+2(c-\sigma_2^2)$다. 이를 0으로 놓으면 다음과 같다.
\par

\[
w_{\mathrm{GMV}}=\frac{\sigma_2^2-c}{\sigma_1^2+\sigma_2^2-2c},\qquad
\sigma_{\mathrm{GMV}}^2=\frac{\sigma_1^2\sigma_2^2-c^2}{\sigma_1^2+\sigma_2^2-2c}.
\]

$A>0$이면 이차식이 볼록하므로 이 정상점이 전역 최소다. 최소분산 비중은 **기대수익률 $\mu_1,\mu_2$에 의존하지 않는다.** 공매도가 금지되면 구한 비중을 $[0,1]$로 제한해야 하므로 내부해가 바깥에 있으면 경계의 단일 자산이 최소분산이 된다.
\par

\paragraph{공분산 추정이 비중을 바꾸는 방향}

개별 분산을 고정하고 공분산 $c$만 조금 바꾼다면 $w_G=(\sigma_2^2-c)/(\sigma_1^2+\sigma_2^2-2c)$의 민감도는
\par

\[
\frac{\partial w_G}{\partial c}
=\frac{\sigma_2^2-\sigma_1^2}
{(\sigma_1^2+\sigma_2^2-2c)^2}.
\]

둘째 자산의 분산이 더 크면 공분산 전망이 높아질수록 첫째 자산의 최소분산 비중도 높아진다. 두 자산이 비슷하게 움직인다는 정보가 강해질수록 상대적으로 덜 흔들리는 자산에 기댄다는 해석이다. $A$가 작으면 분모가 작아져 비중이 입력 오차에 민감해진다.
\par

\paragraph{최소분산 비중이 100\%를 넘는다면 계산이 틀렸나?}

$\sigma_1=10\%$, $\sigma_2=20\%$, $\rho=0.9$이면 공분산은 0.018이다. 첫 자산의 최소분산 비중은 $(0.04-0.018)/(0.01+0.04-0.036)\simeq1.57$이다. 첫 자산을 157\%, 둘째를 -57\% 보유하는 \textbf{공매도 허용} 해다. 공매도가 금지되면 이 해는 실행할 수 없고, 허용구간 $0\leq w\leq1$의 최소분산점은 경계 $w=1$이다.
\par

\paragraph{최소분산 해와 공매도 제약}

목적함수는 $Aw^2+2(c-\sigma_2^2)w+\sigma_2^2$다. $A>0$이면 이차계수는 양수라 극솟값이 하나다. $w^*=(\sigma_2^2-c)/A$가 $0$과 $1$ 사이면 두 자산을 양의 비중으로 보유한다. $w^*<0$이면 공매도 금지 조건에서는 $w=0$, $w^*>1$이면 $w=1$이 제약된 최소분산해다.
\par

특수하게 $A=\operatorname{Var}(R_1-R_2)=0$이면 두 자산의 수익률 차이가 확률적으로 일정하다. 그때 분산식의 이차항이 사라지므로 일반적인 '유일한 내부 최소값' 공식을 분모 0으로 계산해서는 안 된다.
\par

\subsection{슬라이드 15: 최소분산 포트폴리오 숫자 검산}

$\sigma_1=.15,\sigma_2=.20,c=.006$을 넣으면 $w_1^*=(.04-.006)/(.0225+.04-.012)=.673267$, $w_2^*=.326733$이다. 평균은 $.673267(.10)+.326733(.14)=.113069$이다. 분산은 $.0171089$, 표준편차는 약 $.13080$이다.
\par

최소분산점의 기대수익 11.31\%는 두 자산 평균 사이에 있다. 이 점에서 비중을 조금 바꾸면 분산은 $A(w-w_{\mathrm{GMV}})^2$만큼 증가한다. 최소분산점 부근의 곡선이 평평해 보이는 이유다.
\par

\paragraph{정답을 경계값과 비교하기}

예시의 최소분산 비중 $w_1\simeq.6733$은 0과 1 사이이므로 공매도 없이 구현 가능한 내부해다. 자산 1만 보유하면 표준편차 15\%, 자산 2만 보유하면 20\%인데 최소분산 조합은 약 13.08\%다. 계산한 비중을 원래 식에 다시 넣어 분산 $.0171089$를 얻고 두 경계의 분산 $.0225,.04$보다 작은지 확인하면 도출 과정의 부호 실수를 잡을 수 있다.
\par

최소분산점 근처에서는 비중을 조금 바꿔도 분산 변화가 이차적으로 작다. 예를 들어 $w=.6$과 최솟값의 차이는 $A(.6-.673267)^2\simeq.0002711$이며 실제로 $.01738-.0171089$와 같다. 이 안정성은 \textbf{분산값의 변화}에 관한 것이지 예상 수익이나 미래 최적 비중의 추정 오차가 작다는 뜻은 아니다.
\par

\paragraph{과거 상관이 조금만 달라져도 최소분산 비중이 바뀌나?}

$\sigma_1=15\%$, $\sigma_2=20\%$에서 상관이 0.2이면 첫 자산의 최소분산 비중은 약 0.673이다. 상관이 0.8이면 공분산은 0.024, 비중은 $(0.04-0.024)/(0.0225+0.04-0.048)\simeq1.10$으로 바뀐다. 과거 상관 추정치가 달라지면 최적 비중이 \textbf{공매도 여부까지 건드릴 만큼} 바뀔 수 있다.
\par

\subsection{슬라이드 16: 효율적인 가지의 방향}

이 예에서는 자산 2의 기대수익 $.14$가 자산 1의 $.10$보다 높다. 최소분산점 $w_1=.6733$에서 자산 1 비중을 줄이면 기대수익과 위험이 함께 올라가는 \textbf{효율적 가지}를 따라간다. 어느 점을 선택할지는 위험회피 정도에 달려 있다. 예컨대 $U=\mu_p-\tfrac a2\sigma_p^2$를 최대화하는 투자자는 $a$가 클수록 최소분산점에 가까운 점을 선호한다.
\par

\paragraph{지배를 두 좌표로 판정하기}

포트폴리오 A가 B보다 평균이 높고 표준편차가 낮거나 같다면 평균-분산 선호에서는 A가 B를 지배한다. 두 수치가 모두 높은 경우에는 투자자의 위험 선호가 있어야 선택할 수 있다. 최소분산점에서 자산 1 비중을 더 높이는 쪽은 이 예에서 평균이 떨어지고 위험이 커져 비효율적이다.
\par

예를 들어 $w_1=.9$의 평균은 $.104$다. 최소분산 비중 $.673267$을 중심으로 대칭인 $w_1\simeq.446534$는 같은 분산을 갖지만 평균은 약 $.12214$다. 같은 위험에 기대수익이 약 1.81\%포인트 높으므로 후자가 전자를 지배한다. 이 비교는 \textbf{기하적으로 같은 분산 높이의 두 해}를 읽은 결과다.
\par

\paragraph{효율적이지 않은 포트폴리오를 누군가 보유하면 비합리적인가?}

평균·분산 그림에서 지배된다는 평가는 두 투자자가 같은 기대수익률과 공분산을 믿고, 같은 거래 제약을 가지며, 평균과 분산만 중요하게 여긴다는 조건에 따른다. 한 투자자는 세금을 피하려고 오래 보유할 수 있고, 다른 투자자는 특정 시점의 원화 지급 의무를 맞추려 할 수 있다. 개인의 \textbf{현금흐름과 제약}이 다르면 단순한 평균·분산 순위만으로 행동을 판단할 수 없다.
\par

\subsection{슬라이드 17: 지배 포트폴리오를 직접 찾기}

자산 1 단독 보유의 분산은 $.15^2=.0225$다. 이 분산과 같은 $w$를 구하기 위해 $.0505w^2-.068w+.04=.0225$를 푼다. 한 해는 자산 1 자체인 $w=1$이고 다른 해는 $w\simeq.346535$다.
\par

두 번째 해의 기대수익은 $.346535(.10)+.653465(.14)\simeq.12614$로 자산 1의 10\%보다 높다. 즉 \textbf{같은 15\% 표준편차로 기대수익을 약 12.61\%까지 높일 수 있다.} 이것이 지배의 의미다. 제약이 바뀌어 두 번째 조합을 만들 수 없으면 실제 선택 가능성은 달라진다.
\par

\paragraph{꼭짓점 대칭으로 다른 비중 찾기}

분산을 $\sigma_p^2=A(w-w_G)^2+\sigma_G^2$로 쓰면 같은 위험을 갖는 두 비중은 $w_G$에서 같은 거리만큼 떨어져 있다. 한 해가 자산 1 단독 보유인 $w=1$이므로 다른 해는 $w'=2w_G-1$이다. $w_G\simeq.673267$을 대입하면 $w'\simeq.346534$가 된다. 근의 공식 없이도 짝이 되는 조합을 찾을 수 있다.
\par

이 두 조합은 분산이 같지만 기대수익은 각각 10\%와 약 12.61\%다. 어느 쪽이 위쪽 가지인지는 두 자산의 예상 평균이 정한다. \textbf{분산의 대칭성}은 상대 비중만 알려 주고 효율성 판단에는 평균 정보가 더 필요하다.
\par

\paragraph{분산이 같은 두 조합에서 지배 관계를 어떻게 찾나?}

최소분산점 양쪽에 있는 두 조합은 같은 표준편차를 가질 수 있다. 자산 1의 기대수익률이 8\%, 자산 2가 12\%라면 자산 1의 비중이 낮은 조합이 평균 수익률이 높다. 같은 위험에서 더 높은 평균을 주는 조합을 택하면, 낮은 평균의 점은 효율적 프런티어에서 제외된다. \textbf{위험만 최소화}하는 것과 주어진 위험에서 평균을 최대화하는 것은 서로 다른 문제다.
\par

\subsection{슬라이드 18: 상관계수와 프런티어의 이동}

$0<w<1$에서 분산식의 상관 관련 항은 $2w(1-w)\rho\sigma_1\sigma_2$다. 따라서 비중을 고정하고 $\rho$를 낮추면 분산이 감소한다. $\rho=1$이면 표준편차가 $w\sigma_1+(1-w)\sigma_2$가 되어 완전한 동조가 나타난다.
\par

$\rho=-1$이고 공매도 없이 $w=\sigma_2/(\sigma_1+\sigma_2)$를 선택하면 $w\sigma_1-(1-w)\sigma_2=0$이어서 수익률 변동을 완전히 상쇄한다. 실제 자산의 상관은 -1로 고정되지 않으며, 짧은 표본에서 추정한 낮은 상관이 위기에 유지된다고 가정해서는 안 된다. 공매도가 가능하면 비중 범위가 넓어져 그래프의 양 끝도 연장된다.
\par

\paragraph{상관의 변화가 분산에 미치는 크기}

비중과 개별 표준편차를 고정하면 $\partial\sigma_p^2/\partial\rho=2w(1-w)\sigma_1\sigma_2$다. $0<w<1$이면 이 값은 양수이므로 상관이 낮아질수록 분산도 낮아진다. 예시의 $w=.6$에서는 미분값이 $2(.6)(.4)(.15)(.20)=.0144$다. 상관이 $.20$에서 0으로 내려가면 분산은 $.0144\times.20=.00288$만큼 낮아져 $.01450$이 된다.
\par

$\rho=-1$이고 두 자산이 완전히 반대로 움직인다면 $w_1=.20/(.15+.20)\simeq.5714$에서 변동을 상쇄할 수 있다. 하지만 과거 추정 상관이 -1에 가까웠다는 사실만으로 미래 모든 상태에서 같은 상쇄가 일어난다는 보장은 없다. 실무에서는 특히 함께 크게 하락하는 상태를 따로 점검한다.
\par

\paragraph{과거에 분산투자가 잘됐는데 위기 때도 같을까?}

두 자산의 상관은 경제 상황에 따라 달라지는 추정치다. 평소에는 지역별 주식이 서로 다른 소식에 움직이다가 금융위기에는 자금 조달 충격으로 함께 하락할 수 있다. 따라서 과거 평균 상관으로 그린 프런티어 외에 \textbf{높은 상관과 큰 변동성}을 함께 넣은 스트레스 시나리오도 계산해야 한다.
\par

\paragraph{상관이 낮아도 분산투자 이익이 사라질 수 있나?}

평상시 자료로 추정한 상관이 낮더라도 위기 시 공통의 경기·유동성 충격 때문에 함께 하락할 수 있다. 또한 투자 비중을 지나치게 한 자산에 몰면 $2w(1-w)$가 작아져 상관의 완화 효과도 줄어든다. 상관은 투자 가능 자산의 가격 움직임을 요약할 뿐, 투자자의 노동소득이나 주택가격과의 공변동까지 자동으로 포함하지 않는다.
\par

따라서 개인의 전체 위험을 보려면 금융 포트폴리오 밖의 소득·부채·소비계획도 고려한다. 미시경제학적으로는 최종 소비의 나쁜 상태와 함께 손실이 나는 자산이 특히 부담스럽다.
\par

\subsection{슬라이드 19: 두 위험자산에 무위험자산 추가}

앞의 두 위험자산 중 비중 $w$로 만든 임의의 묶음을 $P$라 하고, 나머지 자금과 무위험자산 $r_f=.05$를 결합할 수 있다. 그러면 $P$마다 무위험점에서 시작해 $(\sigma_P,\mu_P)$를 지나는 직선이 생긴다.
\par

전체 재산 중 $y$를 $P$에 넣으면 $R_c=(1-y)r_f+yR_P$, $\mu_c=r_f+y(\mu_P-r_f)$, $\sigma_c=|y|\sigma_P$다. $0\le y\le1$이면 위험자산과 무위험자산을 함께 보유하고, $y>1$이면 무위험 금리로 빌려 $P$를 더 산다. $y\ge0$ 구간에서 $y$를 없애면 직선의 기울기는 $(\mu_P-r_f)/\sigma_P$다.
\par

같은 위험을 감수할 때 가장 높은 평균을 주는 직선은 이 기울기가 최대인 직선이다. 이 직선이 위험자산 효율적 프런티어에 접하는 점이 \textbf{접점 포트폴리오}다. 위험자산 간 비율 $w$와 전체 위험자산 규모 $y$는 서로 다른 선택이다.
\par

\paragraph{위험자산 내부 비중과 전체 투자 비중}

두 위험자산을 내부 비중 $w$로 묶은 포트폴리오를 $P$라 하고 전체 자산 중 $y$를 $P$에 둔다. 그러면 최종 수익률은 $R_c=(1-y)r_f+yR_P$다. $w$를 바꾸면 $P$의 평균과 표준편차, 따라서 무위험점에서 그은 직선의 기울기가 달라진다. $y$를 바꾸면 \textbf{선택한 한 직선 위에서} 투자자의 위험 크기만 조절한다.
\par

예를 들어 접점 위험자산 묶음이 정해진 뒤 $y=.5$라면 전체 자산의 절반은 그 묶음에, 절반은 무위험자산에 있다. 이때 접점 묶음 안의 자산 1 비중이 40\%라면 전체 재산에서 자산 1 비중은 $.5\times.4=.2$, 곧 20\%다. 이 곱셈을 모든 위험자산에 적용하면 계좌의 실제 보유량이 완성된다.
\par

\paragraph{접점 포트폴리오 비중과 전체 자산 비중은 같은가?}

접점 포트폴리오 내부에서 자산 1의 비중이 45.84\%, 자산 2가 54.16\%라고 하자. 투자금의 절반만 접점 포트폴리오에, 나머지 절반을 무위험자산에 두면 \textbf{전체 투자금 기준} 비중은 각각 22.92\%, 27.08\%, 50\%다. 내부 비중 $x_i$와 위험자산 전체 비중 $y$를 곱한 $yx_i$가 실제 계좌의 비중이다.
\par

\subsection{슬라이드 20: 접선이 효율적 프런티어를 바꾸는 이유}

위험자산만 있을 때의 위쪽 곡선에 무위험점 $(0,r_f)$에서 여러 직선을 그을 수 있다. 접점 $P$로 향하는 직선은 기울기 $(\mu_P-r_f)/\sigma_P$가 가장 커서 같은 위험에 가장 높은 평균을 준다. 투자자는 $P$의 \textbf{내부 비중}을 정한 뒤 무위험자산과 섞는 \textbf{전체 비중}을 선택한다.
\par

$P$보다 높은 위험의 직선 구간은 무위험 금리로 차입할 수 있을 때 실현된다. 차입금리가 높아지면 그 구간의 기울기와 선택 가능한 집합이 바뀐다.
\par

\paragraph{접선 이후에도 개인의 선택은 남는다}

같은 무위험 금리에서 투자 가능한 위험자산 묶음이 여러 개라면 샤프비율이 가장 큰 묶음의 직선이 다른 직선보다 같은 위험 수준에서 더 높은 평균을 제공한다. 이 기하적 결과가 \textbf{어떤 위험자산 묶음을 선택할지} 알려준다. 그러나 그 묶음에 전체 재산의 몇 퍼센트를 투자할지는 위험회피 정도가 결정한다.
\par

예를 들어 평균-분산 효용을 $U(y)=r_f+y(\mu_T-r_f)-(\gamma/2)y^2\sigma_T^2$로 근사하면 도함수는 $\mu_T-r_f-\gamma y\sigma_T^2$다. 내부 최적해는 $y^*=(\mu_T-r_f)/(\gamma\sigma_T^2)$이고 위험회피계수 $\gamma$가 클수록 위험자산 비중이 낮아진다. 차입 제한이 있으면 계산한 $y^*>1$을 그대로 실행할 수 없다.
\par

\paragraph{빌릴 때 금리가 더 높으면 접선도 하나인가?}

무위험자산에 맡길 때 7\%, 빌릴 때 10\%를 적용한다고 하자. 기대수익률 15\%, 표준편차 22\%의 위험포트폴리오를 섞으면 저축 구간 기울기는 $(15-7)/22\simeq0.364$지만 차입 구간은 $(15-10)/22\simeq0.227$이다. 차입을 시작하는 지점에서 선이 꺾인다. \textbf{같은 금리로 무제한 차입·대출}할 수 있다는 단순 모형의 직선 하나가 현실에서는 성립하지 않는다.
\par

\subsection{슬라이드 21: 접점 비중과 샤프비율의 유도}

평균 벡터, 공분산행렬, 무위험 수익률 5\%를 뺀 초과수익 벡터를 각각 다음처럼 둔다.
\par

\[
\boldsymbol{\mu}=(.10,.14)^\top,\qquad
\boldsymbol{e}=(.05,.09)^\top.
\]

\[
\boldsymbol{\Sigma}=
\begin{pmatrix}.0225&.006\\.006&.04\end{pmatrix}.
\]

위험자산 비중 벡터 $\boldsymbol{w}$의 샤프비율은 초과수익 $\boldsymbol{w}^\top\boldsymbol{e}$를 표준편차 $\sqrt{\boldsymbol{w}^\top\boldsymbol{\Sigma}\boldsymbol{w}}$로 나눈 값이다. 공분산으로 길이를 잰 코시-슈바르츠 부등식에서 이 비율은 $\boldsymbol{w}\propto\boldsymbol{\Sigma}^{-1}\boldsymbol{e}$일 때 최대다. 비중 합을 1로 정규화하면 다음과 같다.
\par

\[
\begin{aligned}
\boldsymbol{w}_T
&=\frac{\boldsymbol{\Sigma}^{-1}\boldsymbol{e}}
{\boldsymbol{1}^{\top}\boldsymbol{\Sigma}^{-1}\boldsymbol{e}}\\
&=(.45840,.54160)^\top.
\end{aligned}
\]

따라서 $\mu_T\simeq.121664$, $\sigma_T\simeq.139429$, $S_T=(.121664-.05)/.139429\simeq\mathbf{.51398}$이다. $.5416$은 \textbf{둘째 자산의 비중}에 가깝고 샤프비율이 아니다. 계산에는 수익률과 표준편차를 소수 단위로 일관되게 넣는다.
\par

\paragraph{접점 공식의 중간 계산까지 전개하기}

초과수익 벡터를 $\boldsymbol e=\boldsymbol\mu-r_f\boldsymbol1$로 두고, 일정한 위험 $\boldsymbol w^\top\Sigma\boldsymbol w=1$에서 $\boldsymbol w^\top\boldsymbol e$를 최대화하자. 라그랑주 일차조건은 $\boldsymbol e-\lambda\Sigma\boldsymbol w=0$이므로 $\boldsymbol w$는 $\Sigma^{-1}\boldsymbol e$에 비례한다. 비중의 합을 1로 정규화하면 본문의 접점 공식이 나온다. 공분산행렬이 가역이고 분모가 0이 아닌 기준 모형의 해다.
\par

숫자를 넣으면 공통 분모를 제외한 두 성분은 각각 $.04(.05)-.006(.09)=.001460$과 $.0225(.09)-.006(.05)=.001725$다. 합 $.003185$로 나누면 $(.45840,.54160)$이 된다. 그 비중으로 평균과 분산을 \textbf{원래의 두 자산 식에 다시 대입}해 $\mu_T\simeq.121664$, $\sigma_T\simeq.139429$를 확인한다. 둘째 자산 비중은 $.54160$, 조합의 샤프비율은 $.51398$이다.
\par

\paragraph{역행렬로 얻은 접점 비중이 왜 크게 튈 수 있나?}

두 자산이 거의 같은 수익률로 움직이면 공분산행렬은 특이행렬에 가까워진다. 역행렬 계산에서는 거의 0인 차이를 나누게 되어 작은 평균·상관 추정 오차가 큰 매수·공매도 비중으로 확대될 수 있다. 수식상 해가 존재해도 거래비용과 제약을 넣으면 실행이 어렵다. 결과가 과도한 레버리지를 요구한다면 \textbf{입력 추정치의 민감도}부터 확인한다.
\par

\paragraph{접점 공식의 두 자산형과 조건}

초과수익을 $e_1=\mu_1-r_f$, $e_2=\mu_2-r_f$, 공분산을 $c$라 하면 $\boldsymbol{\Sigma}^{-1}\boldsymbol e$는 공통 분모를 제외하고 $(\sigma_2^2e_1-ce_2,\;\sigma_1^2e_2-ce_1)$이다. 따라서 첫 자산의 접점 비중은 다음과 같다.
\par

\[
w_{1,T}=\frac{\sigma_2^2e_1-ce_2}{\sigma_2^2e_1+\sigma_1^2e_2-c(e_1+e_2)}.
\]

이 식은 공분산행렬이 가역이고 분모가 0이 아닌 경우의 무제약 해다. 공매도 금지, 서로 다른 차입·대출 금리, 기대 초과수익이 모두 음수인 경우에는 해의 의미나 효율적 집합이 달라진다. 실제 추정 평균을 조금만 바꾸어도 분모와 분자가 크게 달라질 수 있어 비중을 그대로 실무 규칙으로 사용하는 것은 위험하다.
\par

\subsection{슬라이드 22: 일반화와 실제 적용}

자산이 $n$개면 비중 벡터 $\boldsymbol{w}$, 평균 벡터 $\boldsymbol{\mu}$, 공분산행렬 $\boldsymbol{\Sigma}$로 $\mu_p=\boldsymbol{w}^\top\boldsymbol{\mu}$, $\sigma_p^2=\boldsymbol{w}^\top\boldsymbol{\Sigma}\boldsymbol{w}$라 쓴다. 두 자산에서 공분산 교차항을 전개한 식의 행렬 버전이다. 최소분산은 $\boldsymbol{w}^\top\boldsymbol{1}=1$이라는 제약 아래 이 이차식을 최소화한다.
\par

실제 포트폴리오에서는 예상 평균과 공분산의 추정 오차, 리밸런싱 비용, 차입금리, 공매도 제약, 집중도 제한을 추가한다. 특히 접점 비중은 기대수익의 작은 추정 변화에도 크게 흔들릴 수 있다. 따라서 이론적 해를 먼저 이해한 뒤 제약과 강건한 추정을 적용해야 한다.
\par

\paragraph{여러 자산으로 넓힐 때의 두 공식}

자산이 $n$개이고 공분산행렬 $\Sigma$가 가역이면, 비중 합 $\boldsymbol1^\top\boldsymbol w=1$ 아래 최소분산 문제의 라그랑주 일차조건은 $2\Sigma\boldsymbol w-\lambda\boldsymbol1=0$이다. 따라서 $\boldsymbol w$는 $\Sigma^{-1}\boldsymbol1$에 비례하고, 합을 1로 맞추면 $\boldsymbol w_{\mathrm{GMV}}=\Sigma^{-1}\boldsymbol1/(\boldsymbol1^\top\Sigma^{-1}\boldsymbol1)$이다. 두 자산 공식은 이 식을 성분별로 풀어 쓴 경우다.
\par

접점 포트폴리오에서는 $\boldsymbol1$ 대신 초과수익 벡터 $\boldsymbol\mu-r_f\boldsymbol1$이 들어간다. 이론적으로는 선명한 식이지만 역행렬은 입력값의 작은 오차를 큰 비중 변화로 증폭할 수 있다. 공매도 금지나 집중도 한도를 넣으면 단순 역행렬식만으로 해가 나오지 않으므로 제약을 포함한 최적화로 돌아가야 한다.
\par

\paragraph{세 자산으로 늘리면 분산식에 무엇이 추가되나?}

세 자산의 비중이 $w_1,w_2,w_3$이면 분산에는 개별 분산 항 $w_i^2\sigma_i^2$ 세 개와 서로 다른 자산쌍의 공분산 항 $2w_1w_2\sigma_{12}$, $2w_1w_3\sigma_{13}$, $2w_2w_3\sigma_{23}$가 들어간다. 두 자산 식에서 빠졌던 \textbf{셋째 자산과의 두 관계}가 새로 생긴다. 다음 장의 행렬식은 이 모든 항을 한 번에 적는다.
\par

\clearpage

\section{4부 · 여러 자산의 평균·분산 분석}

\subsection{슬라이드 1: 여러 위험자산으로 넓히는 평균·분산 분석}

두 위험자산에서는 첫 자산의 비중 $w$만 정하면 둘째 비중은 $1-w$였다. 자산이 $n$개면 비중은 $w_1,\ldots,w_n$이고 합계가 1이라는 조건 아래 $n-1$개를 독립적으로 고를 수 있다. 그림에 모든 조합을 직접 표시하기 어려워지므로 \textbf{벡터와 행렬}로 평균·분산을 한 번에 계산한다.
\par

기대수익은 비중의 가중평균이고, 위험에는 모든 자산쌍의 공분산이 들어간다. 벡터와 행렬로 이 두 계산을 묶으면 자산 수가 늘어나도 같은 원리로 최소분산 포트폴리오와 효율적 프런티어를 구할 수 있다.
\par

\paragraph{두 자산의 식에서 달라지는 것은 차원이다}

두 자산에서 $w_2=1-w_1$을 대입해 얻은 분산식은 $n$개 자산의 $\mathbf w^\top\Omega\mathbf w$에 $n=2$를 넣은 특수한 경우다. 계산의 목적도 같다. 주어진 평균에서 위험을 가장 적게 만드는 비중을 구한다. 다만 자산이 셋이면 공분산이 세 쌍, 넷이면 여섯 쌍이 되어 항을 하나씩 적는 방식이 금방 번거로워진다.
\par

행렬은 비중과 수익률을 질서 있게 묶고 모든 공분산 항을 포함하는 표기다. 공매도와 차입의 허용 여부는 최적화의 \textbf{제약 조건}에 담긴다.
\par

\paragraph{자산이 10개면 공분산을 몇 개나 알아야 하나?}

각 자산의 분산 10개와 서로 다른 두 자산쌍의 공분산 $10\times9/2=45$개, 총 55개의 분산·공분산 원소가 필요하다. 기대수익률 10개까지 더하면 65개 입력값이다. 행렬 $\Omega$는 10×10으로 100칸이지만 대칭이어서 중복 칸은 같은 값이다. 자산 수가 늘수록 분산투자의 기회와 \textbf{추정해야 할 입력값의 수}가 함께 늘어난다.
\par

\subsection{슬라이드 2: 평균·분산 선택의 질문을 먼저 고정하기}

평균·분산 분석은 가능한 모든 비중 가운데 \textbf{같은 기대수익률에서 분산이 가장 낮은 조합}을 찾는 문제다. 이렇게 얻은 점들의 집합이 최소분산 경계다. 그중 최소분산점보다 기대수익률이 높은 쪽을 효율적 프런티어라고 부른다. 위험에 비해 기대수익이 낮은 아래쪽 가지는 같은 위험에서 더 높은 평균을 주는 조합에 지배된다.
\par

자산의 과거 평균을 미래 기대수익률로, 과거 공분산을 미래 공분산으로 쓸 때는 추정 오차가 생긴다. 여기서는 먼저 입력값이 주어졌다고 가정해 \textbf{기준 해}를 유도하고, 마지막에 공매도·차입금리·추정 오차가 그 해를 어떻게 바꾸는지 구분한다. 수학적으로 가능한 비중과 실제 거래 가능한 비중도 같은 뜻이 아니다.
\par

\paragraph{최소분산 경계와 효율적 경계의 차이}

목표수익률이 8\%, 10\%, 12\%일 때 각각 가장 낮은 분산을 주는 비중을 구해 연결하면 최소분산 경계가 된다. 이 중 8\% 점과 12\% 점의 분산이 같다면 8\% 점은 지배된다. 따라서 경계 전체가 효율적이지 않다. GMVP보다 평균이 높은 가지가 평균·분산 선호에 맞는 효율적 부분이다.
\par

위험을 싫어하는 정도가 서로 다른 투자자는 효율적 가지의 서로 다른 점을 고를 수 있다. \textbf{경계 계산}은 가능한 최선의 조합을 추리는 작업이고, 그중 어느 점을 살지는 투자자의 선호와 지급 목표가 정한다.
\par

\paragraph{효율적 경계에서는 누가 어느 점을 고르나?}

두 효율적 포트폴리오 A와 B의 평균·표준편차가 각각 $(10\%,12\%)$, $(13\%,20\%)$라고 하자. 평균·분산 효용 $U=\mu-(a/2)\sigma^2$에서 $a=2$인 투자자는 $U_A=.0856$, $U_B=.09$이므로 B를 고른다. $a=5$이면 $U_A=.064$, $U_B=.03$이어서 A를 고른다. \textbf{효율적 경계}는 지배된 조합을 제거하고, 위험회피계수는 남은 점 중 개인의 선택을 정한다.
\par

\subsection{슬라이드 3: 네 단계로 읽는 계산 지도}

첫째, 비중 벡터 $\mathbf w$, 기대수익률 벡터 $\boldsymbol\mu$, 공분산행렬 $\Omega$를 정한다. 둘째, 무위험자산이 없을 때 목표 기대수익률 $m$을 고정하고 분산을 최소화한다. 셋째, $m$을 바꾸며 경계를 그린 뒤 그중 글로벌 최소분산점(GMVP)을 찾는다. 넷째, 무위험자산을 넣으면 가장 높은 기대 초과수익을 위험 한 단위당 주는 접점 포트폴리오를 구한다.
\par

세 최적화의 목적을 구분한다. 목표수익률별 해는 \textbf{평균을 고정한 상태에서} 분산을 최소화한다. GMVP는 그 경계에서 \textbf{분산 자체}가 가장 낮은 점이고, 접점 포트폴리오는 무위험수익률을 기준으로 \textbf{샤프비율}이 가장 높은 점이다.
\par

\paragraph{순서를 지키면 세 최적화가 섞이지 않는다}

목표 평균 $m$을 고정한 뒤 분산을 최소화하면 함수 $\sigma_p^2(m)$이 나온다. 다시 이 함수를 $m$에 대해 최소화하면 GMVP가 나온다. 그 뒤에 무위험수익률 $r_f$를 추가해 $[\mu_p(m)-r_f]/\sigma_p(m)$의 기울기가 최대인 점을 찾으면 접점이다. 같은 $\Omega$를 사용해도 목적이 달라 해가 다르다.
\par

실제 분석에서는 입력값을 추정한 시점과 투자기간을 통일해야 한다. 월별 $\boldsymbol\mu$에 연간 $r_f$를 빼거나 일간 $\Omega$를 그대로 쓰면 계산 결과의 단위가 맞지 않는다.
\par

\paragraph{GMVP와 접점 포트폴리오는 언제 같을 수 있나?}

GMVP는 $\Omega$와 총비중만으로 정하고, 접점은 초과수익률 $\boldsymbol\mu-r_f\mathbf1$까지 사용한다. 보통 다르다. 예컨대 모든 위험자산의 기대수익률이 같다면 비중을 바꿔 평균을 높일 수 없으므로 그 평균에서 분산이 가장 작은 GMVP가 샤프비율도 높다($\mu_i>r_f$일 때). 이는 특수한 경우다. 목표수익률과 샤프비율은 서로 다른 \textbf{최적화 기준}임을 먼저 기억한다.
\par

\subsection{슬라이드 4: 벡터·전치·내적을 수익률에 연결하기}

$n$개 자산의 기대수익률을 열벡터 $\boldsymbol\mu=(\mu_1,\ldots,\mu_n)^\top$, 비중을 $\mathbf w=(w_1,\ldots,w_n)^\top$로 쓴다. 전치 $\mathbf w^\top$는 $n\times1$ 열벡터를 $1\times n$ 행벡터로 바꾼다. 그러므로 행벡터와 열벡터의 곱은 하나의 숫자인 내적이 된다.
\par

\[
\mu_p=\mathbf w^\top\boldsymbol\mu=\sum_{i=1}^n w_i\mu_i,\qquad
\mathbf 1^\top\mathbf w=\sum_{i=1}^n w_i=1.
\]

예를 들어 $\mathbf w=(.5,.3,.2)^\top$, $\boldsymbol\mu=(.08,.12,.16)^\top$이면 $\mu_p=.5(.08)+.3(.12)+.2(.16)=.108$이다. $\mathbf w^\top\boldsymbol\mu$는 $(1\times n)(n\times1)$의 곱이어서 결과가 하나의 숫자다. 그 숫자를 전치해도 같으므로 $\boldsymbol\mu^\top\mathbf w=\mathbf w^\top\boldsymbol\mu$다. 비중에 대한 gradient는 $\nabla_{\mathbf w}(\mathbf w^\top\boldsymbol\mu)=\boldsymbol\mu$로, 각 비중을 조금 높였을 때 평균이 해당 자산의 $\mu_i$만큼 움직인다는 뜻이다. 전치 기호 $\top$는 미분 표시가 아니다.
\par

\paragraph{모든 자산의 예상 수익이 같이 바뀐다면}

모든 자산의 기대수익률 전망이 $d$만큼 함께 높아지면 새 평균 벡터는 $\boldsymbol\mu+d\mathbf1$이다. 포트폴리오 평균의 변화는 $\mathbf w^\top(\boldsymbol\mu+d\mathbf1)-\mathbf w^\top\boldsymbol\mu=d\mathbf w^\top\mathbf1=d$다. 비중 합이 1인 한 어떤 조합이든 평균이 정확히 $d$만큼 오른다. 공통 뉴스와 자산별 상대 전망의 변화를 분리할 때 유용하다.
\par

반대로 비중을 $\delta\mathbf w$만큼 바꾸면서 총비중을 유지하면 $\mathbf1^\top\delta\mathbf w=0$이다. 평균의 변화는 $\delta\mathbf w^\top\boldsymbol\mu$이고, 공통 상승분 $d\mathbf1$은 이 이동의 효과에서 상쇄된다. 어떤 자산을 더 사는지는 \textbf{자산 간 예상 수익의 차이}가 결정한다.
\par

\paragraph{비중을 옮길 때 평균은 얼마나 변하나?}

두 자산의 평균이 8\%, 12\%이고 처음 비중이 $(.5,.5)$이면 포트폴리오 평균은 10\%다. 첫 자산에서 둘째 자산으로 전체 재산의 1\%포인트를 옮기면 비중은 $(.49,.51)$, 평균은 $.49(.08)+.51(.12)=.1004$로 0.04\%포인트 오른다. 예산 제약을 지키는 작은 이동의 효과는 $\Delta\mu_p=\varepsilon(\mu_2-\mu_1)$이다. 각 비중에 대한 편미분 $\mu_i$를 \textbf{함께 적용한 결과}다.
\par

\subsection{슬라이드 5: 공분산행렬과 이차형식의 뜻}

공분산행렬 $\Omega$의 대각 원소는 각 자산의 분산 $\sigma_i^2$, 비대각 원소는 $\sigma_{ij}=\operatorname{Cov}(R_i,R_j)$이다. 공분산은 순서를 바꾸어도 같으므로 $\Omega$는 대칭행렬이다. 세 자산의 분산식은 다음과 같이 전개된다.
\par

\[
\sigma_p^2=\mathbf w^\top\Omega\mathbf w
=\sum_{i=1}^3w_i^2\sigma_i^2
+2w_1w_2\sigma_{12}+2w_1w_3\sigma_{13}+2w_2w_3\sigma_{23}.
\]

공분산 항의 계수 2는 $(w_1X_1+w_2X_2+w_3X_3)^2$를 전개할 때 $X_iX_j$와 $X_jX_i$가 함께 나오기 때문이다. $\Omega\mathbf w$의 $i$번째 성분은 $\sum_jw_j\operatorname{Cov}(R_i,R_j)=\operatorname{Cov}(R_i,R_p)$다. 따라서 분산의 gradient $2\Omega\mathbf w$는 각 자산의 비중을 조금 늘릴 때의 \textbf{한계위험}을 나타낸다. 라그랑지안에서 $\tfrac12\mathbf w^\top\Omega\mathbf w$를 쓰면 미분 결과가 $\Omega\mathbf w$로 간결해진다. 임의의 비중에 대해 분산이 음수가 될 수 없으므로 공분산행렬은 양의 준정부호여야 한다.
\par

\paragraph{분산의 미분을 성분별로 검산하기}

$\mathbf w^\top\Omega\mathbf w=\sum_i\sum_j w_i\sigma_{ij}w_j$를 $w_k$로 미분하면 $k$가 왼쪽에 있는 항에서 $\sum_j\sigma_{kj}w_j$, 오른쪽에 있는 항에서 $\sum_iw_i\sigma_{ik}$가 나온다. 대칭성 $\sigma_{ik}=\sigma_{ki}$ 때문에 두 합은 같아 $2(\Omega\mathbf w)_k$다. 그래서 벡터 전체의 gradient가 $2\Omega\mathbf w$다.
\par

실제로 총비중 1을 유지하려면 자산 $k$를 늘린 만큼 다른 자산 $j$를 줄여야 한다. $w_k$를 $\varepsilon$만큼 늘리고 $w_j$를 줄일 때 분산의 일차 변화는 $2\varepsilon[(\Omega\mathbf w)_k-(\Omega\mathbf w)_j]$다. 따라서 제약을 무시한 한계위험 $2(\Omega\mathbf w)_k$와 \textbf{예산 제약 아래 교체 거래의 한계위험}은 다르다. 최적화에서 제약식을 함께 미분하는 이유가 여기에 있다.
\par

\paragraph{공분산행렬의 대각선만 알면 포트폴리오 위험을 구할 수 있나?}

두 자산의 표준편차가 각각 10\%이고 절반씩 보유한다면 대각선 항만으로 계산한 분산은 $.25(.01)+.25(.01)=.005$다. 상관이 1이면 공분산 .01의 교차항 $.5(.01)=.005$가 더해져 전체 분산은 .01, 표준편차는 10\%다. 상관이 -1이면 교차항이 -.005여서 분산은 0이다. \textbf{비대각선 원소}를 빼면 같은 자산으로도 정반대 결론이 나온다.
\par

\subsection{슬라이드 6: 세 위험자산의 평균과 분산을 직접 계산하기}

세 자산을 비중 $(.5,.3,.2)$로 담고 기대수익률이 $(.08,.12,.16)$이면 기대수익률은 앞에서 계산한 10.8\%다. 표준편차가 각각 10\%, 15\%, 20\%이고 공분산이 $\sigma_{12}=.003$, $\sigma_{13}=.002$, $\sigma_{23}=.009$라면 분산은 다음과 같다.
\par

\[
\begin{aligned}
\sigma_p^2={}&(.5)^2(.10)^2+(.3)^2(.15)^2+(.2)^2(.20)^2\\
&+2(.5)(.3)(.003)+2(.5)(.2)(.002)+2(.3)(.2)(.009)\\
={}&.008505,\qquad \sigma_p\simeq9.22\%.
\end{aligned}
\]

세 자산 중 표준편차가 가장 낮은 자산도 10\%인데 조합의 표준편차는 그보다 낮다. 개별 자산의 평균과 분산만으로는 이 결과를 예측할 수 없고 \textbf{서로 함께 움직이는 정도}를 알아야 한다. 비중의 합이 1인지, 공분산행렬이 대칭인지, 분산이 음수가 아닌지도 계산마다 검산한다.
\par

\paragraph{셋째 자산이 낮은 위험의 이유가 될 수 있다}

세 자산의 분산식에는 $\sigma_{12},\sigma_{13},\sigma_{23}$가 모두 들어간다. 둘째 자산과 셋째 자산의 상관이 높으면 세 번째 자산을 추가해도 기대만큼 위험이 줄지 않을 수 있다. 반대로 개별 변동성이 큰 자산이어도 기존 두 자산과 약하게 움직이면 전체 분산을 낮추는 비중이 있을 수 있다.
\par

예시 비중 $(.5,.3,.2)$의 분산 $.008505$는 소수 단위다. 백분율의 제곱으로 적으면 $85.05$라는 숫자가 되며, 제곱근을 취한 표준편차는 9.22\%다. 분산과 표준편차의 단위를 섞지 않는 것도 중요한 검산이다.
\par

\paragraph{새 자산을 추가하면 최소위험이 반드시 낮아지나?}

기존 두 자산의 비중을 유지하고 새 자산 비중을 0으로 둘 수 있다면, 세 자산의 가능한 조합 안에 기존 조합이 그대로 포함된다. 따라서 수학적 \textbf{최소분산 값은 증가할 수 없다}. 다만 새 자산이 기존 자산과 완전히 중복되면 값이 그대로일 수 있다. 매수 의무, 최소 보유 비중, 거래비용을 새로 붙이면 가능한 조합이나 순수익이 달라져 이 결론을 그대로 적용할 수 없다.
\par

\subsection{슬라이드 7: 자산이 늘면 투자기회집합이 어떻게 달라지나}

두 위험자산의 모든 조합은 비중 하나로 움직이는 곡선이다. 셋 이상이면 비중을 둘 이상 독립적으로 바꿀 수 있어 평균·표준편차 평면에 더 많은 조합이 생긴다. 그러나 그중 효율적인 \textbf{바깥쪽 경계}만 투자자의 평균·분산 선택에 남는다. 그림에서 개별 자산 A·B·C보다 더 높은 평균과 더 낮은 위험을 함께 주는 혼합 비중이 있을 수 있다.
\par

여러 자산을 섞으면 개별 자산보다 유리한 조합이 생기는 경우가 많다. 개별 자산이 경계에 놓이는지는 기대수익률과 공분산 구조에 달려 있다. 공매도를 금지하면 허용 비중의 영역이 좁아져 개별 자산이 제한된 효율적 경계의 끝점이 될 수도 있다.
\par

\paragraph{세 자산이라는 숫자에 특별한 정리가 있는 것은 아니다}

공매도를 허용하면 어떤 자산의 비중도 음수 또는 100\% 초과가 될 수 있으므로 가능한 조합이 넓어진다. 그렇다고 개별 자산이 수학적으로 \textbf{절대} 효율적 경계에 있을 수 없는 것은 아니다. 예컨대 공분산행렬이 단위행렬이고 기대수익률이 $(0,1,0)$인 세 자산에서 평균 1을 달성하는 최소분산 포트폴리오는 둘째 자산 100\%다.
\par

이 사례는 개별 자산이 경계에 있을 수 있음을 보여준다. 그림의 메시지는 여러 자산을 섞어 더 나은 조합을 찾을 가능성이 커진다는 \textbf{일반적인 직관}으로 읽는 편이 정확하다.
\par

\paragraph{개별 자산이 효율적 경계에 있는지 어떻게 판정하나?}

자산 $i$의 평균 $\mu_i$를 목표로 한 최소분산 문제를 푼다. 그 최소분산 값이 자산 $i$ 단독 보유의 분산 $\sigma_i^2$와 같으면 단독 보유가 \textbf{최소분산 경계}에 있다. 이어 $\mu_i$가 GMVP 평균 이상인지 확인하면 위쪽 \textbf{효율적 가지}에 있는지도 판정할 수 있다. 이 두 조건은 그림에 표시된 점의 위치보다 정확한 기준이다.
\par

\subsection{슬라이드 8: 무위험자산이 없을 때 목표수익률별 최소분산 해}

목표 기대수익률을 $m$이라고 하자. 총비중 $\mathbf 1^\top\mathbf w=1$과 평균 $\boldsymbol\mu^\top\mathbf w=m$을 동시에 만족하면서 $\tfrac12\mathbf w^\top\Omega\mathbf w$를 최소화한다. 두 제약의 승수 $\lambda,\gamma$를 둔 라그랑지안은 다음과 같다.
\par

\[
L=\tfrac12\mathbf w^\top\Omega\mathbf w
+\gamma(m-\boldsymbol\mu^\top\mathbf w)
+\lambda(1-\mathbf 1^\top\mathbf w).
\]

비중에 대한 일계조건은 $\Omega\mathbf w-\gamma\boldsymbol\mu-\lambda\mathbf1=0$이다. $\Omega$가 역행렬을 가지면 $\mathbf w=\Omega^{-1}(\gamma\boldsymbol\mu+\lambda\mathbf1)$을 얻는다. 이 비중은 반드시 \textbf{두 제약}을 만족해야 하므로 다시 평균식과 합계식에 대입한다. $A=\mathbf1^\top\Omega^{-1}\mathbf1$, $B=\boldsymbol\mu^\top\Omega^{-1}\mathbf1$, $C=\boldsymbol\mu^\top\Omega^{-1}\boldsymbol\mu$라 두면
\par

\[
\begin{pmatrix}C&B\\B&A\end{pmatrix}
\begin{pmatrix}\gamma\\\lambda\end{pmatrix}
=\begin{pmatrix}m\\1\end{pmatrix}.
\]

첫째 식에 $A$, 둘째 식에 $B$를 곱해 빼면 $(AC-B^2)\gamma=Am-B$다. 둘째 식에 $C$, 첫째 식에 $B$를 곱해 빼면 $(AC-B^2)\lambda=C-Bm$이다. 따라서
\par

\[
\Delta=AC-B^2,\qquad
\gamma=\frac{Am-B}{\Delta},\qquad
\lambda=\frac{C-Bm}{\Delta}
\]

를 얻는다. $A,B,C$는 긴 행렬 곱을 세 숫자로 요약한 표기다. 공분산행렬이 양의 정부호이고 자산들의 기대수익률이 모두 같지 않으면 $\Delta>0$이다. 이 조건은 목표 평균과 비중 합계가 독립적인 제약임을 보장한다.
\par

\paragraph{$\Delta$가 왜 양수인지}

$\Omega$가 양의 정부호이면 $\langle\mathbf x,\mathbf y\rangle=\mathbf x^\top\Omega^{-1}\mathbf y$가 내적을 정의한다. 이 내적에서 $A$와 $C$는 각각 $\mathbf1$과 $\boldsymbol\mu$의 길이 제곱, $B$는 둘 사이의 내적이다. 코시-슈바르츠 부등식은 $B^2\le AC$를 준다. 기대수익률 벡터가 모든 성분에서 같지 않으면 두 벡터가 비례하지 않아 $\Delta=AC-B^2>0$이다.
\par

평균이 거의 같아 $\Delta$가 0에 가까우면 승수 $\gamma=(Am-B)/\Delta$가 크게 변할 수 있다. 매우 비슷한 예상 평균의 차이로 높은 목표수익을 만들려면 매수와 공매도 비중을 크게 상쇄해야 할 수 있다. 이때 \textbf{역행렬 해의 숫자 안정성}과 공매도 가능성을 함께 확인한다.
\par

\paragraph{목표수익률을 고정하면 어떤 비중을 비교하나?}

서로 독립인 세 자산의 기대수익률이 $(6\%,10\%,14\%)$, 분산이 $(.01,.04,.09)$라고 하자. 목표 평균 10\%와 비중 합 1을 함께 적용하면 $w_1=w_3=t$, $w_2=1-2t$다. 공매도를 금지하면 $0\le t\le.5$다. 이 범위의 모든 조합은 \textbf{같은 10\% 평균}을 주므로 이제 위험만 비교한다.
\par

\[
\sigma_p^2(t)=.01t^2+.04(1-2t)^2+.09t^2
=.26t^2-.16t+.04.
\]

미분한 $.52t-.16=0$에서 $t=4/13$을 얻는다. 최적 비중은 $(4/13,5/13,4/13)$, 분산은 $1/65\simeq.01538$, 표준편차는 약 12.4\%다. 자산 2만 보유한 $(0,1,0)$도 두 제약을 만족하지만 표준편차가 20\%다. 제약은 \textbf{비교할 후보}를 정하고 목적함수는 그중 한 조합을 고른다.
\par

\subsection{슬라이드 9: 평균·분산 경계의 방정식을 끝까지 유도하기}

앞에서 구한 $\mathbf w$에 $\Omega$를 곱하면 $\Omega\mathbf w=\gamma\boldsymbol\mu+\lambda\mathbf1$이다. 왼쪽에 $\mathbf w^\top$를 곱하면 분산이 $\sigma_p^2=\gamma\mathbf w^\top\boldsymbol\mu+\lambda\mathbf w^\top\mathbf1=\gamma m+\lambda$가 된다. 승수 값을 대입해 정리하면
\par

\[
\sigma_p^2(m)=\frac{Am^2-2Bm+C}{\Delta}
=\frac1A+\frac{A}{\Delta}\left(m-\frac BA\right)^2.
\]

두 번째 식은 완전제곱으로 얻는다. $\mu_G=B/A$, $\sigma_G^2=1/A$라 쓰면 분산을 가로축에 둔 그림에서는 $\sigma_p^2=\sigma_G^2+(A/\Delta)(m-\mu_G)^2$가 \textbf{옆으로 열린 포물선}이다. 표준편차를 가로축에 두면 같은 식을
\par

\[
\frac{\sigma_p^2}{\sigma_G^2}
-\frac{(m-\mu_G)^2}{\Delta/A^2}=1
\]

로 써서 쌍곡선의 두 가지로 읽는다. $m=\mu_G+d$와 $m=\mu_G-d$는 위험이 같고, 평균이 높은 쪽이 다른 쪽을 지배한다. 따라서 효율적 프런티어는 $m\ge\mu_G$인 \textbf{위쪽 가지}다. $\Delta=0$이면 서로 다른 목표 평균으로 경계를 그리는 이 방식이 성립하지 않는다.
\par

\paragraph{목표수익률의 한계위험}

경계의 분산을 목표 평균 $m$에 대해 미분하면 $d\sigma_p^2/dm=2(A/\Delta)(m-\mu_G)$, 다시 미분하면 $d^2\sigma_p^2/dm^2=2A/\Delta>0$이다. 평균 목표를 GMVP의 $\mu_G$에서 멀리 올릴수록 필요한 위험이 커진다. 처음에는 작은 평균 변화에 따른 분산 증가가 작지만, 더 높은 목표에서는 증가 속도도 커진다.
\par

계수 $A/\Delta$는 경계의 \textbf{휘어짐}을 정한다. 자산들의 기대수익과 공분산이 바뀌면 이 값도 바뀐다. 그래서 같은 1\%포인트의 목표수익 상승이라도 어떤 투자기회집합에서는 위험을 조금, 다른 집합에서는 크게 늘려야 한다.
\par

\paragraph{목표 평균을 더 높일수록 위험은 얼마나 늘어나나?}

설명을 위한 경계식이 $\sigma_p^2=.01+2(m-.08)^2$라고 하자. GMVP의 평균은 8\%, 표준편차는 10\%다. 목표 평균을 10\%로 높이면 분산 $.0108$, 표준편차 약 10.39\%이고, 12\%로 높이면 분산 $.0132$, 표준편차 약 11.49\%다. 평균을 각각 2\%포인트 올렸지만 추가 분산은 $.0008$ 다음 $.0024$로 커진다. \textbf{더 높은 목표수익의 한계위험}이 증가하는 이차식의 모양을 보여 준다.
\par

\subsection{슬라이드 10: 글로벌 최소분산 포트폴리오를 찾기}

앞의 완전제곱식은 $m=B/A$일 때 분산이 가장 작다고 곧바로 알려준다. 따라서 GMVP의 기대수익률과 분산은 $\mu_G=B/A$, $\sigma_G^2=1/A$다. 이 값을 승수식에 넣으면 $\gamma=0$, $\lambda=1/A$가 되어 비중도 얻는다.
\par

\[
\mathbf w_G=\frac{\Omega^{-1}\mathbf1}{\mathbf1^\top\Omega^{-1}\mathbf1},
\qquad \mathbf1^\top\mathbf w_G=1.
\]

GMVP 비중은 주어진 공분산행렬만으로 결정된다. 기대수익률 벡터 $\boldsymbol\mu$는 이 포트폴리오의 평균 $B/A$를 평가하고 효율적 가지의 다른 점을 선택할 때 필요하다. 공매도 금지 제약이 있다면 음의 비중이 나온 경우 제약을 포함해 다시 최적화한다.
\par

\paragraph{GMVP의 평균을 몰라도 비중은 구할 수 있다}

$\mathbf w_G$는 $\mathbf1^\top\mathbf w=1$만을 제약으로 분산을 최소화해도 직접 얻는다. $L=\tfrac12\mathbf w^\top\Omega\mathbf w+\lambda(1-\mathbf1^\top\mathbf w)$를 미분하면 $\mathbf w=\lambda\Omega^{-1}\mathbf1$이고, 합계 제약으로 $\lambda=1/A$다. 이를 평균식에 넣은 결과가 $\mu_G=B/A$다.
\par

즉 평균을 목표로 먼저 넣고 다시 최소화하는 길과 평균 제약 없이 바로 풀어 가는 길이 일치한다. GMVP 비중은 공분산 추정에 민감할 수 있으므로 입력 행렬을 바꾼 민감도 분석이 필요하다.
\par

\paragraph{기대수익률 예측이 달라지면 GMVP도 바뀌나?}

공분산행렬 $\Omega$가 그대로라면 $\mathbf w_G=\Omega^{-1}\mathbf1/A$에는 기대수익률 $\boldsymbol\mu$가 없으므로 비중은 바뀌지 않는다. 하지만 예측된 평균이 바뀌면 GMVP의 기대수익률 $\boldsymbol\mu^\top\mathbf w_G$와 그 점이 투자자 목표에 적합한지는 달라진다. 또한 실제로 평균 전망을 바꾸는 뉴스가 공분산 전망도 바꾼다면 GMVP 자체가 바뀔 수 있다.
\par

\subsection{슬라이드 11: 두 포트폴리오로 경계 전체를 만드는 이유}

목표수익률 $m$에 따른 해는 $\mathbf w(m)=\Omega^{-1}(\gamma\boldsymbol\mu+\lambda\mathbf1)$이다. 승수 $\gamma,\lambda$는 $m$의 일차함수이므로, 서로 다른 목표수익률 $m_a,m_b$의 경계 포트폴리오를 알면 다른 $m$의 경계 포트폴리오는 두 비중의 선형결합으로 표현된다.
\par

\[
\mathbf w(m)=t\mathbf w(m_a)+(1-t)\mathbf w(m_b),
\qquad t=\frac{m_b-m}{m_b-m_a}.
\]

예를 들어 $\mathbf w_a=.30\mathbf w_G+.70\mathbf w_D$이고 $\mathbf w_1=.20\mathbf w_G+.80\mathbf w_D$, $\mathbf w_2=.60\mathbf w_G+.40\mathbf w_D$라면 $\mathbf w_a=.75\mathbf w_1+.25\mathbf w_2$다. 서로 다른 두 \textbf{경계 포트폴리오}로 경계의 다른 점을 만들 수 있다는 성질이 \textbf{two-fund separation}이다. 목표 평균이 두 기준 평균의 바깥에 있으면 계수 하나가 음수가 되어 차입 또는 공매도가 필요할 수 있다.
\par

\paragraph{두 펀드 분리의 조건과 해석}

두 기준 경계 포트폴리오의 기대수익률 $m_a\ne m_b$여야 둘 사이에서 목표 $m$의 비중 $t=(m_b-m)/(m_b-m_a)$를 정할 수 있다. $m$이 그 사이에 있으면 $0\le t\le1$이고, 벗어나면 한 기준 포트폴리오를 공매도하는 형식이 된다. 분리정리는 비중의 \textbf{선형 구조}에 관한 말이며, 모든 투자자가 같은 비중을 골라야 한다는 말은 아니다.
\par

무위험자산을 허용하면 기준 펀드 하나를 무위험자산으로 선택할 수 있다. 그때 다른 하나가 접점 포트폴리오가 되어 투자자는 둘의 혼합 비율만 바꾼다. 두 상황을 같은 이름으로 부르더라도 구성 펀드는 다르다.
\par

\paragraph{two-fund separation은 자산 두 개만 사라는 말인가?}

여기서 fund는 여러 주식을 담은 \textbf{완성된 포트폴리오}일 수 있다. 예컨대 펀드 A가 자산 10개, 펀드 B가 같은 자산 10개를 다른 비중으로 담고 있어도, 두 펀드의 비중을 조절해 최소분산 경계 위 다른 점을 만들 수 있다. 목표 평균이 두 펀드 평균 사이에 있으면 둘을 양의 비중으로 섞지만, 바깥 목표에는 한 펀드의 음의 비중이 필요할 수 있다.
\par

\subsection{슬라이드 12: 무위험자산을 넣으면 생기는 접선}

기대수익률 $r_f$가 확정된 자산과 위험자산 포트폴리오 $P$를 섞자. 전체 자금 중 $y$를 $P$에 넣으면 $\mu_c=r_f+y(\mu_P-r_f)$이고, $y\ge0$에서 $\sigma_c=y\sigma_P$다. $y$를 제거하면 $\mu_c=r_f+[(\mu_P-r_f)/\sigma_P]\sigma_c$가 되어 무위험점에서 $P$를 지나는 직선이 나온다.
\par

어느 $P$를 고를지는 그 직선의 기울기, 곧 샤프비율 $(\mu_P-r_f)/\sigma_P$를 비교해 정한다. 기울기가 가장 큰 직선은 위험자산의 효율적 프런티어에 접한다. 이 접점 $P$가 정해지면 투자자는 위험자산 내부 비중을 그대로 유지하고 $y$만 바꾸어 자신의 위험 수준을 선택한다. 그림의 직선 연장은 **같은 $r_f$로 차입까지 가능**하다는 이상화된 가정에 의존한다.
\par

\paragraph{같은 위험에서 두 직선을 비교하기}

같은 전체 계좌 표준편차 10\%를 목표로 두자. 샤프비율 $.5$의 위험자산 조합을 쓰면 기대 초과수익은 5\%포인트, 샤프비율 $.7$의 조합을 쓰면 7\%포인트다. 무위험수익률이 같다면 후자의 기대수익이 2\%포인트 높다. 접선의 선택은 \textbf{같은 위험에서 기대수익을 비교하는 계산}이다.
\par

접점보다 낮은 위험은 무위험자산을 함께 보유해 만든다. 접점보다 높은 위험을 직선 위에서 만들려면 통상 차입이 필요하다. 차입금리와 대출금리가 다르면 실제 투자기회집합은 접점에서 꺾이거나 다른 접점을 가질 수 있다.
\par

\paragraph{무위험자산을 더하면 왜 샤프비율을 비교하나?}

무위험수익률이 3\%이고 위험포트폴리오 A가 평균 9\%, 표준편차 12\%라면 기울기는 $(9-3)/12=.5$다. B가 평균 10\%, 표준편차 10\%라면 기울기는 $.7$이다. 전체 위험을 6\%로 맞추면 A를 50\% 보유한 계좌의 기대수익은 6\%, B를 60\% 보유한 계좌는 7.2\%다. 같은 위험에서 B를 포함한 직선이 위에 있다는 뜻이다.
\par

\subsection{슬라이드 13: 접점을 찾는 첫 단계: 초과수익의 위험 최소화}

각 위험자산의 초과 기대수익률을 $\mathbf e=\boldsymbol\mu-r_f\mathbf1$이라 하자. 위험자산에 넣는 전체 비중 $\mathbf w_c$가 목표 초과수익률 $e_c=\mathbf e^\top\mathbf w_c$를 만들 때 분산이 최소인 비중을 찾는다. 이 단계에서는 \textbf{위험자산 비중의 합을 1로 고정하지 않는다}. 나머지 금액이 무위험자산에 들어가기 때문이다.
\par

\[
L=\tfrac12\mathbf w_c^\top\Omega\mathbf w_c
+\theta(e_c-\mathbf e^\top\mathbf w_c),
\quad \Omega\mathbf w_c=\theta\mathbf e,
\quad \theta=\frac{e_c}{\mathbf e^\top\Omega^{-1}\mathbf e}.
\]

따라서 $\mathbf w_c=[e_c/(\mathbf e^\top\Omega^{-1}\mathbf e)]\Omega^{-1}\mathbf e$다. 목표 초과수익 $e_c$는 전체 크기만 바꾸고 위험자산 \textbf{서로 간의 비율}은 $\Omega^{-1}\mathbf e$가 정한다. 이 해가 최소분산인 이유는 양의 정부호 $\Omega$ 아래 목적함수가 볼록하기 때문이다.
\par

\paragraph{비중 합 제약을 먼저 넣지 않는 이유}

$\mathbf w_c$는 전체 계좌에서 위험자산에 배분한 비중 벡터다. 예를 들어 위험자산에 계좌의 60\%만 넣으면 $\mathbf1^\top\mathbf w_c=.6$이고 무위험자산 비중은 .4다. 이 단계에서 합계를 1로 고정하면 무위험자산을 섞는 자유도를 없애 버린다.
\par

$\Omega^{-1}\mathbf e$는 공분산으로 조정된 초과수익의 방향이다. $e_c$를 두 배로 올리면 $\mathbf w_c$와 표준편차도 같은 비율로 커진다. 이 선형성이 자본배분선의 직선 모양으로 이어진다.
\par

\paragraph{위험자산 비중의 합이 1보다 작아도 되는 이유는?}

두 위험자산의 초과 평균이 $\mathbf e=(.04,.08)^\top$, 공분산행렬이 $\operatorname{diag}(.01,.04)$라고 하자. $\Omega^{-1}\mathbf e=(4,2)^\top$이고 $\mathbf e^\top\Omega^{-1}\mathbf e=.32$다. 계좌 목표 초과수익 $e_c=.032$를 넣으면 $\mathbf w_c=(.032/.32)(4,2)=(.4,.2)$다. 위험자산 합계는 .6이고 남은 .4는 무위험자산이다. 실제 초과수익도 $.4(.04)+.2(.08)=.032$다. 위험자산 \textbf{내부 비율}은 $(2/3,1/3)$이지만 계좌 비중은 $(.4,.2)$다.
\par

\subsection{슬라이드 14: 둘째 단계: 위험자산 내부 비중을 1로 정규화하기}

접점 포트폴리오 $P$는 위험자산만으로 이루어진 \textbf{내부 조합}이므로 $\mathbf1^\top\mathbf w_P=1$이어야 한다. 앞의 방향 벡터 $\Omega^{-1}\mathbf e$를 합계가 1이 되도록 나누면
\par

\[
\mathbf w_P=\frac{\Omega^{-1}(\boldsymbol\mu-r_f\mathbf1)}
{\mathbf1^\top\Omega^{-1}(\boldsymbol\mu-r_f\mathbf1)}
=\frac{\Omega^{-1}(\boldsymbol\mu-r_f\mathbf1)}{B-Ar_f}.
\]

$B-Ar_f>0$이면 이 정규화의 초과 기대수익이 양수이며 샤프비율은 $\sqrt{\mathbf e^\top\Omega^{-1}\mathbf e}$다. $B-Ar_f=0$이면 방향벡터의 비중 합이 0이어서 유한한 합계 1 포트폴리오로 정규화할 수 없다. 전체 계좌에서는 위험자산에 $y\mathbf w_P$, 무위험자산에 $1-y$를 배분한다.
\par

\paragraph{두 종류의 비중을 구분하기}

$\mathbf w_P$는 위험자산 \textbf{내부에서} 합계가 1인 비중이고, $\mathbf w_c=y\mathbf w_P$는 \textbf{전체 계좌에서} 위험자산에 들어간 비중이다. $y=.6$이면 무위험자산이 .4, $y=1.4$이면 무위험자산이 -.4다. 뒤의 경우는 무위험 금리로 .4를 빌려 접점 포트폴리오를 1.4만큼 보유한다는 뜻이다.
\par

분모 $B-Ar_f=\mathbf1^\top\Omega^{-1}\mathbf e$는 위험자산 방향벡터의 성분합이다. 0에 가까우면 정규화된 접점 비중이 크게 튈 수 있으므로 입력값과 공매도 가능성을 확인해야 한다.
\par

\paragraph{접점 비중의 분모 $B-Ar_f$가 0이면 무엇이 문제인가?}

$\Omega^{-1}(\boldsymbol\mu-r_f\mathbf1)$의 성분합이 0이라는 뜻이다. 방향벡터 자체는 0이 아닐 수 있지만, 합계를 1로 만들려고 0으로 나눌 수는 없다. 그런 방향은 위험자산끼리 매수·공매도를 상쇄한 \textbf{자기자본 0의 롱숏 포지션}이 될 수 있다. 표준 접점 공식이 유한한 위험자산 내부 포트폴리오를 주지 않으므로 별도 제약이나 기회집합 분석이 필요하다.
\par

\subsection{슬라이드 15: 두 위험자산 공식과 같은 답인지 검산하기}

두 자산이면 $\mathbf e=(e_1,e_2)^\top$이고 $\Omega=\begin{pmatrix}\sigma_1^2&\sigma_{12}\\\sigma_{12}&\sigma_2^2\end{pmatrix}$이다. 역행렬에 $\mathbf e$를 곱한 뒤 두 성분의 합으로 나누면 첫 자산의 접점 비중은
\par

\[
w_{P,1}=\frac{\sigma_2^2e_1-\sigma_{12}e_2}
{(\sigma_2^2-\sigma_{12})e_1+(\sigma_1^2-\sigma_{12})e_2},
\qquad w_{P,2}=1-w_{P,1}.
\]

예를 들어 $\mu=(.10,.14)$, $r_f=.05$, $\sigma_1=.15$, $\sigma_2=.20$, $\rho=.20$이면 $\sigma_{12}=.006$이고 $e_1=.05,e_2=.09$다. 분자는 $.04(.05)-.006(.09)=.00146$, 분모는 $.034(.05)+.0165(.09)=.003185$이므로 $w_{P,1}\simeq45.84\%$, $w_{P,2}\simeq54.16\%$다. 두 자산에서 직접 미분해 얻은 비중과 일치한다. 역행렬 식의 공통 분모 $\sigma_1^2\sigma_2^2-\sigma_{12}^2$는 정규화 과정에서 약분된다.
\par

\paragraph{2×2 역행렬에서 공통분모가 사라지는 과정}

$D=\sigma_1^2\sigma_2^2-\sigma_{12}^2$라 두면 $\Omega^{-1}\mathbf e=D^{-1}(\sigma_2^2e_1-\sigma_{12}e_2,\,\sigma_1^2e_2-\sigma_{12}e_1)^\top$이다. 성분의 합으로 나눌 때 $D^{-1}$이 분자와 분모에서 함께 사라진다. 따라서 최종 비중식에 행렬식 $D$가 보이지 않는 것은 항을 빠뜨린 것이 아니다.
\par

숫자 예시의 첫 비중은 $.00146/.003185\simeq.4584$다. 이 값이 앞 장의 직접 미분 해와 같아야 일반화한 행렬 공식이 일관적이다.
\par

\paragraph{상관이 0인 두 자산에 공식을 적용하면?}

무위험수익률이 3\%, 두 위험자산의 평균이 9\%, 13\%, 표준편차가 10\%, 20\%, 상관이 0이라고 하자. 초과수익은 $(.06,.10)$, 공분산행렬은 $\operatorname{diag}(.01,.04)$이다. 역행렬을 곱한 방향벡터는 $(6,2.5)$이므로 접점 비중은 $(6/8.5,2.5/8.5)=(12/17,5/17)$이다. 두 자산 공식에 직접 넣어도 첫 비중은 $.04(.06)/[.04(.06)+.01(.10)]=12/17$로 같다. \textbf{일반식과 특수식의 일치}를 숫자로 확인한다.
\par

\subsection{슬라이드 16: 두 프런티어를 한 그림에서 비교하기}

무위험자산이 없으면 최소분산 경계는 $\sigma_p^2(m)=(Am^2-2Bm+C)/\Delta$이며, GMVP는 $(\sigma_G,\mu_G)=(1/\sqrt A,B/A)$다. 무위험자산을 허용하고 접점 $P$가 존재하면 새 효율적 선택선은 $r_f$에서 시작해 $P$를 지나는 직선이다.
\par

\[
\mu_c=r_f+\frac{\mu_P-r_f}{\sigma_P}\,\sigma_c
\quad(y\ge0),\qquad
\frac{\mu_c-r_f}{\sigma_c}=\frac{\mu_P-r_f}{\sigma_P}.
\]

$\sigma_c$는 전체 계좌의 표준편차, $\sigma_P$는 접점 포트폴리오 자체의 표준편차다. 투자자는 위험자산 내부 비중 $\mathbf w_P$를 고정한 채 전체 위험자산 비중 $y$를 조절한다. $y>1$ 구간은 무위험 금리로 차입할 수 있을 때 직선으로 이어지며, 차입금리가 높으면 그 구간의 기울기가 낮아진다.
\par

\paragraph{그림의 기호를 계좌의 금액으로 번역하기}

접점 $P$의 내부 비중이 45.84\%와 54.16\%일 때 $y=.5$로 보유하면 전체 계좌에는 각각 22.92\%와 27.08\%, 무위험자산 50\%가 들어간다. 그림에서 $P$보다 왼쪽의 점은 이렇게 위험자산 전체 규모를 줄인 조합이다. 경계 위의 다른 위험자산 비중으로 옮긴 결과가 아니다.
\par

$\sigma_c=y\sigma_P$는 $y\ge0$일 때이고 일반적으로는 $|y|\sigma_P$다. 위험을 가로축으로 그린 선의 왼쪽 연장까지 해석할 때 절댓값 조건을 잊으면 표준편차가 음수가 되는 오류가 생긴다.
\par

\paragraph{접점보다 오른쪽의 직선을 실제로 만들 수 있나?}

접점의 위험이 12\%이고 계좌 목표 위험이 18\%라면 단순 모형에서 $y=18/12=1.5$다. 계좌 자기자본 100만원에 50만원을 빌려 접점 포트폴리오 150만원을 보유하는 셈이다. 대출금리보다 차입금리가 높고 담보를 추가로 요구할 수 있으므로 실제 기대수익은 무위험 금리 직선을 연장한 값보다 낮을 수 있다. \textbf{선 위의 점과 실행 가능한 거래}를 구분한다.
\par

\subsection{슬라이드 17: 평균·분산 선택에서 CAPM으로 넘어가기}

지금까지는 기대수익률 $\boldsymbol\mu$와 공분산 $\Omega$가 주어졌을 때 투자자 한 명의 최적 비중을 구했다. CAPM은 투자자들이 공통 기대를 가지고, 위험자산의 총수요가 실제 공급과 같다는 \textbf{시장청산} 조건을 더한다. 이때 접점 포트폴리오를 시장포트폴리오에 연결하면 개별 자산의 \textbf{시장과 함께 움직이는 위험}이 기대수익률에 어떻게 반영되는지 설명할 수 있다.
\par

\paragraph{시장청산에서 베타 관계까지}

공통 기대와 무위험 차입·대출을 가정하고 시장청산으로 접점 비중이 시장포트폴리오 $\mathbf w_M$와 같아졌다고 하자. 접점의 방향 조건 $\Omega\mathbf w_M=k(\boldsymbol\mu-r_f\mathbf1)$을 바꾸어 쓰면 $\boldsymbol\mu-r_f\mathbf1=\lambda\Omega\mathbf w_M$다. 양변에 $\mathbf w_M^\top$를 곱하면 $\mu_M-r_f=\lambda\sigma_M^2$이므로 $\lambda=(\mu_M-r_f)/\sigma_M^2$다.
\par

$i$번째 성분 $[\Omega\mathbf w_M]_i=\operatorname{Cov}(R_i,R_M)$을 쓰면 $\mu_i-r_f=\beta_i(\mu_M-r_f)$, $\beta_i=\operatorname{Cov}(R_i,R_M)/\sigma_M^2$가 나온다. 개별 자산의 총표준편차 대신 \textbf{시장과 함께 움직이는 위험}이 기대수익에 들어가는 이유가 이 최적화·시장청산 연결에 있다.
\par

\paragraph{모두 같은 위험자산 묶음을 고르면 시장에는?}

두 위험자산의 시가총액 비중이 $(.6,.4)$라고 하자. 투자자들이 같은 접점 묶음 $P$를 사되 위험자산에 넣는 \textbf{총액}만 다르게 정하면, 합친 위험자산 보유량의 내부 비중도 $P$다. 시장청산에서는 이 합계가 공급 비중 $(.6,.4)$와 같아야 하므로 $P=(.6,.4)$가 된다. 개인의 위험회피 정도에 따라 총 투자액은 달라도 시장의 위험자산 내부 비중은 같다.
\par

\end{document}
